% !TEX program = pdflatex
% ==============================================================================
% TỔNG HỢP TOÀN DIỆN KIẾN THỨC GIẢI TÍCH 1 (CALCULUS 1 FOR ENGINEERS)
% Mã học phần: UET.MAT1050 - Trường Đại học Công nghệ, ĐHQGHN (VNU-UET)
% Tài liệu tham chiếu: James Stewart, "Calculus: Early Transcendentals" (9th ed.)
% Biên soạn tự động chuẩn LaTeX + TikZ + tcolorbox
% ==============================================================================
\documentclass[11pt,a4paper,oneside]{book}

% --- Gói ngôn ngữ tiếng Việt và Font Toán học ---
\usepackage[utf8]{inputenc}
\usepackage[T5]{fontenc}
\usepackage[vietnamese]{babel}
\usepackage{amsmath,amssymb,amsfonts,amsthm,mathtools}

% --- Gói căn chỉnh trang và Bảng biểu ---
\usepackage[left=2.2cm,right=2.2cm,top=2.4cm,bottom=2.4cm]{geometry}
\usepackage{enumitem}
\usepackage{longtable,booktabs,array}
\usepackage[table,dvipsnames,svgnames]{xcolor}
\usepackage{fancyhdr}
\usepackage{hyperref}

% --- Gói vẽ hình TikZ và Đóng khung định lý tcolorbox ---
\usepackage{tikz}
\usetikzlibrary{arrows.meta,calc,patterns,angles,quotes,decorations.markings}
\usepackage[most]{tcolorbox}

% --- Định nghĩa màu sắc chủ đạo ---
\definecolor{uetblue}{RGB}{12, 74, 110}
\definecolor{uetaccent}{RGB}{2, 132, 199}
\definecolor{slate}{RGB}{51, 65, 85}

% --- Cấu hình liên kết PDF ---
\hypersetup{
    colorlinks=true,
    linkcolor=uetblue,
    urlcolor=uetaccent,
    pdftitle={Tong Hop Kien Thuc Giai Tich 1 - UET.MAT1050},
    pdfauthor={Tro Ly Hoc Tap Giai Tich 1 - VNU-UET}
}

% --- Các môi trường hộp kiến thức (tcolorbox) ---
\newtcolorbox{congthuc}[1]{
    enhanced, breakable,
    colback=blue!3!white,
    colframe=uetblue,
    fonttitle=\bfseries,
    title={#1}
}

\newtcolorbox{vidu}[1]{
    enhanced, breakable,
    colback=teal!4!white,
    colframe=teal!80!black,
    fonttitle=\bfseries,
    title={Ví dụ minh họa: #1}
}

\newtcolorbox{ungdung}[1]{
    enhanced, breakable,
    colback=orange!5!white,
    colframe=orange!80!black,
    fonttitle=\bfseries,
    title={Ứng dụng thực tế: #1}
}

% --- Header & Footer ---
\pagestyle{fancy}
\fancyhf{}
\fancyhead[L]{\small\textbf{UET.MAT1050 -- GIẢI TÍCH 1 (CALCULUS 1)}}
\fancyhead[R]{\small\leftmark}
\fancyfoot[C]{\thepage}

\begin{document}

% ==============================================================================
% TRANG BÌA TÀI LIỆU (TITLE PAGE)
% ==============================================================================
\begin{titlepage}
    \centering
    \vspace*{1cm}
    {\Large \textbf{ĐẠI HỌC QUỐC GIA HÀ NỘI}}\\[4pt]
    {\Large \textbf{TRƯỜNG ĐẠI HỌC CÔNG NGHỆ (VNU-UET)}}\\[1.5cm]
    \rule{\textwidth}{1.5pt}\\[0.5cm]
    {\Huge \bfseries TỔNG HỢP TOÀN DIỆN KIẾN THỨC\\[8pt] GIẢI TÍCH 1 (CALCULUS 1)}\\[0.4cm]
    {\large \textbf{Mã học phần: UET.MAT1050 (5 Tín chỉ -- 15 Buổi học)}}\\[0.4cm]
    \rule{\textwidth}{1.5pt}\\[1.2cm]
    
    \begin{tcolorbox}[colback=slate!4!white,colframe=uetblue,width=0.92\textwidth,title={\textbf{Cấu trúc Tài liệu Chuẩn Đề cương UET.MAT1050}}]
    \begin{itemize}[leftmargin=1.5em, itemsep=4pt]
        \item \textbf{Hệ thống 11 Chương (Buổi 1 -- Buổi 15):} Bám sát 100\% đề cương chi tiết Giải tích 1 và giáo trình \textit{James Stewart, Calculus: Early Transcendentals (9th Ed.)}.
        \item \textbf{Trọng tâm Đạo hàm \& Tích phân:} Đầy đủ công thức, chú giải điều kiện, ý nghĩa cơ học/hình học và ví dụ giải chi tiết từng bước.
        \item \textbf{Hình vẽ minh họa TikZ vector:} Mã nguồn TikZ trực quan cho hàm ngược, giới hạn $\varepsilon$-$\delta$, tiếp tuyến, định lý Lagrange, tổng Riemann, vật thể tròn xoay, trường hướng Logistic và chuỗi Taylor.
        \item \textbf{Ứng dụng thực tế \& Bài tập trắc nghiệm:} Tích hợp mô hình thực tiễn (Khoa học máy tính, Điện - Điện tử, Cơ học, Kinh tế, Xác suất) và bộ câu hỏi trắc nghiệm có lời giải sau mỗi chương.
    \end{itemize}
    \end{tcolorbox}
    
    \vfill
    {\large Hà Nội -- Tài liệu Chuẩn hóa Biên dịch bằng \LaTeX{} (pdflatex / Overleaf)}
\end{titlepage}

% ==============================================================================
% MỤC LỤC TỔNG QUÁT (TABLE OF CONTENTS)
% ==============================================================================
\tableofcontents
\newpage

% ======================================================================
% BẢNG TRA CỨU NHANH CÁC CÔNG THỨC GIẢI TÍCH 1 CƠ BẢN & TRỌNG TÂM
% ======================================================================
\chapter*{Bảng Tra Cứu Nhanh Công Thức Giải Tích 1}
\addcontentsline{toc}{chapter}{Bảng Tra Cứu Nhanh Công Thức Giải Tích 1}

\section*{1. Giới hạn cơ bản \& Vô cùng bé tương đương (khi x $\to$  0)}
\begin{longtable}{|p{0.25\textwidth}|p{0.48\textwidth}|p{0.21\textwidth}|}
\hline
\rowcolor{blue!10} \textbf{Tên công thức} & \textbf{Biểu thức Toán học} & \textbf{Chú giải / Điều kiện} \\ \hline
\endfirsthead
\hline
\rowcolor{blue!10} \textbf{Tên công thức} & \textbf{Biểu thức Toán học} & \textbf{Chú giải / Điều kiện} \\ \hline
\endhead
\textbf{Giới hạn lượng giác cơ bản} & $\displaystyle \lim_{x \to 0} \frac{\sin x}{x} = 1, \quad \lim_{x \to 0} \frac{1 - \cos x}{x^2} = \frac{1}{2}$ & x tính bằng radian, x $\to$ 0 \\ \hline
\textbf{Giới hạn số e (Chuẩn Euler)} & $\displaystyle \lim_{x \to 0} (1 + x)^{\frac{1}{x}} = e, \quad \lim_{n \to \infty} \left(1 + \frac{1}{n}\right)^n = e$ & Dạng vô định $1^\infty$ \\ \hline
\textbf{Giới hạn mũ và Logarithm} & $\displaystyle \lim_{x \to 0} \frac{e^x - 1}{x} = 1, \quad \lim_{x \to 0} \frac{\ln(1 + x)}{x} = 1$ & Áp dụng quy tắc vô cùng bé tương đương \\ \hline
\textbf{Vô cùng bé tương đương cấp cao} & $\displaystyle \tan x \sim x, \quad \arcsin x \sim x, \quad (1+x)^\alpha - 1 \sim \alpha x$ & Chỉ thay thế trong tích/thương khi x $\to$ 0 \\ \hline
\end{longtable}

\section*{2. Bảng Đạo hàm cơ bản \& Quy tắc đạo hàm (Trọng tâm)}
\begin{longtable}{|p{0.25\textwidth}|p{0.48\textwidth}|p{0.21\textwidth}|}
\hline
\rowcolor{blue!10} \textbf{Tên công thức} & \textbf{Biểu thức Toán học} & \textbf{Chú giải / Điều kiện} \\ \hline
\endfirsthead
\hline
\rowcolor{blue!10} \textbf{Tên công thức} & \textbf{Biểu thức Toán học} & \textbf{Chú giải / Điều kiện} \\ \hline
\endhead
\textbf{Quy tắc đạo hàm tích, thương, hàm hợp} & $\displaystyle (uv)' = u'v + uv', \quad \left(\frac{u}{v}\right)' = \frac{u'v - uv'}{v^2}, \quad [f(g(x))]' = f'(g(x))\cdot g'(x)$ & u, v khả vi; v(x) $\neq$ 0 \\ \hline
\textbf{Đạo hàm lũy thừa, mũ, logarithm} & $\displaystyle (x^\alpha)' = \alpha x^{\alpha-1}, \quad (a^x)' = a^x \ln a, \quad (\ln |x|)' = \frac{1}{x}$ & a > 0, a $\neq$ 1, x $\neq$ 0 \\ \hline
\textbf{Đạo hàm hàm lượng giác} & $\displaystyle (\sin x)' = \cos x, \quad (\cos x)' = -\sin x, \quad (\tan x)' = \sec^2 x = \frac{1}{\cos^2 x}$ & $\cos$ x $\neq$ 0 đối với $\tan$ x \\ \hline
\textbf{Đạo hàm hàm lượng giác ngược} & $\displaystyle (\arcsin x)' = \frac{1}{\sqrt{1-x^2}}, \quad (\arccos x)' = -\frac{1}{\sqrt{1-x^2}}, \quad (\arctan x)' = \frac{1}{1+x^2}$ & |x| < 1 với $\arcsin$, $\arccos$ \\ \hline
\textbf{Đạo hàm hàm Hyperbolic} & $\displaystyle (\sinh x)' = \cosh x, \quad (\cosh x)' = \sinh x, \quad (\tanh x)' = \text{sech}^2 x = \frac{1}{\cosh^2 x}$ & $\sin$h x = $\frac{e^x - e^{-x}}{2}$, $\cos$h x = $\frac{e^x + e^{-x}}{2}$ \\ \hline
\textbf{Công thức Leibniz (Đạo hàm cấp n của tích)} & $\displaystyle (u \cdot v)^{(n)} = \sum_{k=0}^{n} C_n^k u^{(n-k)} v^{(k)}$ & u, v khả vi đến cấp n \\ \hline
\end{longtable}

\section*{3. Bảng Nguyên hàm \& Tích phân thông dụng (Trọng tâm)}
\begin{longtable}{|p{0.25\textwidth}|p{0.48\textwidth}|p{0.21\textwidth}|}
\hline
\rowcolor{blue!10} \textbf{Tên công thức} & \textbf{Biểu thức Toán học} & \textbf{Chú giải / Điều kiện} \\ \hline
\endfirsthead
\hline
\rowcolor{blue!10} \textbf{Tên công thức} & \textbf{Biểu thức Toán học} & \textbf{Chú giải / Điều kiện} \\ \hline
\endhead
\textbf{Tích phân hàm lũy thừa và phân thức đơn giản} & $\displaystyle \int x^\alpha \,dx = \frac{x^{\alpha+1}}{\alpha+1} + C \;(\alpha \neq -1), \quad \int \frac{dx}{x} = \ln|x| + C$ & x $\neq$ 0 \\ \hline
\textbf{Tích phân hàm mũ và lượng giác} & $\displaystyle \int e^{kx}\,dx = \frac{1}{k}e^{kx} + C, \quad \int \sin x\,dx = -\cos x + C, \quad \int \cos x\,dx = \sin x + C$ & k $\neq$ 0 \\ \hline
\textbf{Tích phân đưa về arctan và ln} & $\displaystyle \int \frac{dx}{x^2 + a^2} = \frac{1}{a}\arctan\left(\frac{x}{a}\right) + C, \quad \int \frac{dx}{x^2 - a^2} = \frac{1}{2a}\ln\left|\frac{x-a}{x+a}\right| + C$ & a > 0 \\ \hline
\textbf{Tích phân chứa căn thức bậc hai} & $\displaystyle \int \frac{dx}{\sqrt{a^2 - x^2}} = \arcsin\left(\frac{x}{a}\right) + C, \quad \int \frac{dx}{\sqrt{x^2 \pm a^2}} = \ln\left|x + \sqrt{x^2 \pm a^2}\right| + C$ & a > 0, |x| < a cho arcsin \\ \hline
\textbf{Công thức Tích phân từng phần} & $\displaystyle \int u\,dv = u\cdot v - \int v\,du, \quad \int_a^b u\,dv = \Big[u\cdot v\Big]_a^b - \int_a^b v\,du$ & Thứ tự ưu tiên đặt u: Nhất Lô - Nhì Đa - Tam Lượng - Tứ Mũ \\ \hline
\end{longtable}

\section*{4. Ứng dụng Hình học \& Vật lý của Tích phân}
\begin{longtable}{|p{0.25\textwidth}|p{0.48\textwidth}|p{0.21\textwidth}|}
\hline
\rowcolor{blue!10} \textbf{Tên công thức} & \textbf{Biểu thức Toán học} & \textbf{Chú giải / Điều kiện} \\ \hline
\endfirsthead
\hline
\rowcolor{blue!10} \textbf{Tên công thức} & \textbf{Biểu thức Toán học} & \textbf{Chú giải / Điều kiện} \\ \hline
\endhead
\textbf{Diện tích miền phẳng \& Thể tích tròn xoay (Phương pháp Đĩa/Vòng đệm)} & $\displaystyle S = \int_a^b |f(x) - g(x)|\,dx, \quad V_{Ox} = \pi \int_a^b \left|[f(x)]^2 - [g(x)]^2\right| dx$ & Quay quanh trục Ox trên [a, b] \\ \hline
\textbf{Thể tích vỏ hình trụ (Cylindrical Shells)} & $\displaystyle V_{Oy} = 2\pi \int_a^b x \cdot |f(x)|\,dx$ & Miền giới hạn bởi y=f(x), y=0, x=a, x=b (0 $\le$ a < b) quay quanh Oy \\ \hline
\textbf{Độ dài cung \& Diện tích mặt tròn xoay} & $\displaystyle L = \int_a^b \sqrt{1 + [f'(x)]^2}\,dx, \quad S_{Ox} = 2\pi \int_a^b |f(x)|\sqrt{1 + [f'(x)]^2}\,dx$ & f'(x) liên tục trên [a, b] \\ \hline
\textbf{Độ dài cung \& Diện tích trong Tọa độ cực r = r($\theta$)} & $\displaystyle S = \frac{1}{2}\int_\alpha^\beta [r(\theta)]^2\,d\theta, \quad L = \int_\alpha^\beta \sqrt{[r(\theta)]^2 + [r'(\theta)]^2}\,d\theta$ & $\alpha \le \theta \le \beta$ \\ \hline
\end{longtable}

\section*{5. Phương trình Vi phân Cấp 1 \& Chuỗi Taylor / Maclaurin}
\begin{longtable}{|p{0.25\textwidth}|p{0.48\textwidth}|p{0.21\textwidth}|}
\hline
\rowcolor{blue!10} \textbf{Tên công thức} & \textbf{Biểu thức Toán học} & \textbf{Chú giải / Điều kiện} \\ \hline
\endfirsthead
\hline
\rowcolor{blue!10} \textbf{Tên công thức} & \textbf{Biểu thức Toán học} & \textbf{Chú giải / Điều kiện} \\ \hline
\endhead
\textbf{PTVP Tuyến tính cấp 1: y' + P(x)y = Q(x)} & $\displaystyle y(x) = e^{-\int P(x)dx} \left( \int Q(x) e^{\int P(x)dx}\,dx + C \right)$ & Thừa số tích phân I(x) = $e^{\int P(x)dx}$ \\ \hline
\textbf{Phương trình tăng trưởng Logistic} & $\displaystyle \frac{dP}{dt} = kP\left(1 - \frac{P}{M}\right) \implies P(t) = \frac{M}{1 + A e^{-kt}}, \quad A = \frac{M - P_0}{P_0}$ & M: sức chứa tối đa (carrying capacity) \\ \hline
\textbf{Khai triển Maclaurin hàm mũ \& lượng giác} & $\displaystyle e^x = \sum_{n=0}^\infty \frac{x^n}{n!}, \quad \sin x = \sum_{n=0}^\infty (-1)^n \frac{x^{2n+1}}{(2n+1)!}, \quad \cos x = \sum_{n=0}^\infty (-1)^n \frac{x^{2n}}{(2n)!}$ & Bán kính hội tụ R = $\infty$ ($\forall$ x $\in \mathbb{R}$) \\ \hline
\textbf{Khai triển Maclaurin phân thức \& Logarithm} & $\displaystyle \frac{1}{1-x} = \sum_{n=0}^\infty x^n, \quad \ln(1+x) = \sum_{n=1}^\infty (-1)^{n-1}\frac{x^n}{n}, \quad (1+x)^\alpha = 1 + \alpha x + \frac{\alpha(\alpha-1)}{2!}x^2 + \cdots$ & |x| < 1 (riêng $\ln$(1+x) hội tụ tại x = 1) \\ \hline
\end{longtable}

% ======================================================================
% CHƯƠNG 1: HÀM SỐ VÀ MÔ HÌNH (Buổi 1)
% ======================================================================
\chapter{Hàm số và Mô hình}
\label{chap:chuong-1}

\begin{tcolorbox}[colback=slate!5!white,colframe=slate!75!black,title={\textbf{Mục lục Chương 1 \& Chuẩn đầu ra (Buổi 1)}}]
\textbf{Tổng quan:} Hàm số là đối tượng trung tâm của giải tích một biến. Chương này hệ thống hóa 4 phương pháp biểu diễn hàm số, các phép biến đổi đồ thị, tính chất đơn ánh, hàm ngược, hàm mũ, hàm logarithm tự nhiên và các hàm lượng giác ngược dùng để xây dựng mô hình toán học thực tế.\\[6pt]
\textbf{Nội dung chính trong chương:}
\begin{itemize}[leftmargin=1.5em, itemsep=2pt]
  \item \textbf{Mục 1.1:} Bốn cách biểu diễn hàm số \& Tập xác định - Tập giá trị
  \item \textbf{Mục 1.2 - 1.3:} Các mô hình toán học cơ bản \& Phép biến đổi đồ thị hàm số
  \item \textbf{Mục 1.4 - 1.9:} Hàm mũ, Đơn ánh, Hàm ngược và Logarithm cơ số tự nhiên
  \item \textbf{Mục 1.10:} Các hàm lượng giác ngược (Arcsin, Arccos, Arctan)
\end{itemize}
\textbf{Chuẩn đầu ra (LLO):}
\begin{itemize}[leftmargin=1.5em, itemsep=2pt]
  \item LLO1.1: Nắm vững và biết vận dụng khái niệm hàm số, đồ thị hàm số, các cách biểu diễn hàm số.
  \item LLO1.2: Kiểm tra được tính đơn ánh, tìm được hàm ngược và khảo sát tốc độ tăng của hàm mũ - logarithm.
\end{itemize}
\end{tcolorbox}

\section{[Mục 1.1] Bốn cách biểu diễn hàm số \& Tập xác định - Tập giá trị}
Một hàm số f là một quy tắc gán mỗi phần tử x thuộc tập D (tập xác định - Domain) với duy nhất một phần tử f(x) thuộc tập E. Tập hợp tất cả các giá trị f(x) được gọi là tập giá trị (Range). Hàm số được biểu diễn qua 4 cách: bằng lời (mô tả), bằng bảng số liệu, bằng đồ thị, và bằng công thức đại số.

\begin{congthuc}{Các công thức trọng tâm (1.1)}
\noindent\textbf{Định nghĩa ánh xạ hàm số và Tập giá trị:}
\[ f: D \to E, \quad x \mapsto y = f(x), \quad \text{Range}(f) = \{f(x) \mid x \in D\} \]
\noindent\textit{Chú giải:} Tiêu chuẩn đường thẳng đứng (Vertical Line Test): Mọi đường thẳng x = a cắt đồ thị tối đa tại 1 điểm.\\[4pt]
\noindent\textbf{Hàm chẵn và Hàm lẻ (Tính đối xứng):}
\[ f(-x) = f(x) \;(\text{Hàm chẵn - đối xứng qua } Oy), \quad f(-x) = -f(x) \;(\text{Hàm lẻ - đối xứng qua } O) \]
\noindent\textit{Chú giải:} Tập xác định D phải là tập đối xứng (x $\in$ D $\implies$ -x $\in$ D).\\[4pt]
\end{congthuc}

\noindent\textbf{Chú giải \& Lưu ý kỹ thuật:}
\begin{itemize}[leftmargin=1.5em, itemsep=2pt]
  \item Khi tìm tập xác định của hàm biểu diễn bằng công thức: Mẫu thức phải khác 0, biểu thức dưới căn bậc chẵn phải không âm ($\ge$  0), và biểu thức trong logarithm phải dương (> 0).
  \item Hàm hợp (Composite Function): (f $\circ$  g)(x) = f(g(x)) xác định khi x thuộc tập xác định của g và g(x) thuộc tập xác định của f.
\end{itemize}

\begin{vidu}{Tìm tập xác định và xét tính chẵn lẻ của hàm hợp}
\textbf{Đề bài:} Cho hàm số f(x) = $\ln\left(\frac{1 + x}{1 - x}\right)$. Hãy tìm tập xác định D và xét tính chẵn lẻ của hàm số.
\[ f(x) = \ln\left(\frac{1 + x}{1 - x}\right) \]
\textbf{Lời giải chi tiết:}
\begin{enumerate}[leftmargin=1.5em, itemsep=2pt]
  \item Bước 1: Tìm điều kiện xác định của biểu thức trong logarithm và mẫu số.
  \item Bước 2: Kiểm tra tính đối xứng của tập xác định D = (-1, 1).
  \item Bước 3: Tính f(-x) và so sánh với f(x) bằng tính chất logarithm.
\end{enumerate}
\[ \text{Điều kiện: } \frac{1+x}{1-x} > 0 \iff (1+x)(1-x) > 0 \iff -1 < x < 1 \implies D = (-1, 1) \]
\[ \forall x \in (-1, 1) \implies -x \in (-1, 1). \text{ Ta có: } f(-x) = \ln\left(\frac{1-x}{1+x}\right) = \ln\left(\frac{1+x}{1-x}\right)^{-1} = -\ln\left(\frac{1+x}{1-x}\right) = -f(x) \]
\noindent\textbf{Kết luận:} Vậy hàm số đã cho là hàm lẻ trên khoảng xác định D = (-1, 1).
\end{vidu}

\section{[Mục 1.2 - 1.3] Các mô hình toán học cơ bản \& Phép biến đổi đồ thị hàm số}
Từ một đồ thị chuẩn y = f(x), ta có thể dựng nhanh đồ thị của các hàm liên quan thông qua các phép tịnh tiến (ngang/dọc), phép co giãn (tỉ lệ) và phép lấy đối xứng qua các trục tọa độ.

\begin{congthuc}{Các công thức trọng tâm (1.2 - 1.3)}
\noindent\textbf{Phép tịnh tiến đồ thị (c > 0):}
\[ y = f(x) + c \;(\uparrow c), \quad y = f(x) - c \;(\downarrow c), \quad y = f(x - c) \;(\rightarrow c), \quad y = f(x + c) \;(\leftarrow c) \]
\noindent\textit{Chú giải:} Tịnh tiến theo trục hoành ngược dấu với hằng số cộng trực tiếp vào biến x.\\[4pt]
\noindent\textbf{Phép co giãn \& Phép đối xứng (c > 1):}
\[ y = c f(x) \;(\text{giãn dọc } c \text{ lần}), \quad y = f(cx) \;(\text{nén ngang } c \text{ lần}), \quad y = -f(x) \;(\text{đối xứng } Ox) \]
\noindent\textit{Chú giải:} Đồ thị y = f(|x|) giữ nguyên phần x $\ge$  0 rồi lấy đối xứng phần đó qua Oy.\\[4pt]
\end{congthuc}

\noindent\textbf{Chú giải \& Lưu ý kỹ thuật:}
\begin{itemize}[leftmargin=1.5em, itemsep=2pt]
  \item Mô hình tuyến tính y = mx + b có tốc độ thay đổi (hệ số góc m) không đổi.
  \item Mô hình đa thức P(x) = $a_n x^n$ + ... + $a_0$, hàm phân thức hữu tỉ R(x) = P(x)/Q(x), và hàm lũy thừa y = $x^a$ là nền tảng của xấp xỉ kỹ thuật.
\end{itemize}

\begin{vidu}{Biến đổi đồ thị từ hàm cơ sở}
\textbf{Đề bài:} Từ đồ thị parabol cơ bản y = $x^2$, hãy nêu quy trình dựng đồ thị hàm số y = -2(x - $3)^2$ + 5.
\[ y = x^2 \longrightarrow y = -2(x - 3)^2 + 5 \]
\textbf{Lời giải chi tiết:}
\begin{enumerate}[leftmargin=1.5em, itemsep=2pt]
  \item Bước 1: Tịnh tiến đồ thị y = $x^2$ sang phải 3 đơn vị để được y = (x - $3)^2$.
  \item Bước 2: Giãn dọc đồ thị theo hệ số 2 để được y = 2(x - $3)^2$.
  \item Bước 3: Lấy đối xứng qua trục hoành Ox để được y = -2(x - $3)^2$, sau đó tịnh tiến lên trên 5 đơn vị.
\end{enumerate}
\[ x^2 \xrightarrow{\text{phải } 3} (x-3)^2 \xrightarrow{\text{giãn dọc } 2} 2(x-3)^2 \xrightarrow{\text{đối xứng } Ox} -2(x-3)^2 \xrightarrow{\text{lên } 5} -2(x-3)^2 + 5 \]
\noindent\textbf{Kết luận:} Đỉnh của parabol dịch chuyển từ O(0, 0) tới điểm I(3, 5), bề lõm quay xuống dưới.
\end{vidu}

\section{[Mục 1.4 - 1.9] Hàm mũ, Đơn ánh, Hàm ngược và Logarithm cơ số tự nhiên}
Hàm số f được gọi là đơn ánh (one-to-one) nếu hai giá trị x khác nhau luôn cho hai giá trị f(x) khác nhau. Mọi hàm đơn ánh từ D lên R đều có hàm ngược $f^{-1}$ từ R lên D, và đồ thị của $f^{-1}$ đối xứng với đồ thị của f qua đường phân giác y = x.

\begin{congthuc}{Các công thức trọng tâm (1.4 - 1.9)}
\noindent\textbf{Định nghĩa đơn ánh \& Quan hệ hàm ngược:}
\[ x_1 \neq x_2 \implies f(x_1) \neq f(x_2), \qquad y = f(x) \iff x = f^{-1}(y) \]
\noindent\textit{Chú giải:} Tiêu chuẩn đường thẳng ngang (Horizontal Line Test): Mọi đường y = c cắt đồ thị tối đa 1 điểm.\\[4pt]
\noindent\textbf{Hàm Logarithm tự nhiên \& Công thức đổi cơ số:}
\[ y = \ln x \iff e^y = x, \qquad \log_a x = \frac{\ln x}{\ln a} \quad (a > 0, a \neq 1, x > 0) \]
\noindent\textit{Chú giải:} Tốc độ tăng khi x -> +$\infty$: ln x << $x^\alpha$  << $a^x$ (với $\alpha$  > 0, a > 1).\\[4pt]
\end{congthuc}

\noindent\textbf{Chú giải \& Lưu ý kỹ thuật:}
\begin{itemize}[leftmargin=1.5em, itemsep=2pt]
  \item Quy trình 3 bước tìm hàm ngược y = $f^{-1}$(x): (1) Viết y = f(x); (2) Giải phương trình tìm x theo y; (3) Đổi chỗ ký hiệu x và y để thu được y = $f^{-1}$(x).
  \item Cẩn thận tránh nhầm lẫn ký hiệu: $f^{-1}$(x) là hàm ngược chứ KHÔNG phải nghịch đảo $[f(x)]^{-1}$ = 1/f(x).
\end{itemize}

\begin{vidu}{Tìm hàm ngược của hàm phân thức mũ}
\textbf{Đề bài:} Chứng minh hàm số f(x) = $\frac{e^x}{1 + 2e^x}$ là đơn ánh trên $\mathbb{R}$ và tìm công thức hàm ngược $f^{-1}$(x) cùng tập xác định của nó.
\[ y = f(x) = \frac{e^x}{1 + 2e^x} \]
\textbf{Lời giải chi tiết:}
\begin{enumerate}[leftmargin=1.5em, itemsep=2pt]
  \item Bước 1: Biến đổi đại số biểu thức y = $e^x$ / (1 + $2e^x$) để rút $e^x$ theo y.
  \item Bước 2: Xác định miền giá trị của y (vì $e^x$ > 0 với mọi x $\in$  R).
  \item Bước 3: Lấy logarithm tự nhiên hai vế và hoán đổi biến x, y.
\end{enumerate}
\[ y(1 + 2e^x) = e^x \iff y = e^x(1 - 2y) \iff e^x = \frac{y}{1 - 2y} \]
\[ \text{Do } e^x > 0 \iff \frac{y}{1 - 2y} > 0 \iff 0 < y < \frac{1}{2} \implies \text{Range}(f) = \left(0, \frac{1}{2}\right) \]
\[ x = \ln\left(\frac{y}{1 - 2y}\right) \implies f^{-1}(x) = \ln\left(\frac{x}{1 - 2x}\right), \quad D_{f^{-1}} = \left(0, \frac{1}{2}\right) \]
\noindent\textbf{Kết luận:} Hàm ngược là $f^{-1}$(x) = $\ln\left(\frac{x}{1 - 2x}\right)$ với tập xác định (0, 1/2).
\end{vidu}

\section{[Mục 1.10] Các hàm lượng giác ngược (Arcsin, Arccos, Arctan)}
Vì các hàm lượng giác tuần hoàn nên không phải là đơn ánh trên toàn trục số. Để định nghĩa hàm ngược, ta thu hẹp tập xác định của chúng về một đoạn/khoảng chuẩn sao cho hàm đơn ánh và đạt toàn bộ tập giá trị.

\begin{congthuc}{Các công thức trọng tâm (1.10)}
\noindent\textbf{Miền xác định và Miền giá trị chuẩn của hàm lượng giác ngược:}
\[ \begin{aligned} y &= \arcsin x \iff \sin y = x, && x \in [-1, 1], \; y \in \left[-\frac{\pi}{2}, \frac{\pi}{2}\right] \\ y &= \arccos x \iff \cos y = x, && x \in [-1, 1], \; y \in [0, \pi] \\ y &= \arctan x \iff \tan y = x, && x \in \mathbb{R}, \; y \in \left(-\frac{\pi}{2}, \frac{\pi}{2}\right) \end{aligned} \]
\noindent\textit{Chú giải:} Hằng đẳng thức quan trọng: arcsin(x) + arccos(x) = $\pi$ /2 với mọi x $\in  [-1, 1]$.\\[4pt]
\end{congthuc}

\noindent\textbf{Chú giải \& Lưu ý kỹ thuật:}
\begin{itemize}[leftmargin=1.5em, itemsep=2pt]
  \item Khi rút gọn biểu thức dạng cos(arcsin x) hoặc tan(arccos x), hãy đặt $\theta$  = arcsin x và vẽ một tam giác vuông có cạnh đối là x, cạnh huyền là 1, từ đó cạnh kề là $\sqrt{1 - x^2}$.
\end{itemize}

\begin{vidu}{Rút gọn biểu thức chứa hàm lượng giác ngược}
\textbf{Đề bài:} Rút gọn biểu thức đại số A = $\cos$($\arctan$ x) và B = $\sin$(2$\arcsin$ x) với x $\in$ (-1, 1).
\[ A = \cos(\arctan x), \quad B = \sin(2\arcsin x) \]
\textbf{Lời giải chi tiết:}
\begin{enumerate}[leftmargin=1.5em, itemsep=2pt]
  \item Bước 1: Đặt $\theta$  = arctan x $\in$  (-$\pi$ /2, $\pi$ /2), suy ra tan $\theta$  = x và cos $\theta$  > 0.
  \item Bước 2: Sử dụng hệ thức 1 + $tan^2 \theta$  = 1 / $cos^2 \theta$  để tính cos $\theta$ .
  \item Bước 3: Đặt $\alpha$  = arcsin x $\in$  (-$\pi$ /2, $\pi$ /2), áp dụng công thức nhân đôi sin(2$\alpha$ ) = 2 sin $\alpha$  cos $\alpha$ .
\end{enumerate}
\[ \cos\theta = \frac{1}{\sqrt{1 + \tan^2\theta}} = \frac{1}{\sqrt{1 + x^2}} \implies A = \frac{1}{\sqrt{1 + x^2}} \]
\[ B = \sin(2\alpha) = 2\sin\alpha\cos\alpha = 2x\sqrt{1 - \sin^2\alpha} = 2x\sqrt{1 - x^2} \]
\noindent\textbf{Kết luận:} Ta thu được công thức đại số tường minh A = (1 + $x^2 )^{-1/2}$ và B = 2x$\sqrt{1 - x^2}$.
\end{vidu}

\section{Minh họa Hình học bằng TikZ}
\begin{figure}[htbp]
\centering
\begin{tikzpicture}[scale=1.05, >=stealth]
  % Lưới và trục tọa độ
  \draw[very thin, color=gray!25] (-2.5,-2.5) grid (3.5,3.5);
  \draw[->, thick] (-2.6,0) -- (3.7,0) node[right] {$x$};
  \draw[->, thick] (0,-2.6) -- (0,3.7) node[above] {$y$};
  
  % Đường phân giác y = x
  \draw[dashed, color=gray!80, thick] (-2.3,-2.3) -- (3.3,3.3) node[above right] {$y = x$};
  
  % Đồ thị hàm mũ y = e^x
  \draw[domain=-2.3:1.25, smooth, variable=\x, blue!80!black, very thick] 
    plot ({\x}, {exp(\x)}) node[above left] {$y = e^x$};
    
  % Đồ thị hàm logarithm y = ln(x)
  \draw[domain=0.1:3.4, smooth, variable=\x, red!80!black, very thick] 
    plot ({\x}, {ln(\x)}) node[below right] {$y = \ln x$};
    
  % Các điểm đối xứng đặc biệt
  \filldraw[blue!80!black] (0,1) circle (2pt) node[above left] {$(0,1)$};
  \filldraw[red!80!black] (1,0) circle (2pt) node[below right] {$(1,0)$};
  \draw[dotted, thick, teal] (0,1) -- (1,0);
\end{tikzpicture}
\caption{Hình 1.1: Tính đối xứng giữa đồ thị hàm mũ y = $e^x$ và hàm ngược y = ln x qua đường phân giác y = x.}
\label{fig:tikz-ch1-inverse}
\end{figure}

\section{Ứng dụng Thực tế \& Tối ưu hóa (Khoa học Máy tính \& Xử lý Tín hiệu)}
\begin{ungdung}{Thang đo Decibel và Độ phức tạp Thuật toán Logarithm}
Trong kỹ thuật âm thanh và viễn thông, cường độ tín hiệu thay đổi qua nhiều bậc độ lớn nên được nén bằng hàm logarithm (Decibel). Trong Khoa học Máy tính, thuật toán Tìm kiếm Nhị phân (Binary Search) là hàm ngược của quá trình nhân đôi kích thước mảng, giúp giảm số phép so sánh từ N xuống $\log_2$(N).
\[ L_{\text{dB}} = 10 \log_{10}\left(\frac{P}{P_0}\right) = \frac{10}{\ln 10} \ln\left(\frac{P}{P_0}\right), \qquad T(N) = \mathcal{O}(\log_2 N) \]
\noindent\textbf{Ví dụ tính toán thực tế:} Với cơ sở dữ liệu N = 1,000,000 bản ghi đã sắp xếp, tìm kiếm tuyến tính cần tối đa $10^6$ bước, trong khi tìm kiếm nhị phân chỉ cần $\lceil \log_2 (10^6$)$\rceil = \lceil$ 6 ln(10)/ln(2)$\rceil$  = 20 bước so sánh.
\end{ungdung}

\section{Bài tập Trắc nghiệm Củng cố Chương 1}
\begin{enumerate}[label=\textbf{Câu \arabic*:}, leftmargin=2.2em, itemsep=8pt]
  \item Tập xác định của hàm số f(x) = $\arcsin$(2x - 1) + $\ln$(x) là:
  \begin{itemize}[label={}, leftmargin=1em, itemsep=2pt]
    \item \textbf{A.} Đoạn [0, 1] \quad ($[0, 1]$)
    \item \textbf{B.} Nửa khoảng (0, 1] \quad ($(0, 1]$)
    \item \textbf{C.} Đoạn [-1, 1] \quad ($[-1, 1]$)
    \item \textbf{D.} Khoảng (0, 1/2] \quad ($(0, 1/2]$)
  \end{itemize}
  \vspace{2pt}
  \noindent{\color{teal!85!black}\textbf{Đáp án đúng: B.}} \textit{Giải thích:} Điều kiện của arcsin(2x - 1) là -1 $\le$  2x - 1 $\le$  1 $\iff$  0 $\le$  x $\le$  1. Kết hợp với điều kiện của ln(x) là x > 0, ta được x $\in$  (0, 1].
  \[ \begin{cases} -1 \le 2x - 1 \le 1 \\ x > 0 \end{cases} \iff \begin{cases} 0 \le x \le 1 \\ x > 0 \end{cases} \iff x \in (0, 1] \]
  \item Cho hàm số f(x) = $\ln$(x + $\sqrt{x^2 + 1}$). Khẳng định nào sau đây là ĐÚNG?
  \begin{itemize}[label={}, leftmargin=1em, itemsep=2pt]
    \item \textbf{A.} f(x) là hàm lẻ trên toàn trục số thực R.
    \item \textbf{B.} f(x) là hàm chẵn trên toàn trục số thực R.
    \item \textbf{C.} f(x) không chẵn cũng không lẻ.
    \item \textbf{D.} f(x) chỉ xác định khi x > 0.
  \end{itemize}
  \vspace{2pt}
  \noindent{\color{teal!85!black}\textbf{Đáp án đúng: A.}} \textit{Giải thích:} Vì $\sqrt{x^2 + 1} > \sqrt{x^2}$ = |x| $\ge$  -x nên x + $\sqrt{x^2 + 1}$ > 0 với mọi x $\in$  R. Khi đó f(-x) + f(x) = ln[(-x + $\sqrt{x^2+1}$)(x + $\sqrt{x^2+1}$)] = $ln(x^2$ + 1 - $x^2$) = ln(1) = 0, suy ra f(-x) = -f(x).
  \[ f(-x) = \ln\left(\sqrt{x^2+1} - x\right) = \ln\left(\frac{1}{\sqrt{x^2+1}+x}\right) = -f(x) \]
  \item Hàm ngược của hàm số y = f(x) = 3 + $\ln$(2x - 1) với x > 1/2 là:
  \begin{itemize}[label={}, leftmargin=1em, itemsep=2pt]
    \item \textbf{A.} $f^{-1}$(x) = (1 + $e^{x - 3}$) / 2 \quad ($f^{-1}(x) = \frac{1 + e^{x-3}}{2}$)
    \item \textbf{B.} $f^{-1}$(x) = $(e^{x + 3}$ - 1) / 2 \quad ($f^{-1}(x) = \frac{e^{x+3} - 1}{2}$)
    \item \textbf{C.} $f^{-1}$(x) = 1 / (3 + $\ln$(2x - 1)) \quad ($f^{-1}(x) = \frac{1}{3 + \ln(2x-1)}$)
    \item \textbf{D.} $f^{-1}$(x) = $e^{2x - 1}$ - 3 \quad ($f^{-1}(x) = e^{2x-1} - 3$)
  \end{itemize}
  \vspace{2pt}
  \noindent{\color{teal!85!black}\textbf{Đáp án đúng: A.}} \textit{Giải thích:} Từ y = 3 + ln(2x - 1) $\iff$  ln(2x - 1) = y - 3 $\iff$  2x - 1 = $e^{y - 3} \iff$  x = (1 + $e^{y - 3}$)/2. Đổi chỗ x và y ta có đáp án A.
  \[ y - 3 = \ln(2x-1) \iff 2x - 1 = e^{y-3} \implies f^{-1}(x) = \frac{1 + e^{x-3}}{2} \]
\end{enumerate}

% ======================================================================
% CHƯƠNG 2: GIỚI HẠN VÀ ĐẠO HÀM (Buổi 1 - 2)
% ======================================================================
\chapter{Giới hạn và Đạo hàm}
\label{chap:chuong-2}

\begin{tcolorbox}[colback=slate!5!white,colframe=slate!75!black,title={\textbf{Mục lục Chương 2 \& Chuẩn đầu ra (Buổi 1 - 2)}}]
\textbf{Tổng quan:} Chương nền tảng giải quyết hai bài toán cổ điển dẫn tới phép tính vi phân: Bài toán tiếp tuyến và Bài toán vận tốc tức thời. Trình bày từ trực quan đến định nghĩa chính xác bằng ngôn ngữ $\varepsilon - \delta$ , tính liên tục, giới hạn tại vô cùng và định nghĩa đạo hàm như một hàm số.\\[6pt]
\textbf{Nội dung chính trong chương:}
\begin{itemize}[leftmargin=1.5em, itemsep=2pt]
  \item \textbf{Mục 2.1 - 2.3:} Bài toán tiếp tuyến, Vận tốc \& Các quy tắc tính giới hạn
  \item \textbf{Mục 2.4:} Định nghĩa chính xác của giới hạn (Ngôn ngữ $\varepsilon - \delta$ )
  \item \textbf{Mục 2.5 - 2.6:} Tính liên tục của hàm số \& Giới hạn tại vô cùng (Tiệm cận)
  \item \textbf{Mục 2.7 - 2.8:} Đạo hàm tại một điểm \& Đạo hàm như một hàm số (Trọng tâm)
\end{itemize}
\textbf{Chuẩn đầu ra (LLO):}
\begin{itemize}[leftmargin=1.5em, itemsep=2pt]
  \item LLO2.1: Giải thích được định nghĩa chính xác về giới hạn thông qua ngôn ngữ epsilon - delta ($\varepsilon - \delta$ ).
  \item LLO2.2: Giải được các bài toán tính giới hạn tại một điểm, giới hạn tại vô cùng và xét tính liên tục của hàm số.
  \item LLO2.3: Vận dụng được định nghĩa đạo hàm để giải bài toán tiếp tuyến, vận tốc tức thời và tốc độ thay đổi.
\end{itemize}
\end{tcolorbox}

\section{[Mục 2.1 - 2.3] Bài toán tiếp tuyến, Vận tốc \& Các quy tắc tính giới hạn}
Hệ số góc của tiếp tuyến là giới hạn của hệ số góc dây cung (cát tuyến) khi điểm thứ hai tiến dần về điểm tiếp xúc. Tương tự, vận tốc tức thời là giới hạn của vận tốc trung bình khi khoảng thời gian $\Delta$ t tiến về 0.

\begin{congthuc}{Các công thức trọng tâm (2.1 - 2.3)}
\noindent\textbf{Giới hạn một phía \& Điều kiện tồn tại giới hạn:}
\[ \lim_{x \to a} f(x) = L \iff \lim_{x \to a^-} f(x) = \lim_{x \to a^+} f(x) = L \]
\noindent\textit{Chú giải:} Nếu giới hạn trái khác giới hạn phải thì giới hạn tại x = a không tồn tại.\\[4pt]
\noindent\textbf{Định lý Kẹp (Squeeze Theorem):}
\[ g(x) \le f(x) \le h(x) \text{ và } \lim_{x \to a} g(x) = \lim_{x \to a} h(x) = L \implies \lim_{x \to a} f(x) = L \]
\noindent\textit{Chú giải:} Thường áp dụng khi xuất hiện sin(1/x) hoặc cos(1/x) nhân với đại lượng tiến về 0.\\[4pt]
\end{congthuc}

\noindent\textbf{Chú giải \& Lưu ý kỹ thuật:}
\begin{itemize}[leftmargin=1.5em, itemsep=2pt]
  \item Để khử dạng vô định 0/0 của hàm phân thức hoặc căn thức: Phân tích đa thức thành nhân tử, nhân lượng liên hợp, hoặc sử dụng giới hạn cơ bản sin(x)/x $\to$  1.
\end{itemize}

\begin{vidu}{Khử dạng vô định 0/0 kết hợp lượng liên hợp và Định lý Kẹp}
\textbf{Đề bài:} Tính các giới hạn sau: (a) $I_1$ = \$\lim_{x \to 0} \frac{\sqrt{1 + 3x} - 1}{x}$ và (b) $I_2$ = \$\lim_{x \to 0} x^2 \cos\left(\frac{1}{x}\right)$.
\[ I_1 = \lim_{x \to 0} \frac{\sqrt{1 + 3x} - 1}{x}, \qquad I_2 = \lim_{x \to 0} x^2 \cos\left(\frac{1}{x}\right) \]
\textbf{Lời giải chi tiết:}
\begin{enumerate}[leftmargin=1.5em, itemsep=2pt]
  \item Câu (a): Nhân cả tử và mẫu với biểu thức liên hợp ($\sqrt{1 + 3x}$ + 1) để triệt tiêu nhân tử x.
  \item Câu (b): Đánh giá -1 $\le$  cos(1/x) $\le$  1 với mọi x $\neq$  0, nhân hai vế với $x^2 \ge$  0 và áp dụng Định lý Kẹp.
\end{enumerate}
\[ I_1 = \lim_{x \to 0} \frac{(1 + 3x) - 1}{x(\sqrt{1 + 3x} + 1)} = \lim_{x \to 0} \frac{3}{\sqrt{1 + 3x} + 1} = \frac{3}{1 + 1} = \frac{3}{2} \]
\[ -x^2 \le x^2 \cos\left(\frac{1}{x}\right) \le x^2, \text{ mà } \lim_{x \to 0} (\pm x^2) = 0 \implies I_2 = 0 \]
\noindent\textbf{Kết luận:} Kết quả: $I_1$ = 3/2 và $I_2$ = 0.
\end{vidu}

\section{[Mục 2.4] Định nghĩa chính xác của giới hạn (Ngôn ngữ $\varepsilon - \delta$ )}
Ta nói giới hạn của f(x) khi x tiến tới a bằng L nếu với mọi sai số cho phép $\varepsilon$  > 0 (dù nhỏ tùy ý), luôn tìm được bán kính lân cận $\delta$  > 0 sao cho mọi x thỏa mãn 0 < |x - a| < $\delta$  đều kéo theo |f(x) - L| < $\varepsilon$ .

\begin{congthuc}{Các công thức trọng tâm (2.4)}
\noindent\textbf{Định nghĩa hình thức Cauchy ($\varepsilon - \delta$ ):}
\[ \lim_{x \to a} f(x) = L \iff \forall \varepsilon > 0, \, \exists \delta = \delta(\varepsilon) > 0 : \forall x \in D, \; 0 < |x - a| < \delta \implies |f(x) - L| < \varepsilon \]
\noindent\textit{Chú giải:} Điều kiện 0 < |x - a| có nghĩa là x $\neq$  a (giá trị tại x = a không ảnh hưởng đến giới hạn).\\[4pt]
\end{congthuc}

\noindent\textbf{Chú giải \& Lưu ý kỹ thuật:}
\begin{itemize}[leftmargin=1.5em, itemsep=2pt]
  \item Phương pháp chứng minh bằng $\varepsilon - \delta$  gồm 2 giai đoạn: (1) Phân tích sơ bộ (Scratch work) đánh giá |f(x) - L| theo |x - a| để tìm $\delta$  theo $\varepsilon$ ; (2) Trình bày chứng minh chính thức.
\end{itemize}

\begin{vidu}{Chứng minh giới hạn bằng định nghĩa $\varepsilon - \delta$ }
\textbf{Đề bài:} Dùng định nghĩa $\varepsilon - \delta$ để chứng minh rằng \$\lim_{x \to 3}$ (4x - 5) = 7.
\[ \lim_{x \to 3} (4x - 5) = 7 \]
\textbf{Lời giải chi tiết:}
\begin{enumerate}[leftmargin=1.5em, itemsep=2pt]
  \item Bước 1 (Phân tích): Biến đổi |f(x) - L| = |(4x - 5) - 7| = |4x - 12| = 4|x - 3|.
  \item Bước 2 (Chọn $\delta$ ): Để 4|x - 3| < $\varepsilon$ , ta chỉ cần |x - 3| < $\varepsilon$ /4. Vậy chọn $\delta = \varepsilon$ /4.
  \item Bước 3 (Chứng minh): Với mọi $\varepsilon$  > 0, chọn $\delta = \varepsilon$ /4 > 0, suy ra điều phải chứng minh.
\end{enumerate}
\[ \forall \varepsilon > 0, \text{ chọn } \delta = \frac{\varepsilon}{4} > 0. \]
\[ \text{Với mọi } x \text{ thỏa mãn } 0 < |x - 3| < \delta \implies |(4x - 5) - 7| = 4|x - 3| < 4\delta = 4\left(\frac{\varepsilon}{4}\right) = \varepsilon. \]
\noindent\textbf{Kết luận:} Vậy theo định nghĩa $\varepsilon - \delta$ , ta có $\lim_{x \to  3}$ (4x - 5) = 7.
\end{vidu}

\section{[Mục 2.5 - 2.6] Tính liên tục của hàm số \& Giới hạn tại vô cùng (Tiệm cận)}
Hàm số f liên tục tại x = a nếu giới hạn của hàm khi x tiến tới a bằng đúng giá trị hàm tại a. Giới hạn tại vô cùng mô tả dáng điệu của hàm số khi x $\to \pm \infty$ và xác định các đường tiệm cận ngang.

\begin{congthuc}{Các công thức trọng tâm (2.5 - 2.6)}
\noindent\textbf{Điều kiện liên tục tại một điểm \& Định lý Giá trị Trung gian (IVT):}
\[ \lim_{x \to a} f(x) = f(a) \iff \lim_{x \to a^-} f(x) = \lim_{x \to a^+} f(x) = f(a) \]
\noindent\textit{Chú giải:} Hệ quả IVT: Nếu f liên tục trên [a, b] và f(a)$\cdot$ f(b) < 0 thì tồn tại c $\in$  (a, b) sao cho f(c) = 0.\\[4pt]
\noindent\textbf{Đường tiệm cận đứng và Tiệm cận ngang:}
\[ \lim_{x \to a^\pm} f(x) = \pm\infty \implies x = a \text{ (TCĐ)}, \qquad \lim_{x \to \pm\infty} f(x) = L \implies y = L \text{ (TCN)} \]
\noindent\textit{Chú giải:} Chú ý: Khi x $\to - \infty$, ta có $\sqrt{x^2}$ = |x| = -x.\\[4pt]
\end{congthuc}

\noindent\textbf{Chú giải \& Lưu ý kỹ thuật:}
\begin{itemize}[leftmargin=1.5em, itemsep=2pt]
  \item Ba loại điểm gián đoạn thường gặp: Gián đoạn bỏ được (removable - tồn tại giới hạn hữu hạn nhưng khác f(a)), gián đoạn bước nhảy (jump - giới hạn trái khác giới hạn phải), và gián đoạn vô cực.
\end{itemize}

\begin{vidu}{Tìm tham số để hàm ghép liên tục và tìm tiệm cận ngang chứa căn thức}
\textbf{Đề bài:} Tìm hằng số k để hàm số f(x) = $\begin{cases} \frac{\sin(3x)}{x} & \text{khi } x \neq 0 \\ 2k - 1 & \text{khi } x = 0 \end{cases}$ liên tục tại x = 0.
\[ \lim_{x \to 0} \frac{\sin(3x)}{x} = f(0) = 2k - 1 \]
\textbf{Lời giải chi tiết:}
\begin{enumerate}[leftmargin=1.5em, itemsep=2pt]
  \item Bước 1: Tính giới hạn của f(x) khi x $\to$  0 bằng cách nhân và chia cho 3.
  \item Bước 2: Cho giá trị giới hạn bằng f(0) = 2k - 1 để giải tìm k.
\end{enumerate}
\[ \lim_{x \to 0} f(x) = \lim_{x \to 0} 3 \cdot \frac{\sin(3x)}{3x} = 3 \cdot 1 = 3 \]
\[ \text{Để } f \text{ liên tục tại } x = 0 \iff 2k - 1 = 3 \iff k = 2 \]
\noindent\textbf{Kết luận:} Với k = 2 thì hàm số liên tục tại x = 0 (và liên tục trên toàn bộ R).
\end{vidu}

\section{[Mục 2.7 - 2.8] Đạo hàm tại một điểm \& Đạo hàm như một hàm số (Trọng tâm)}
Đạo hàm f'(a) là tốc độ thay đổi tức thời của y = f(x) đối với x tại điểm x = a, đồng thời là hệ số góc của tiếp tuyến với đồ thị tại điểm P(a, f(a)). Khi cho a thay đổi, f'(x) trở thành một hàm số mới.

\begin{congthuc}{Các công thức trọng tâm (2.7 - 2.8)}
\noindent\textbf{Định nghĩa Đạo hàm qua giới hạn số gia:}
\[ f'(a) = \lim_{h \to 0} \frac{f(a + h) - f(a)}{h} = \lim_{x \to a} \frac{f(x) - f(a)}{x - a} \]
\noindent\textit{Chú giải:} Phương trình tiếp tuyến tại x = a: y = f(a) + f'(a)(x - a).\\[4pt]
\noindent\textbf{Quan hệ giữa Tính khả vi và Tính liên tục:}
\[ f \text{ khả vi tại } a \implies f \text{ liên tục tại } a \quad (\text{Chiều ngược lại KHÔNG đúng!}) \]
\noindent\textit{Chú giải:} Ví dụ kinh điển: Hàm f(x) = |x| liên tục tại x = 0 nhưng không có đạo hàm tại x = 0 (điểm góc).\\[4pt]
\end{congthuc}

\noindent\textbf{Chú giải \& Lưu ý kỹ thuật:}
\begin{itemize}[leftmargin=1.5em, itemsep=2pt]
  \item Hàm số không khả vi tại x = a trong 3 trường hợp hình học: (1) Đồ thị có điểm góc nhọn (corner/cusp, đạo hàm trái $\neq$  đạo hàm phải); (2) Đồ thị bị gián đoạn tại a; (3) Đồ thị có tiếp tuyến thẳng đứng (lim |f'(x)| = $\infty$).
\end{itemize}

\begin{vidu}{Tính đạo hàm bằng định nghĩa và viết phương trình tiếp tuyến}
\textbf{Đề bài:} Dùng định nghĩa giới hạn để tìm hàm đạo hàm f'(x) của f(x) = $\sqrt{x + 1}$ (với x > -1) và viết phương trình tiếp tuyến tại điểm có hoành độ $x_0$ = 3.
\[ f'(x) = \lim_{h \to 0} \frac{\sqrt{x + h + 1} - \sqrt{x + 1}}{h} \]
\textbf{Lời giải chi tiết:}
\begin{enumerate}[leftmargin=1.5em, itemsep=2pt]
  \item Bước 1: Thiết lập tỉ số số gia [f(x + h) - f(x)] / h và nhân lượng liên hợp của tử thức.
  \item Bước 2: Rút gọn h ở tử và mẫu rồi cho h $\to$  0 để thu được f'(x).
  \item Bước 3: Thay $x_0$ = 3 để tính f(3) = 2, f'(3) = 1/4 và lập phương trình tiếp tuyến.
\end{enumerate}
\[ f'(x) = \lim_{h \to 0} \frac{(x + h + 1) - (x + 1)}{h\left(\sqrt{x + h + 1} + \sqrt{x + 1}\right)} = \lim_{h \to 0} \frac{1}{\sqrt{x + h + 1} + \sqrt{x + 1}} = \frac{1}{2\sqrt{x + 1}} \]
\[ \text{Tại } x_0 = 3: f(3) = 2, \; f'(3) = \frac{1}{4} \implies y = 2 + \frac{1}{4}(x - 3) \iff y = \frac{1}{4}x + \frac{5}{4} \]
\noindent\textbf{Kết luận:} Hàm đạo hàm là f'(x) = 1/(2$\sqrt{x+1}$) và tiếp tuyến tại (3, 2) là y = x/4 + 5/4.
\end{vidu}

\section{Minh họa Hình học bằng TikZ}
\begin{figure}[htbp]
\centering
\begin{tikzpicture}[scale=1.1, >=stealth]
  % Trục tọa độ
  \draw[->, thick] (-0.5,0) -- (5.2,0) node[right] {$x$};
  \draw[->, thick] (0,-0.5) -- (0,4.2) node[above] {$y$};
  
  % Dải băng epsilon và delta
  \fill[blue!8] (0, 1.8) rectangle (4.8, 2.8);
  \fill[red!10] (2.1, 0) rectangle (3.1, 3.8);
  \fill[purple!20] (2.1, 1.8) rectangle (3.1, 2.8);
  
  % Đồ thị hàm số y = f(x)
  \draw[domain=0.4:4.5, smooth, variable=\x, blue!80!black, very thick] 
    plot ({\x}, {0.22*(\x)^2 + 0.8}) node[right] {$y = f(x)$};
    
  % Các đường gióng tại x = a, y = L
  \draw[dashed, thick, black!70] (2.6,0) node[below] {$a$} -- (2.6,2.287) -- (0,2.287) node[left] {$L$};
  \draw[dotted, red!80!black] (2.1,0) node[below left, scale=0.8] {$a-\delta$} -- (2.1,2.8);
  \draw[dotted, red!80!black] (3.1,0) node[below right, scale=0.8] {$a+\delta$} -- (3.1,2.8);
  \draw[dotted, blue!80!black] (0,1.8) node[left, scale=0.8] {$L-\varepsilon$} -- (4.8,1.8);
  \draw[dotted, blue!80!black] (0,2.8) node[left, scale=0.8] {$L+\varepsilon$} -- (4.8,2.8);
  
  % Điểm giới hạn (a, L)
  \filldraw[fill=white, draw=blue!80!black, very thick] (2.6,2.287) circle (2.5pt);
\end{tikzpicture}
\caption{Hình 2.1: Minh họa hình học định nghĩa giới hạn $\varepsilon - \delta$  và sự hội tụ của cát tuyến PQ về tiếp tuyến tại P.}
\label{fig:tikz-ch2-epsilon-delta}
\end{figure}

\section{Ứng dụng Thực tế \& Tối ưu hóa (Cơ điện tử \& Điều khiển Tự động)}
\begin{ungdung}{Cảm biến Vận tốc Tức thời và Khâu Vi phân (D) trong Bộ điều khiển PID}
Trong xe tự hành và robot công nghiệp, bộ mã hóa vòng quay (Encoder) đo vị trí rời rạc $s(t_k$). Bộ điều khiển PID tính xấp xỉ đạo hàm de(t)/dt của sai số vị trí để dự đoán xu hướng vượt quá (overshoot) và tạo lực hãm sớm giúp hệ thống ổn định nhanh.
\[ u(t) = K_p e(t) + K_i \int_0^t e(\tau)\,d\tau + K_d \frac{de(t)}{dt}, \qquad v(t) = \lim_{\Delta t \to 0} \frac{s(t + \Delta t) - s(t)}{\Delta t} \]
\noindent\textbf{Ví dụ tính toán thực tế:} Nếu phương trình chuyển động của drone hạ cánh là h(t) = 20 - 1.$5t^2$ (m), thì vận tốc chạm đất tại t = 2s là v(2) = h'(2) = -3(2) = -6 m/s (hướng xuống).
\end{ungdung}

\section{Bài tập Trắc nghiệm Củng cố Chương 2}
\begin{enumerate}[label=\textbf{Câu \arabic*:}, leftmargin=2.2em, itemsep=8pt]
  \item Giá trị của giới hạn L = \$\lim_{x \to -\infty} \frac{\sqrt{4x^2 + 3x}}{2x - 1}$ bằng:
  \begin{itemize}[label={}, leftmargin=1em, itemsep=2pt]
    \item \textbf{A.} 1 \quad ($1$)
    \item \textbf{B.} -1 \quad ($-1$)
    \item \textbf{C.} 2 \quad ($2$)
    \item \textbf{D.} 0 \quad ($0$)
  \end{itemize}
  \vspace{2pt}
  \noindent{\color{teal!85!black}\textbf{Đáp án đúng: B.}} \textit{Giải thích:} Khi x $\to - \infty$, ta có $\sqrt{x^2}$ = |x| = -x. Do đó $\sqrt{4x^2 + 3x}$ = -x$\sqrt{4 + 3/x}$. Chia cả tử và mẫu cho x, ta được lim [-$\sqrt{4 + 3/x}$ / (2 - 1/x)] = -2/2 = -1.
  \[ \lim_{x \to -\infty} \frac{-x\sqrt{4 + 3/x}}{x(2 - 1/x)} = \frac{-\sqrt{4}}{2} = -1 \]
  \item Cho hàm số f(x) = |x - 2| tại điểm x = 2. Phát biểu nào sau đây là ĐÚNG?
  \begin{itemize}[label={}, leftmargin=1em, itemsep=2pt]
    \item \textbf{A.} Hàm số liên tục tại x = 2 nhưng không khả vi tại x = 2.
    \item \textbf{B.} Hàm số vừa liên tục vừa khả vi tại x = 2 với f'(2) = 0.
    \item \textbf{C.} Hàm số bị gián đoạn tại x = 2.
    \item \textbf{D.} Hàm số khả vi tại x = 2 nhưng không liên tục.
  \end{itemize}
  \vspace{2pt}
  \noindent{\color{teal!85!black}\textbf{Đáp án đúng: A.}} \textit{Giải thích:} Ta có $\lim_{x\to 2}$ |x - 2| = 0 = f(2) nên hàm liên tục tại x = 2. Tuy nhiên đạo hàm bên phải $f'_+$(2) = 1 khác đạo hàm bên trái $f'_-$(2) = -1 nên f'(2) không tồn tại.
  \[ f'_+(2) = \lim_{h \to 0^+} \frac{h}{h} = 1 \neq f'_-(2) = \lim_{h \to 0^-} \frac{-h}{h} = -1 \]
  \item Giới hạn \$\lim_{h \to 0} \frac{\ln(e + h) - 1}{h}$ biểu thị giá trị đạo hàm nào sau đây?
  \begin{itemize}[label={}, leftmargin=1em, itemsep=2pt]
    \item \textbf{A.} Đạo hàm của f(x) = ln(x) tại x = e, bằng 1/e. \quad ($f'(e) = \frac{1}{e}$)
    \item \textbf{B.} Đạo hàm của f(x) = ln(x) tại x = 1, bằng 1. \quad ($f'(1) = 1$)
    \item \textbf{C.} Đạo hàm của f(x) = $e^x$ tại x = 1, bằng e. \quad ($f'(1) = e$)
    \item \textbf{D.} Giới hạn bằng 0. \quad ($0$)
  \end{itemize}
  \vspace{2pt}
  \noindent{\color{teal!85!black}\textbf{Đáp án đúng: A.}} \textit{Giải thích:} Vì 1 = ln(e), biểu thức đã cho chính là định nghĩa đạo hàm $\lim_{h\to 0} [ln(e + h) - ln(e)]$ / h = (ln $x)'|_{x=e}$ = 1/e.
  \[ \lim_{h \to 0} \frac{\ln(e+h) - \ln e}{h} = \left.(\ln x)'\right|_{x=e} = \frac{1}{e} \]
\end{enumerate}

% ======================================================================
% CHƯƠNG 3: CÁC QUY TẮC TÍNH ĐẠO HÀM (TRỌNG TÂM) (Buổi 3)
% ======================================================================
\chapter{Các quy tắc tính Đạo hàm (Trọng tâm)}
\label{chap:chuong-3}

\begin{tcolorbox}[colback=slate!5!white,colframe=slate!75!black,title={\textbf{Mục lục Chương 3 \& Chuẩn đầu ra (Buổi 3)}}]
\textbf{Tổng quan:} Chương trọng tâm cung cấp hệ thống công cụ giải tích mạnh mẽ để tính đạo hàm của mọi hàm sơ cấp và tổ hợp của chúng mà không cần quay lại tính giới hạn: Quy tắc tích/thương, Quy tắc dây chuyền (hàm hợp), Đạo hàm hàm ẩn, Đạo hàm logarithm, Hàm hyperbolic, Đạo hàm cấp cao và Xấp xỉ tuyến tính - Vi phân.\\[6pt]
\textbf{Nội dung chính trong chương:}
\begin{itemize}[leftmargin=1.5em, itemsep=2pt]
  \item \textbf{Mục 3.1 - 3.3:} Đạo hàm đa thức, hàm mũ, tích, thương và hàm lượng giác
  \item \textbf{Mục 3.4 - 3.5:} Quy tắc hàm hợp (Chain Rule) \& Đạo hàm của hàm ẩn
  \item \textbf{Mục 3.6 - 3.7:} Đạo hàm hàm Logarithm, Đạo hàm dạng $u(x)^v$(x) \& Hàm Hyperbolic
  \item \textbf{Mục 3.8 - 3.9:} Đạo hàm bậc cao (Công thức Leibniz) \& Xấp xỉ tuyến tính - Vi phân
\end{itemize}
\textbf{Chuẩn đầu ra (LLO):}
\begin{itemize}[leftmargin=1.5em, itemsep=2pt]
  \item LLO3.1: Tính thành thạo đạo hàm của hàm đa thức, hàm mũ, tích/thương, lượng giác, hàm hợp, hàm ẩn, hàm logarithm, hàm hyperbolic và đạo hàm cấp cao.
  \item LLO3.2: Thiết lập được xấp xỉ tuyến tính của hàm tại một điểm và vận dụng vi phân để ước lượng sai số tối đa trong đo lường kỹ thuật.
\end{itemize}
\end{tcolorbox}

\section{[Mục 3.1 - 3.3] Đạo hàm đa thức, hàm mũ, tích, thương và hàm lượng giác}
Hàm số mũ tự nhiên f(x) = $e^x$ có tính chất đặc biệt duy nhất là đạo hàm của nó bằng chính nó: $(e^x$)' = $e^x$. Kết hợp với quy tắc đạo hàm của tổng, hiệu, tích, thương và các hàm lượng giác, ta dễ dàng khảo sát tốc độ biến thiên của các hàm cơ bản.

\begin{congthuc}{Các công thức trọng tâm (3.1 - 3.3)}
\noindent\textbf{Quy tắc Đạo hàm của Tích và Thương:}
\[ \frac{d}{dx}[f(x)g(x)] = f'(x)g(x) + f(x)g'(x), \qquad \frac{d}{dx}\left[\frac{f(x)}{g(x)}\right] = \frac{f'(x)g(x) - f(x)g'(x)}{[g(x)]^2} \]
\noindent\textit{Chú giải:} Tuyệt đối tránh sai lầm phổ biến: (f$\cdot$ g)' $\neq$  f'$\cdot$ g' và (f/g)' $\neq$  f'/g'.\\[4pt]
\noindent\textbf{Đạo hàm 6 hàm lượng giác chuẩn:}
\[ (\sin x)' = \cos x, \; (\cos x)' = -\sin x, \; (\tan x)' = \sec^2 x, \; (\cot x)' = -\csc^2 x, \; (\sec x)' = \sec x\tan x \]
\noindent\textit{Chú giải:} Trong đó sec x = 1/cos x và csc x = 1/sin x.\\[4pt]
\end{congthuc}

\noindent\textbf{Chú giải \& Lưu ý kỹ thuật:}
\begin{itemize}[leftmargin=1.5em, itemsep=2pt]
  \item Tiếp tuyến nằm ngang (Horizontal tangent) xuất hiện tại các điểm mà đạo hàm f'(x) = 0.
\end{itemize}

\begin{vidu}{Vận dụng quy tắc tích và thương tìm tiếp tuyến song song}
\textbf{Đề bài:} Tính đạo hàm của hàm số y = $\frac{x^2 e^x}{x + 1}$ tại điểm x = 1.
\[ y = \frac{x^2 e^x}{x + 1} \]
\textbf{Lời giải chi tiết:}
\begin{enumerate}[leftmargin=1.5em, itemsep=2pt]
  \item Bước 1: Đặt tử số u(x) = $x^2 e^x$, tính u'(x) bằng quy tắc tích: u'(x) = (2x + $x^2 )e^x$.
  \item Bước 2: Áp dụng quy tắc đạo hàm của thương với mẫu v(x) = x + 1, v'(x) = 1.
  \item Bước 3: Rút gọn biểu thức và thay x = 1.
\end{enumerate}
\[ y' = \frac{[(2x + x^2)e^x](x + 1) - x^2 e^x (1)}{(x + 1)^2} = \frac{x e^x (x^2 + 2x + 2)}{(x + 1)^2} \]
\[ y'(1) = \frac{1 \cdot e^1 (1 + 2 + 2)}{(1 + 1)^2} = \frac{5e}{4} \]
\noindent\textbf{Kết luận:} Giá trị đạo hàm tại x = 1 là y'(1) = 5e/4.
\end{vidu}

\section{[Mục 3.4 - 3.5] Quy tắc hàm hợp (Chain Rule) \& Đạo hàm của hàm ẩn}
Quy tắc dây chuyền (Chain Rule) cho phép tính đạo hàm của hàm hợp F(x) = f(g(x)) bằng tích đạo hàm hàm ngoài (tính tại hàm trong) với đạo hàm hàm trong. Khi y không được biểu diễn tường minh theo x mà qua phương trình F(x, y) = 0, ta lấy đạo hàm hai vế theo x (coi y = y(x) là hàm hợp) rồi giải tìm y'.

\begin{congthuc}{Các công thức trọng tâm (3.4 - 3.5)}
\noindent\textbf{Quy tắc dây chuyền (Ký hiệu Lagrange và Leibniz):}
\[ F'(x) = f'(g(x)) \cdot g'(x) \qquad \text{hoặc} \qquad \frac{dy}{dx} = \frac{dy}{du} \cdot \frac{du}{dx} \]
\noindent\textit{Chú giải:} Với hàm lũy thừa tổng quát: $([u(x)]^n$)' = $n[u(x)]^{n-1} \cdot$  u'(x).\\[4pt]
\noindent\textbf{Công thức tính nhanh Đạo hàm hàm ẩn F(x, y) = 0:}
\[ \frac{d}{dx}[F(x, y(x))] = 0 \implies y'(x) = -\frac{F'_x(x, y)}{F'_y(x, y)} \quad (F'_y \neq 0) \]
\noindent\textit{Chú giải:} Khi đạo hàm các số hạng chứa y theo x, luôn phải nhân thêm y' (ví dụ: $(y^3 )'_x = 3y^2 \cdot$  y').\\[4pt]
\end{congthuc}

\noindent\textbf{Chú giải \& Lưu ý kỹ thuật:}
\begin{itemize}[leftmargin=1.5em, itemsep=2pt]
  \item Quy tắc đạo hàm hàm ẩn cũng là phương pháp chuẩn để chứng minh công thức đạo hàm của các hàm lượng giác ngược như (arcsin x)' = 1/$\sqrt{1 - x^2}$.
\end{itemize}

\begin{vidu}{Đạo hàm hàm ẩn trên đường cong Lá Descartes (Folium of Descartes)}
\textbf{Đề bài:} Cho đường cong xác định bởi phương trình hàm ẩn $x^3 + y^3$ = 6xy. Hãy tính đạo hàm y' = dy/dx và viết phương trình tiếp tuyến tại điểm M(3, 3).
\[ x^3 + y^3 = 6xy \quad \text{tại } M(3, 3) \]
\textbf{Lời giải chi tiết:}
\begin{enumerate}[leftmargin=1.5em, itemsep=2pt]
  \item Bước 1: Đạo hàm cả hai vế của phương trình theo biến x (chú ý vế phải 6xy là tích của hai hàm theo x).
  \item Bước 2: Nhóm các số hạng chứa y' sang một vế và rút y' theo x, y.
  \item Bước 3: Thay tọa độ x = 3, y = 3 vào biểu thức y' và viết phương trình tiếp tuyến.
\end{enumerate}
\[ 3x^2 + 3y^2 y' = 6y + 6x y' \iff (y^2 - 2x)y' = 2y - x^2 \implies y' = \frac{2y - x^2}{y^2 - 2x} \]
\[ \text{Tại } M(3, 3): y'(3, 3) = \frac{2(3) - 3^2}{3^2 - 2(3)} = \frac{6 - 9}{9 - 6} = -1 \]
\[ \text{Phương trình tiếp tuyến: } y - 3 = -1(x - 3) \iff x + y - 6 = 0 \]
\noindent\textbf{Kết luận:} Hệ số góc tiếp tuyến tại M(3, 3) là -1 và phương trình tiếp tuyến là x + y - 6 = 0.
\end{vidu}

\section{[Mục 3.6 - 3.7] Đạo hàm hàm Logarithm, Đạo hàm dạng $u(x)^v$(x) \& Hàm Hyperbolic}
Phương pháp lấy đạo hàm logarithm (Logarithmic Differentiation) giúp đơn giản hóa việc tính đạo hàm của các biểu thức tích/thương phức tạp hoặc hàm lũy thừa - mũ dạng y = $[u(x)]^{v(x)}$ bằng cách lấy ln hai vế trước khi đạo hàm.

\begin{congthuc}{Các công thức trọng tâm (3.6 - 3.7)}
\noindent\textbf{Đạo hàm Logarithm và Hàm mũ-lũy thừa $u(x)^v$(x):}
\[ (\ln|u(x)|)' = \frac{u'(x)}{u(x)}, \qquad \left(u^v\right)' = u^v \left(v' \ln u + v \frac{u'}{u}\right) \]
\noindent\textit{Chú giải:} Bản chất: Viết $u^v = e^{v ln u}$ rồi áp dụng quy tắc hàm hợp.\\[4pt]
\noindent\textbf{Định nghĩa và Đạo hàm các hàm Hyperbolic:}
\[ \sinh x = \frac{e^x - e^{-x}}{2}, \; \cosh x = \frac{e^x + e^{-x}}{2}, \; \cosh^2 x - \sinh^2 x = 1, \; (\sinh x)' = \cosh x, \; (\cosh x)' = \sinh x \]
\noindent\textit{Chú giải:} Đường dây cáp treo tự do (dây truyền tải điện) có hình dạng chính xác là đồ thị hàm y = a cosh(x/a) (đường dây xích - Catenary).\\[4pt]
\end{congthuc}

\noindent\textbf{Chú giải \& Lưu ý kỹ thuật:}
\begin{itemize}[leftmargin=1.5em, itemsep=2pt]
  \item Hàm ngược của sinh x là argsinh(x) = ln(x + $\sqrt{x^2 + 1}$), có đạo hàm là 1/$\sqrt{x^2 + 1}$.
\end{itemize}

\begin{vidu}{Tính đạo hàm của hàm lũy thừa - mũ bằng phương pháp Logarithm}
\textbf{Đề bài:} Tính đạo hàm của hàm số y = $x^{\sin x}$ với x > 0.
\[ y = x^{\sin x} \quad (x > 0) \]
\textbf{Lời giải chi tiết:}
\begin{enumerate}[leftmargin=1.5em, itemsep=2pt]
  \item Bước 1: Lấy logarithm tự nhiên cả hai vế: ln y = $ln(x^{sin x}$) = sin(x) $\cdot$  ln(x).
  \item Bước 2: Đạo hàm hai vế theo x (vế trái là hàm ẩn theo y, vế phải là tích hai hàm).
  \item Bước 3: Nhân hai vế với y = $x^{sin x}$ để thu được kết quả cuối cùng.
\end{enumerate}
\[ \ln y = \sin x \cdot \ln x \implies \frac{y'}{y} = \cos x \cdot \ln x + \frac{\sin x}{x} \]
\[ y' = x^{\sin x} \left( \cos x \ln x + \frac{\sin x}{x} \right) \]
\noindent\textbf{Kết luận:} Đạo hàm của hàm số là y' = $x^{\sin x}$ (cos x ln x + sin x / x).
\end{vidu}

\section{[Mục 3.8 - 3.9] Đạo hàm bậc cao (Công thức Leibniz) \& Xấp xỉ tuyến tính - Vi phân}
Đạo hàm cấp n được định nghĩa quy nạp: $f^{(n)}$(x) = $[f^{(n-1)}$(x)]'. Trong kỹ thuật, đường tiếp tuyến L(x) tại x = a được dùng để xấp xỉ tuyến tính hàm f(x) trong lân cận điểm a, và vi phân dy = f'(x)dx dùng để ước lượng sai số lan truyền.

\begin{congthuc}{Các công thức trọng tâm (3.8 - 3.9)}
\noindent\textbf{Đạo hàm cấp n cơ bản \& Công thức Leibniz:}
\[ (\sin x)^{(n)} = \sin\left(x + \frac{n\pi}{2}\right), \quad \left(\frac{1}{x+a}\right)^{(n)} = \frac{(-1)^n n!}{(x+a)^{n+1}}, \quad (uv)^{(n)} = \sum_{k=0}^n C_n^k u^{(n-k)} v^{(k)} \]
\noindent\textit{Chú giải:} Công thức Leibniz cực kỳ hiệu quả khi một trong hai hàm là đa thức bậc thấp (vì đạo hàm từ cấp m+1 trở đi bằng 0).\\[4pt]
\noindent\textbf{Xấp xỉ tuyến tính L(x) và Vi phân hàm một biến:}
\[ f(x) \approx L(x) = f(a) + f'(a)(x - a), \qquad \Delta y \approx dy = f'(x)\,dx \]
\noindent\textit{Chú giải:} Sai số tương đối trong đo lường: $\Delta$ y / y $\approx$  dy / y = d(ln y).\\[4pt]
\end{congthuc}

\noindent\textbf{Chú giải \& Lưu ý kỹ thuật:}
\begin{itemize}[leftmargin=1.5em, itemsep=2pt]
  \item Khi áp dụng công thức Leibniz cho tích P(x)$\cdot$ g(x) với P(x) là đa thức bậc 2, tổng chuỗi Leibniz chỉ còn đúng 3 số hạng ứng với k = 0, 1, 2.
\end{itemize}

\begin{vidu}{Ứng dụng xấp xỉ tuyến tính và tính đạo hàm cấp n bằng Leibniz}
\textbf{Đề bài:} (a) Tính đạo hàm cấp 50 của hàm y = $x^2 e^{2x}$. (b) Dùng vi phân tính gần đúng giá trị $\sqrt[3]{8.06}$.
\[ y = x^2 e^{2x} \implies y^{(50)} = ? \qquad \sqrt[3]{8.06} \approx ? \]
\textbf{Lời giải chi tiết:}
\begin{enumerate}[leftmargin=1.5em, itemsep=2pt]
  \item Câu (a): Đặt u = $e^{2x} \Rightarrow u^{(m)} = 2^m e^{2x}$ và v = $x^2 \Rightarrow$  v' = 2x, v'' = 2, $v^{(k)}$ = 0 với k $\ge$  3.
  \item Câu (b): Xét f(x) = $x^{1/3}$ tại a = 8 với $\Delta$ x = dx = 0.06. Tính f(8) = 2 và f'(8) = 1/(3$\cdot 8^{2/3}$) = 1/12.
\end{enumerate}
\[ y^{(50)} = C_{50}^0 (e^{2x})^{(50)} x^2 + C_{50}^1 (e^{2x})^{(49)} (2x) + C_{50}^2 (e^{2x})^{(48)} (2) = 2^{48} e^{2x} (4x^2 + 200x + 2450) \]
\[ \sqrt[3]{8.06} \approx f(8) + f'(8)\cdot 0.06 = 2 + \frac{1}{12}(0.06) = 2 + 0.005 = 2.005 \]
\noindent\textbf{Kết luận:} Giá trị gần đúng là 2.005 (giá trị máy tính chính xác đến 6 chữ số thập phân là 2.004988).
\end{vidu}

\section{Minh họa Hình học bằng TikZ}
\begin{figure}[htbp]
\centering
\begin{tikzpicture}[scale=1.15, >=stealth]
  % Trục tọa độ
  \draw[->, thick] (-0.4,0) -- (4.8,0) node[right] {$x$};
  \draw[->, thick] (0,-0.4) -- (0,3.8) node[above] {$y$};
  
  % Đồ thị hàm số y = f(x)
  \draw[domain=0.4:4.2, smooth, variable=\x, blue!80!black, very thick] 
    plot ({\x}, {0.2*(\x)^2 + 0.5}) node[above left] {$y = f(x)$};
    
  % Tiếp tuyến tại a = 1.5 (y = 0.95 + 0.6(x - 1.5))
  \draw[domain=0.5:4.3, smooth, variable=\x, red!80!black, thick] 
    plot ({\x}, {0.95 + 0.6*(\x - 1.5)}) node[right] {$y = L(x)$};
    
  % Các đoạn hình học tại a = 1.5 và a + dx = 3.5
  \filldraw[black] (1.5, 0.95) circle (2pt) node[above left] {$P(a, f(a))$};
  \draw[dashed] (1.5, 0) node[below] {$a$} -- (1.5, 0.95);
  \draw[dashed] (3.5, 0) node[below] {$a + dx$} -- (3.5, 2.95);
  \draw[dashed] (1.5, 0.95) -- (3.5, 0.95);
  
  % Chú thích dy và Delta y
  \draw[<->, thick, red!80!black] (3.65, 0.95) -- (3.65, 2.15) node[midway, right] {$dy = f'(a)dx$};
  \draw[<->, thick, blue!80!black] (4.35, 0.95) -- (4.35, 2.95) node[midway, right] {$\Delta y$};
  \draw[<->, thick, teal!80!black] (1.5, 0.65) -- (3.5, 0.65) node[midway, below] {$dx = \Delta x$};
\end{tikzpicture}
\caption{Hình 3.1: Ý nghĩa hình học của Xấp xỉ tuyến tính L(x), số gia thực $\Delta$ y của đường cong và vi phân dy trên tiếp tuyến.}
\label{fig:tikz-ch3-linear-approx}
\end{figure}

\section{Ứng dụng Thực tế \& Tối ưu hóa (Học máy (Machine Learning) \& Kỹ thuật Đo lường)}
\begin{ungdung}{Lan truyền Ngược (Backpropagation) và Phân tích Sai số Cảm biến}
Trong Mạng Nơ-ron sâu (Deep Neural Networks), thuật toán Lan truyền ngược (Backpropagation) chính là việc áp dụng liên tiếp Quy tắc đạo hàm hàm hợp (Chain Rule) để tính gradient của hàm mất mát L đối với từng trọng số $w_i$. Ngoài ra, vi phân dy = f'(x)dx được dùng để tính sai số tối đa của thể tích/điện trở khi biết sai số của thước đo.
\[ \frac{\partial \mathcal{L}}{\partial w_1} = \frac{\partial \mathcal{L}}{\partial \hat{y}} \cdot \frac{\partial \hat{y}}{\partial z_2} \cdot \frac{\partial z_2}{\partial a_1} \cdot \frac{\partial a_1}{\partial w_1}, \qquad \frac{\Delta V}{V} \approx \frac{dV}{V} = 3\frac{dR}{R} \]
\noindent\textbf{Ví dụ tính toán thực tế:} Bán kính quả cầu kim loại đo được R = 10 cm với sai số tối đa $\pm$ 0.05 cm (tức dR/R = 0.5\%). Khi đó sai số tương đối tối đa khi tính thể tích V = (4/3)$\pi R^3$ là dV/V = 3(dR/R) = 1.5\%, tương ứng dV $\approx$  4$\pi (10^2$)(0.05) $\approx$  62.83 $cm^3$.
\end{ungdung}

\section{Bài tập Trắc nghiệm Củng cố Chương 3}
\begin{enumerate}[label=\textbf{Câu \arabic*:}, leftmargin=2.2em, itemsep=8pt]
  \item Đạo hàm của hàm hợp f(x) = $\ln$($\cos$(2x)) trên khoảng (-$\pi$/4, $\pi$/4) là:
  \begin{itemize}[label={}, leftmargin=1em, itemsep=2pt]
    \item \textbf{A.} -2 tan(2x) \quad ($-2\tan(2x)$)
    \item \textbf{B.} 2 tan(2x) \quad ($2\tan(2x)$)
    \item \textbf{C.} 1 / cos(2x) \quad ($\frac{1}{\cos(2x)}$)
    \item \textbf{D.} -sin(2x) / cos(2x) \quad ($-\tan(2x)$)
  \end{itemize}
  \vspace{2pt}
  \noindent{\color{teal!85!black}\textbf{Đáp án đúng: A.}} \textit{Giải thích:} Áp dụng quy tắc hàm hợp hai lần: f'(x) = [cos(2x)]' / cos(2x) = -2sin(2x) / cos(2x) = -2tan(2x).
  \[ f'(x) = \frac{(\cos 2x)'}{\cos 2x} = \frac{-2\sin 2x}{\cos 2x} = -2\tan(2x) \]
  \item Cho hàm số y = y(x) xác định ẩn bởi phương trình $e^y$ + xy = e. Giá trị của y'(0) bằng:
  \begin{itemize}[label={}, leftmargin=1em, itemsep=2pt]
    \item \textbf{A.} -1/e \quad ($-\frac{1}{e}$)
    \item \textbf{B.} -e \quad ($-e$)
    \item \textbf{C.} 1/e \quad ($\frac{1}{e}$)
    \item \textbf{D.} 0 \quad ($0$)
  \end{itemize}
  \vspace{2pt}
  \noindent{\color{teal!85!black}\textbf{Đáp án đúng: A.}} \textit{Giải thích:} Tại x = 0, ta có $e^{y(0)}$ + 0 = e $\Rightarrow$  y(0) = 1. Đạo hàm hai vế theo x: $e^y \cdot$  y' + y + x$\cdot$ y' = 0. Thay x = 0, y = 1 vào: $e^1 \cdot$  y'(0) + 1 + 0 = 0 $\Rightarrow$  y'(0) = -1/e.
  \[ e^y y' + y + xy' = 0 \xrightarrow{x=0,\,y=1} e \cdot y'(0) + 1 = 0 \implies y'(0) = -\frac{1}{e} \]
  \item Xấp xỉ tuyến tính L(x) của hàm số f(x) = $\sqrt{1 + 2x}$ tại điểm a = 0 là:
  \begin{itemize}[label={}, leftmargin=1em, itemsep=2pt]
    \item \textbf{A.} L(x) = 1 + x \quad ($L(x) = 1 + x$)
    \item \textbf{B.} L(x) = 1 + 2x \quad ($L(x) = 1 + 2x$)
    \item \textbf{C.} L(x) = 1 + x/2 \quad ($L(x) = 1 + \frac{x}{2}$)
    \item \textbf{D.} L(x) = x \quad ($L(x) = x$)
  \end{itemize}
  \vspace{2pt}
  \noindent{\color{teal!85!black}\textbf{Đáp án đúng: A.}} \textit{Giải thích:} Ta có f(0) = 1 và f'(x) = 2 / (2$\sqrt{1 + 2x}$) = 1 / $\sqrt{1 + 2x}$, suy ra f'(0) = 1. Vậy L(x) = f(0) + f'(0)(x - 0) = 1 + x.
  \[ L(x) = f(0) + f'(0)(x - 0) = 1 + 1\cdot x = 1 + x \]
\end{enumerate}

% ======================================================================
% CHƯƠNG 4: ỨNG DỤNG CỦA ĐẠO HÀM (TRỌNG TÂM) (Buổi 4)
% ======================================================================
\chapter{Ứng dụng của Đạo hàm (Trọng tâm)}
\label{chap:chuong-4}

\begin{tcolorbox}[colback=slate!5!white,colframe=slate!75!black,title={\textbf{Mục lục Chương 4 \& Chuẩn đầu ra (Buổi 4)}}]
\textbf{Tổng quan:} Chương này khai thác sức mạnh của đạo hàm cấp 1 và cấp 2 để giải quyết các bài toán cốt lõi của kỹ thuật và khoa học: Tìm cực trị địa phương và toàn cục, Định lý Giá trị Trung bình (Lagrange), khảo sát tính đơn điệu và tính lồi/lõm (điểm uốn), Quy tắc L'Hospital khử các dạng vô định, Bài toán tối ưu hóa thực tế, Phương pháp lặp Newton giải phương trình phi tuyến và mở đầu về Nguyên hàm.\\[6pt]
\textbf{Nội dung chính trong chương:}
\begin{itemize}[leftmargin=1.5em, itemsep=2pt]
  \item \textbf{Mục 4.1 - 4.2:} Cực trị của hàm số, Định lý Fermat \& Định lý Giá trị Trung bình (MVT)
  \item \textbf{Mục 4.3:} Khảo sát dáng điệu đồ thị: Đơn điệu, Lồi lõm \& Điểm uốn
  \item \textbf{Mục 4.4:} Các dạng vô định và Quy tắc L'Hospital
  \item \textbf{Mục 4.5 - 4.7:} Bài toán Tối ưu hóa, Phương pháp Newton \& Nguyên hàm
\end{itemize}
\textbf{Chuẩn đầu ra (LLO):}
\begin{itemize}[leftmargin=1.5em, itemsep=2pt]
  \item LLO4.1: Ứng dụng đạo hàm tìm cực trị địa phương, cực trị tuyệt đối (toàn cục) trên đoạn đóng.
  \item LLO4.2: Phân tích ảnh hưởng của f' và f'' lên dáng điệu đồ thị (đơn điệu, lồi lõm, điểm uốn).
  \item LLO4.3: Áp dụng thành thạo quy tắc L'Hospital để tìm giới hạn của 7 dạng vô định.
  \item LLO4.4 - 4.6: Giải bài toán tối ưu hóa thực tế, áp dụng thuật toán lặp Newton và tìm nguyên hàm cơ bản.
\end{itemize}
\end{tcolorbox}

\section{[Mục 4.1 - 4.2] Cực trị của hàm số, Định lý Fermat \& Định lý Giá trị Trung bình (MVT)}
Điểm tới hạn (Critical number) của hàm f là điểm c thuộc tập xác định sao cho f'(c) = 0 hoặc f'(c) không tồn tại. Theo Định lý Weierstrass, hàm liên tục trên đoạn đóng [a, b] luôn đạt Giá trị lớn nhất (GTLN) và Giá trị nhỏ nhất (GTNN) tại các điểm tới hạn hoặc tại hai đầu mút a, b.

\begin{congthuc}{Các công thức trọng tâm (4.1 - 4.2)}
\noindent\textbf{Định lý Fermat \& Điểm tới hạn:}
\[ f \text{ đạt cực trị địa phương tại } c \text{ và khả vi tại } c \implies f'(c) = 0 \]
\noindent\textit{Chú giải:} Chiều ngược lại không đúng: f'(0) = 0 với f(x) = $x^3$ nhưng x = 0 không phải cực trị.\\[4pt]
\noindent\textbf{Định lý Giá trị Trung bình Lagrange (Mean Value Theorem):}
\[ f \text{ liên tục trên } [a,b], \text{ khả vi trên } (a,b) \implies \exists c \in (a,b) : f'(c) = \frac{f(b) - f(a)}{b - a} \]
\noindent\textit{Chú giải:} Khi f(a) = f(b), Định lý Lagrange suy biến thành Định lý Rolle: $\exists$ c $\in$  (a, b) sao cho f'(c) = 0.\\[4pt]
\end{congthuc}

\noindent\textbf{Chú giải \& Lưu ý kỹ thuật:}
\begin{itemize}[leftmargin=1.5em, itemsep=2pt]
  \item Quy trình tìm GTLN/GTNN trên đoạn đóng [a, b] (Closed Interval Method): (1) Tìm các điểm tới hạn $c_i \in$  (a, b); (2) Tính $f(c_i$) và f(a), f(b); (3) So sánh các giá trị để kết luận max và min.
\end{itemize}

\begin{vidu}{Ứng dụng Định lý Lagrange chứng minh bất đẳng thức giải tích}
\textbf{Đề bài:} Sử dụng Định lý Giá trị Trung bình Lagrange, chứng minh rằng |$\sin$ a - $\sin$ b| $\le$ |a - b| với mọi số thực a, b.
\[ |\sin a - \sin b| \le |a - b|, \quad \forall a, b \in \mathbb{R} \]
\textbf{Lời giải chi tiết:}
\begin{enumerate}[leftmargin=1.5em, itemsep=2pt]
  \item Bước 1: Nếu a = b thì bất đẳng thức hiển nhiên đúng (0 $\le$  0). Giả sử a < b.
  \item Bước 2: Xét hàm f(x) = sin x liên tục trên [a, b] và khả vi trên (a, b) với f'(x) = cos x.
  \item Bước 3: Áp dụng định lý Lagrange và đánh giá |cos c| $\le$  1.
\end{enumerate}
\[ \exists c \in (a, b) : \frac{\sin b - \sin a}{b - a} = \cos c \implies |\sin a - \sin b| = |\cos c| \cdot |a - b| \le 1 \cdot |a - b| \]
\noindent\textbf{Kết luận:} Bất đẳng thức được chứng minh chặt chẽ (hàm sin x có tính chất Lipschitz với hằng số L = 1).
\end{vidu}

\section{[Mục 4.3] Khảo sát dáng điệu đồ thị: Đơn điệu, Lồi lõm \& Điểm uốn}
Dấu của đạo hàm cấp một f'(x) quyết định chiều biến thiên (đồng biến / nghịch biến), trong khi dấu của đạo hàm cấp hai f''(x) quyết định bề lõm của đồ thị (lõm lên - concave upward khi f'' > 0, lõm xuống - concave downward khi f'' < 0) và xác định điểm uốn.

\begin{congthuc}{Các công thức trọng tâm (4.3)}
\noindent\textbf{Tiêu chuẩn Đạo hàm cấp 2 cho Cực trị địa phương:}
\[ f'(c) = 0 \text{ và } \begin{cases} f''(c) > 0 \implies c \text{ là điểm cực tiểu địa phương} \\ f''(c) < 0 \implies c \text{ là điểm cực đại địa phương} \end{cases} \]
\noindent\textit{Chú giải:} Nếu f''(c) = 0 thì tiêu chuẩn đạo hàm cấp 2 chưa kết luận được, phải dùng bảng xét dấu f'(x).\\[4pt]
\noindent\textbf{Điều kiện Điểm uốn (Inflection Point):}
\[ f''(x) \text{ đổi dấu khi } x \text{ đi qua } x_0 \implies I(x_0, f(x_0)) \text{ là điểm uốn của đồ thị} \]
\noindent\textit{Chú giải:} Tại điểm uốn, đường tiếp tuyến xuyên qua đồ thị hàm số.\\[4pt]
\end{congthuc}

\noindent\textbf{Chú giải \& Lưu ý kỹ thuật:}
\begin{itemize}[leftmargin=1.5em, itemsep=2pt]
  \item Trong giáo trình James Stewart, f''(x) > 0 được gọi là "lõm lên" (Concave upward - hình chữ U), còn f''(x) < 0 được gọi là "lõm xuống" (Concave downward - hình chữ U ngược).
\end{itemize}

\begin{vidu}{Tìm cực trị và điểm uốn của hàm phân thức mũ Gaussian}
\textbf{Đề bài:} Tìm các khoảng lồi/lõm và tọa độ các điểm uốn của đường cong hình chuông f(x) = $e^{-x^2/2}$.
\[ f(x) = e^{-x^2/2}, \quad x \in \mathbb{R} \]
\textbf{Lời giải chi tiết:}
\begin{enumerate}[leftmargin=1.5em, itemsep=2pt]
  \item Bước 1: Tính đạo hàm cấp một f'(x) = -x $e^{-x^2/2}$.
  \item Bước 2: Tính đạo hàm cấp hai f''(x) bằng quy tắc tích và giải phương trình f''(x) = 0.
  \item Bước 3: Xét dấu f''(x) qua các nghiệm x = $\pm$ 1 để kết luận điểm uốn.
\end{enumerate}
\[ f''(x) = -e^{-x^2/2} + (-x)(-x)e^{-x^2/2} = (x^2 - 1)e^{-x^2/2} \]
\[ f''(x) = 0 \iff x^2 - 1 = 0 \iff x = \pm 1 \quad (\text{do } e^{-x^2/2} > 0, \forall x) \]
\noindent\textbf{Kết luận:} Vì f''(x) đổi dấu qua x = -1 và x = 1 nên đồ thị có 2 điểm uốn là (-1, $e^{-1/2}$) và (1, $e^{-1/2}$).
\end{vidu}

\section{[Mục 4.4] Các dạng vô định và Quy tắc L'Hospital}
Quy tắc L'Hospital cho phép tính giới hạn của thương hai hàm số có dạng vô định 0/0 hoặc $\infty / \infty$ bằng cách thay thế bằng giới hạn của thương hai đạo hàm. Các dạng vô định khác (0$\cdot \infty$, $\infty - \infty$, $0^0$, $\infty^0$, $1^\infty$) đều được quy về hai dạng chuẩn này.

\begin{congthuc}{Các công thức trọng tâm (4.4)}
\noindent\textbf{Định lý L'Hospital (Dạng 0/0 hoặc $\infty / \infty$):}
\[ \lim_{x \to a} \frac{f(x)}{g(x)} = \lim_{x \to a} \frac{f'(x)}{g'(x)} \quad \left(\text{nếu giới hạn vế phải tồn tại hoặc bằng } \pm\infty\right) \]
\noindent\textit{Chú giải:} Luôn kiểm tra chắc chắn biểu thức có dạng 0/0 hoặc $\infty / \infty$ trước mỗi lần áp dụng L'Hospital!\\[4pt]
\noindent\textbf{Cách khử dạng vô định lũy thừa $0^0$, $\infty^0$, $1^\infty$:}
\[ \lim_{x \to a} [u(x)]^{v(x)} = \exp\left( \lim_{x \to a} v(x) \ln u(x) \right) = e^L \]
\noindent\textit{Chú giải:} Biến đổi tích v(x)$\cdot$ ln u(x) (dạng 0$\cdot \infty$) thành thương ln u(x) / [1/v(x)] để dùng L'Hospital.\\[4pt]
\end{congthuc}

\noindent\textbf{Chú giải \& Lưu ý kỹ thuật:}
\begin{itemize}[leftmargin=1.5em, itemsep=2pt]
  \item Mẹo tính nhanh: Kết hợp thay thế vô cùng bé tương đương (như sin x ~ x, ln(1+x) ~ x) cho các nhân tử trước khi áp dụng L'Hospital sẽ giảm đáng kể độ phức tạp khi đạo hàm.
\end{itemize}

\begin{vidu}{Khử dạng vô định 0/0 và $1^\infty$ bằng quy tắc L'Hospital}
\textbf{Đề bài:} Tính giới hạn L = \$\lim_{x \to 0} \frac{e^x - e^{-x} - 2x}{x - \sin x}$.
\[ L = \lim_{x \to 0} \frac{e^x - e^{-x} - 2x}{x - \sin x} \quad \left(\text{Dạng } \frac{0}{0}\right) \]
\textbf{Lời giải chi tiết:}
\begin{enumerate}[leftmargin=1.5em, itemsep=2pt]
  \item Bước 1: Kiểm tra khi x $\to$  0 cả tử và mẫu đều tiến về 0 (dạng 0/0). Áp dụng L'Hospital lần 1.
  \item Bước 2: Tiếp tục kiểm tra dạng 0/0 và áp dụng L'Hospital lần 2 và lần 3.
\end{enumerate}
\[ L \stackrel{\text{L'H}}{=} \lim_{x \to 0} \frac{e^x + e^{-x} - 2}{1 - \cos x} \stackrel{\text{L'H}}{=} \lim_{x \to 0} \frac{e^x - e^{-x}}{\sin x} \stackrel{\text{L'H}}{=} \lim_{x \to 0} \frac{e^x + e^{-x}}{\cos x} = \frac{1 + 1}{1} = 2 \]
\noindent\textbf{Kết luận:} Giới hạn cần tìm bằng L = 2.
\end{vidu}

\section{[Mục 4.5 - 4.7] Bài toán Tối ưu hóa, Phương pháp Newton \& Nguyên hàm}
Bài toán tối ưu hóa chuyển mô tả thực tế thành bài toán tìm cực trị toàn cục của hàm mục tiêu. Phương pháp lặp Newton–Raphson sử dụng tiếp tuyến để tìm nghiệm gần đúng của phương trình f(x) = 0 với tốc độ hội tụ bậc hai. Phần cuối giới thiệu Nguyên hàm F(x) thỏa mãn F'(x) = f(x).

\begin{congthuc}{Các công thức trọng tâm (4.5 - 4.7)}
\noindent\textbf{Công thức lặp Newton - Raphson giải phương trình f(x) = 0:}
\[ x_{n+1} = x_n - \frac{f(x_n)}{f'(x_n)}, \quad n = 0, 1, 2, \dots \quad (f'(x_n) \neq 0) \]
\noindent\textit{Chú giải:} Ý nghĩa hình học: $x_{n+1}$ là giao điểm của tiếp tuyến tại $(x_n$, $f(x_n$)) với trục hoành Ox.\\[4pt]
\noindent\textbf{Định nghĩa Họ Nguyên hàm:}
\[ F'(x) = f(x), \forall x \in I \implies \int f(x)\,dx = F(x) + C \]
\noindent\textit{Chú giải:} Mọi nguyên hàm của f(x) trên một khoảng I chỉ sai khác nhau một hằng số cộng C.\\[4pt]
\end{congthuc}

\noindent\textbf{Chú giải \& Lưu ý kỹ thuật:}
\begin{itemize}[leftmargin=1.5em, itemsep=2pt]
  \item Trong bài toán thiết kế hình trụ có thể tích V cho trước sao cho diện tích toàn phần S = 2$\pi r^2$ + 2$\pi$ rh nhỏ nhất: Thay h = V/($\pi r^2$) vào S(r) rồi giải S'(r) = 0 ta luôn thu được tỉ lệ vàng kỹ thuật: chiều cao bằng đường kính (h = 2r).
\end{itemize}

\begin{vidu}{Bài toán tối ưu chi phí vật liệu chế tạo bình chứa hình trụ}
\textbf{Đề bài:} Cần chế tạo một lon nhôm hình trụ kín hai đầu có thể tích cố định V = 16$\pi (cm^3$). Hãy tìm bán kính đáy r và chiều cao h để tổng diện tích bề mặt vỏ nhôm là nhỏ nhất.
\[ V = \pi r^2 h = 16\pi, \qquad S(r) = 2\pi r^2 + 2\pi r h \to \min \]
\textbf{Lời giải chi tiết:}
\begin{enumerate}[leftmargin=1.5em, itemsep=2pt]
  \item Bước 1: Từ ràng buộc thể tích $\pi r^2$ h = 16$\pi$ , rút chiều cao h = 16 / $r^2$ (với r > 0).
  \item Bước 2: Thế h vào hàm diện tích toàn phần S(r) = 2$\pi r^2$ + 32$\pi$ /r.
  \item Bước 3: Tính đạo hàm S'(r), giải phương trình S'(r) = 0 và kiểm tra S''(r) > 0.
\end{enumerate}
\[ S'(r) = 4\pi r - \frac{32\pi}{r^2} = 0 \iff r^3 = 8 \iff r = 2 \text{ (cm)} \]
\[ S''(r) = 4\pi + \frac{64\pi}{r^3} > 0, \forall r > 0 \implies r = 2 \text{ đạt GTNN toàn cục, khi đó } h = \frac{16}{2^2} = 4 \text{ (cm)} \]
\noindent\textbf{Kết luận:} Diện tích vật liệu nhỏ nhất là $S_min$ = 24$\pi cm^2$ khi bán kính đáy r = 2 cm và chiều cao h = 2r = 4 cm.
\end{vidu}

\section{Minh họa Hình học bằng TikZ}
\begin{figure}[htbp]
\centering
\begin{tikzpicture}[scale=1.1, >=stealth]
  % Trục tọa độ
  \draw[->, thick] (-0.4,0) -- (5.0,0) node[right] {$x$};
  \draw[->, thick] (0,-0.4) -- (0,3.8) node[above] {$y$};
  
  % Đồ thị hàm số y = f(x)
  \draw[domain=0.6:4.2, smooth, variable=\x, blue!80!black, very thick] 
    plot ({\x}, {0.9 + 1.8*sin((\x - 0.6)*45) + 0.25*(\x)});
    
  % Dây cung nối A và B
  \filldraw[black] (0.6, 1.05) circle (2pt) node[left] {$A(a,f(a))$};
  \filldraw[black] (4.2, 2.75) circle (2pt) node[right] {$B(b,f(b))$};
  \draw[dashed, thick, teal!80!black] (0.6, 1.05) -- (4.2, 2.75);
  
  % Tiếp tuyến song song tại điểm c = 2.35
  \filldraw[red!80!black] (2.35, 3.25) circle (2.2pt) node[above] {$P(c,f(c))$};
  \draw[red!80!black, thick] (1.1, 2.66) -- (3.6, 3.84);
  
  % Đường gióng xuống trục hoành
  \draw[dotted] (0.6,0) node[below] {$a$} -- (0.6,1.05);
  \draw[dotted] (2.35,0) node[below, red!80!black] {$c$} -- (2.35,3.25);
  \draw[dotted] (4.2,0) node[below] {$b$} -- (4.2,2.75);
\end{tikzpicture}
\caption{Hình 4.1: Định lý Giá trị Trung bình Lagrange (MVT) — Tiếp tuyến tại c song song với dây cung AB.}
\label{fig:tikz-ch4-mvt}
\end{figure}

\section{Ứng dụng Thực tế \& Tối ưu hóa (Đồ họa Máy tính (Game Engine) \& Tối ưu Công nghiệp)}
\begin{ungdung}{Thuật toán Fast Inverse Square Root (Phương pháp Newton) \& Thiết kế Bao bì Tối ưu}
Trong các engine đồ họa 3D (như Quake III và GPU hiện đại), việc chuẩn hóa vector pháp tuyến ánh sáng cần tính 1/$\sqrt{x}$ hàng triệu lần mỗi giây. Thuật toán kinh điển sử dụng 1 bước lặp Newton–Raphson trên hàm f(y) = 1/$y^2$ - x = 0 giúp tăng tốc độ tính toán gấp 4 lần so với phép chia phần cứng.
\[ f(y) = y^{-2} - x = 0 \implies y_{n+1} = y_n - \frac{y_n^{-2} - x}{-2y_n^{-3}} = \frac{1}{2}y_n(3 - x y_n^2) \]
\noindent\textbf{Ví dụ tính toán thực tế:} Để tính 1/$\sqrt{2} \approx$  0.707107 với giá trị khởi tạo $y_0$ = 0.7: Chỉ sau 1 bước lặp Newton ta có $y_1$ = 0.5 $\times$  0.7 $\times$  (3 - 2 $\times$  0.49) = 0.7070 (sai số chỉ 0.015\%).
\end{ungdung}

\section{Bài tập Trắc nghiệm Củng cố Chương 4}
\begin{enumerate}[label=\textbf{Câu \arabic*:}, leftmargin=2.2em, itemsep=8pt]
  \item Giá trị lớn nhất và giá trị nhỏ nhất của hàm số f(x) = $x^3 - 3x^2$ + 1 trên đoạn [-1/2, 4] lần lượt là:
  \begin{itemize}[label={}, leftmargin=1em, itemsep=2pt]
    \item \textbf{A.} GTLN = 17, GTNN = -3 \quad ($\max = 17, \; \min = -3$)
    \item \textbf{B.} GTLN = 1, GTNN = -3 \quad ($\max = 1, \; \min = -3$)
    \item \textbf{C.} GTLN = 17, GTNN = 1/8 \quad ($\max = 17, \; \min = \frac{1}{8}$)
    \item \textbf{D.} GTLN = 1, GTNN = 1/8 \quad ($\max = 1, \; \min = \frac{1}{8}$)
  \end{itemize}
  \vspace{2pt}
  \noindent{\color{teal!85!black}\textbf{Đáp án đúng: A.}} \textit{Giải thích:} Ta có f'(x) = $3x^2$ - 6x = 0 $\iff$  x = 0 hoặc x = 2 (đều thuộc [-1/2, 4]). Tính giá trị: f(-1/2) = 1/8, f(0) = 1, f(2) = -3, f(4) = 17. Vậy max = 17 (tại x = 4) và min = -3 (tại x = 2).
  \[ f(-1/2) = 0.125, \; f(0) = 1, \; f(2) = -3, \; f(4) = 17 \implies \max_{[-1/2, 4]} f = 17, \; \min_{[-1/2, 4]} f = -3 \]
  \item Bằng quy tắc L'Hospital, giới hạn \$\lim_{x \to 0^+} x^{\sin x}$ có giá trị bằng:
  \begin{itemize}[label={}, leftmargin=1em, itemsep=2pt]
    \item \textbf{A.} 1 \quad ($1$)
    \item \textbf{B.} 0 \quad ($0$)
    \item \textbf{C.} e \quad ($e$)
    \item \textbf{D.} +$\infty$ \quad ($+\infty$)
  \end{itemize}
  \vspace{2pt}
  \noindent{\color{teal!85!black}\textbf{Đáp án đúng: A.}} \textit{Giải thích:} Xét ln L = $\lim_{x\to 0^+}$ (sin x $\cdot$  ln x) = $\lim_{x\to 0^+} [ln x / (1/sin x)]$. Vì sin x ~ x nên xét $\lim_{x\to 0^+} [ln x / (1/x)] = \lim_{x\to 0^+} [(1/x) / (-1/x^2)]$ = lim (-x) = 0. Do đó L = $e^0$ = 1.
  \[ \ln L = \lim_{x \to 0^+} \frac{\ln x}{1/x} \stackrel{\text{L'H}}{=} \lim_{x \to 0^+} \frac{1/x}{-1/x^2} = 0 \implies L = e^0 = 1 \]
  \item Áp dụng phương pháp Newton tìm nghiệm gần đúng của phương trình $x^2$ - 2 = 0 với $x_0$ = 1, sau 1 bước lặp ta được $x_1$ bằng:
  \begin{itemize}[label={}, leftmargin=1em, itemsep=2pt]
    \item \textbf{A.} $x_1$ = 1.5 \quad ($x_1 = 1.5$)
    \item \textbf{B.} $x_1$ = 1.4 \quad ($x_1 = 1.4$)
    \item \textbf{C.} $x_1$ = 1.414 \quad ($x_1 = 1.414$)
    \item \textbf{D.} $x_1$ = 2.0 \quad ($x_1 = 2.0$)
  \end{itemize}
  \vspace{2pt}
  \noindent{\color{teal!85!black}\textbf{Đáp án đúng: A.}} \textit{Giải thích:} Với f(x) = $x^2$ - 2 $\Rightarrow$  f'(x) = 2x. Công thức Newton: $x_1 = x_0 - f(x_0$)/$f'(x_0$) = 1 - $(1^2$ - 2)/(2$\cdot$ 1) = 1 - (-1)/2 = 1.5.
  \[ x_1 = x_0 - \frac{x_0^2 - 2}{2x_0} = 1 - \frac{1 - 2}{2} = 1.5 \]
\end{enumerate}

% ======================================================================
% CHƯƠNG 5: TÍCH PHÂN (TRỌNG TÂM) (Buổi 5)
% ======================================================================
\chapter{Tích phân (Trọng tâm)}
\label{chap:chuong-5}

\begin{tcolorbox}[colback=slate!5!white,colframe=slate!75!black,title={\textbf{Mục lục Chương 5 \& Chuẩn đầu ra (Buổi 5)}}]
\textbf{Tổng quan:} Chương này mở đầu phép tính tích phân từ bài toán tính diện tích miền dưới đường cong và bài toán tính quãng đường từ vận tốc thông qua tổng Riemann. Điểm nhấn quan trọng nhất là Định lý Cơ bản của Giải tích (FTC) kết nối phép tính đạo hàm và phép tính tích phân, cùng quy tắc đổi biến số.\\[6pt]
\textbf{Nội dung chính trong chương:}
\begin{itemize}[leftmargin=1.5em, itemsep=2pt]
  \item \textbf{Mục 5.1 - 5.2:} Bài toán Diện tích, Khoảng cách, Tổng Riemann \& Tích phân xác định
  \item \textbf{Mục 5.3:} Định lý Cơ bản của Giải tích (FTC Phần 1 \& Công thức Newton - Leibniz)
  \item \textbf{Mục 5.4 - 5.5:} Tích phân bất định \& Quy tắc đổi biến số (Substitution Rule)
\end{itemize}
\textbf{Chuẩn đầu ra (LLO):}
\begin{itemize}[leftmargin=1.5em, itemsep=2pt]
  \item LLO5.1: Hiểu và vận dụng được khái niệm tích phân xác định thông qua tổng Riemann, bài toán diện tích và khoảng cách.
  \item LLO5.2: Nắm vững Định lý Cơ bản của Giải tích (phần 1 và phần 2) và giải thành thạo tích phân xác định, tích phân bất định bằng phương pháp đổi biến.
\end{itemize}
\end{tcolorbox}

\section{[Mục 5.1 - 5.2] Bài toán Diện tích, Khoảng cách, Tổng Riemann \& Tích phân xác định}
Chia đoạn [a, b] thành n đoạn con đều nhau có bề rộng $\Delta$ x = (b - a)/n và chọn các điểm mẫu $x_i$\^{}*. Tổng diện tích các hình chữ nhật xấp xỉ gọi là tổng Riemann. Khi cho n $\to \infty$, giới hạn của tổng Riemann chính là tích phân xác định của f(x) trên [a, b].

\begin{congthuc}{Các công thức trọng tâm (5.1 - 5.2)}
\noindent\textbf{Định nghĩa Tích phân xác định qua Tổng Riemann:}
\[ \int_a^b f(x)\,dx = \lim_{n \to \infty} \sum_{i=1}^n f(x_i^*) \Delta x, \qquad \Delta x = \frac{b - a}{n}, \; x_i = a + i\Delta x \]
\noindent\textit{Chú giải:} Nếu f(x) nhận cả giá trị âm và dương, tích phân xác định là "diện tích đại số" (phần trên Ox trừ phần dưới Ox).\\[4pt]
\noindent\textbf{Tính chất đối xứng của tích phân trên đoạn đối xứng [-a, a]:}
\[ \int_{-a}^a f(x)\,dx = \begin{cases} 2\int_0^a f(x)\,dx & \text{nếu } f(-x) = f(x) \text{ (Hàm chẵn)} \\ 0 & \text{nếu } f(-x) = -f(x) \text{ (Hàm lẻ)} \end{cases} \]
\noindent\textit{Chú giải:} Tính chất hàm lẻ giúp tính nhẩm ngay nhiều tích phân phức tạp trên đoạn [-a, a] bằng 0.\\[4pt]
\end{congthuc}

\noindent\textbf{Chú giải \& Lưu ý kỹ thuật:}
\begin{itemize}[leftmargin=1.5em, itemsep=2pt]
  \item Trong vật lý, nếu v(t) là hàm vận tốc thì tích phân $\int_a^b$ v(t)dt là độ dịch chuyển (displacement), còn $\int_a^b$ |v(t)|dt là tổng quãng đường thực đi được (total distance).
\end{itemize}

\begin{vidu}{Chuyển giới hạn của tổng về tích phân xác định (Tổng Riemann)}
\textbf{Đề bài:} Tính giới hạn của dãy tổng S = \$\lim_{n \to \infty} \sum_{i=1}^n \frac{n}{n^2 + i^2}$ bằng cách sử dụng định nghĩa tích phân xác định.
\[ S = \lim_{n \to \infty} \sum_{i=1}^n \frac{n}{n^2 + i^2} \]
\textbf{Lời giải chi tiết:}
\begin{enumerate}[leftmargin=1.5em, itemsep=2pt]
  \item Bước 1: Chia cả tử và mẫu của phân thức cho $n^2$ để làm xuất hiện $\Delta$ x = 1/n và điểm mẫu $x_i$ = i/n.
  \item Bước 2: Nhận diện tổng Riemann của hàm f(x) = 1/(1 + $x^2$) trên đoạn [0, 1].
  \item Bước 3: Tính tích phân xác định tương ứng.
\end{enumerate}
\[ S = \lim_{n \to \infty} \sum_{i=1}^n \frac{1}{1 + (i/n)^2} \cdot \frac{1}{n} = \int_0^1 \frac{1}{1 + x^2}\,dx = \Big[\arctan x\Big]_0^1 = \frac{\pi}{4} - 0 = \frac{\pi}{4} \]
\noindent\textbf{Kết luận:} Giới hạn của tổng đã cho bằng $\pi$ /4.
\end{vidu}

\section{[Mục 5.3] Định lý Cơ bản của Giải tích (FTC Phần 1 \& Công thức Newton - Leibniz)}
Định lý Cơ bản của Giải tích (Fundamental Theorem of Calculus - FTC) gồm 2 phần: Phần 1 khẳng định đạo hàm của hàm tích phân cận biến thiên bằng chính hàm dưới dấu tích phân; Phần 2 (Newton–Leibniz) cho phép tính tích phân xác định thông qua hiệu giá trị của một nguyên hàm tại hai cận.

\begin{congthuc}{Các công thức trọng tâm (5.3)}
\noindent\textbf{FTC Phần 1: Đạo hàm của hàm tích phân có cận phụ thuộc x (Quy tắc Leibniz):}
\[ \frac{d}{dx} \left[ \int_a^x f(t)\,dt \right] = f(x), \qquad \frac{d}{dx} \left[ \int_{u(x)}^{v(x)} f(t)\,dt \right] = f(v(x))\cdot v'(x) - f(u(x))\cdot u'(x) \]
\noindent\textit{Chú giải:} Công thức tổng quát kết hợp FTC Phần 1 với Quy tắc đạo hàm hàm hợp (Chain Rule).\\[4pt]
\noindent\textbf{FTC Phần 2: Công thức Newton - Leibniz:}
\[ \int_a^b f(x)\,dx = F(b) - F(a) = \Big[F(x)\Big]_a^b \quad \text{với } F'(x) = f(x) \]
\noindent\textit{Chú giải:} Định lý tốc độ biến thiên tổng cộng (Net Change Theorem): $\int_a^b$ F'(x)dx = F(b) - F(a).\\[4pt]
\end{congthuc}

\noindent\textbf{Chú giải \& Lưu ý kỹ thuật:}
\begin{itemize}[leftmargin=1.5em, itemsep=2pt]
  \item Các bài toán tìm giới hạn dạng 0/0 có chứa tích phân xác định ở tử số (ví dụ $\lim_{x\to 0} [\int_0^{x^2} sin(t)dt] / x^4$) luôn giải cực nhanh bằng cách kết hợp Quy tắc L'Hospital và FTC Phần 1.
\end{itemize}

\begin{vidu}{Đạo hàm của hàm tích phân cận biến thiên kết hợp L'Hospital}
\textbf{Đề bài:} Cho hàm số G(x) = $\int_{\sin x}^{x^2} \sqrt{1 + t^3}$\,dt. Hãy tính đạo hàm G'(x).
\[ G(x) = \int_{\sin x}^{x^2} \sqrt{1 + t^3}\,dt \implies G'(x) = ? \]
\textbf{Lời giải chi tiết:}
\begin{enumerate}[leftmargin=1.5em, itemsep=2pt]
  \item Bước 1: Xác định hàm dưới dấu tích phân f(t) = $\sqrt{1 + t^3}$, cận dưới u(x) = sin x, cận trên v(x) = $x^2$.
  \item Bước 2: Tính đạo hàm hai cận: u'(x) = cos x và v'(x) = 2x.
  \item Bước 3: Áp dụng công thức G'(x) = f(v(x))$\cdot$ v'(x) - f(u(x))$\cdot$ u'(x).
\end{enumerate}
\[ G'(x) = \sqrt{1 + (x^2)^3} \cdot (x^2)' - \sqrt{1 + (\sin x)^3} \cdot (\sin x)' = 2x\sqrt{1 + x^6} - \cos x \sqrt{1 + \sin^3 x} \]
\noindent\textbf{Kết luận:} Đạo hàm của hàm tích phân là G'(x) = 2x$\sqrt{1 + x^6}$ - cos(x)$\sqrt{1 + sin^3 x}$.
\end{vidu}

\section{[Mục 5.4 - 5.5] Tích phân bất định \& Quy tắc đổi biến số (Substitution Rule)}
Quy tắc đổi biến số trong tích phân là phép toán ngược của quy tắc đạo hàm hàm hợp (Chain Rule). Khi gặp tích phân có chứa một biểu thức u = g(x) đi kèm với vi phân du = g'(x)dx của nó, ta đổi sang biến mới u để đưa về tích phân cơ bản.

\begin{congthuc}{Các công thức trọng tâm (5.4 - 5.5)}
\noindent\textbf{Công thức Đổi biến trong Tích phân bất định và xác định:}
\[ \int f(g(x))g'(x)\,dx = \int f(u)\,du, \qquad \int_a^b f(g(x))g'(x)\,dx = \int_{g(a)}^{g(b)} f(u)\,du \]
\noindent\textit{Chú giải:} Khi tính tích phân xác định bằng đổi biến, nhớ ĐỔI CẬN từ x $\in  [a, b]$ sang u $\in  [g(a), g(b)]$ để không cần trả lại biến cũ.\\[4pt]
\end{congthuc}

\noindent\textbf{Chú giải \& Lưu ý kỹ thuật:}
\begin{itemize}[leftmargin=1.5em, itemsep=2pt]
  \item Kinh nghiệm chọn u = g(x): Thường đặt u bằng biểu thức nằm trong căn thức, biểu thức ở mẫu số, biểu thức trong hàm mũ $e^{g(x)}$, hoặc đặt u = ln x khi xuất hiện dx/x.
\end{itemize}

\begin{vidu}{Tính tích phân xác định bằng phương pháp đổi biến và tính chất đối xứng}
\textbf{Đề bài:} Tính tích phân I = $\int_0^{\pi/2} \frac{\sin(2x)}{1 + \cos^2 x}$\,dx.
\[ I = \int_0^{\pi/2} \frac{\sin(2x)}{1 + \cos^2 x}\,dx \]
\textbf{Lời giải chi tiết:}
\begin{enumerate}[leftmargin=1.5em, itemsep=2pt]
  \item Bước 1: Đặt u = 1 + $cos^2$ x, suy ra du = 2 cos x (-sin x) dx = -sin(2x) dx.
  \item Bước 2: Đổi cận: Khi x = 0 thì u = 1 + $1^2$ = 2; khi x = $\pi$ /2 thì u = 1 + 0 = 1.
  \item Bước 3: Thay vào tích phân và tính theo biến u.
\end{enumerate}
\[ I = \int_2^1 \frac{-du}{u} = \int_1^2 \frac{du}{u} = \Big[\ln|u|\Big]_1^2 = \ln 2 - \ln 1 = \ln 2 \]
\noindent\textbf{Kết luận:} Giá trị chính xác của tích phân là I = ln 2 $\approx$  0.6931.
\end{vidu}

\section{Minh họa Hình học bằng TikZ}
\begin{figure}[htbp]
\centering
\begin{tikzpicture}[scale=1.1, >=stealth]
  % Các hình chữ nhật xấp xỉ Riemann
  \foreach \x/\h in {0.6/1.28, 1.3/1.82, 2.0/2.45, 2.7/2.95, 3.4/3.12} {
    \fill[cyan!20, draw=cyan!70!black, thick] (\x, 0) rectangle (\x+0.7, \h);
  }
  % Trục tọa độ
  \draw[->, thick] (-0.3,0) -- (4.8,0) node[right] {$x$};
  \draw[->, thick] (0,-0.3) -- (0,3.8) node[above] {$y$};
  
  % Đồ thị hàm số y = f(x)
  \draw[domain=0.4:4.3, smooth, variable=\x, blue!85!black, very thick] 
    plot ({\x}, {0.8 + 1.6*sin((\x - 0.3)*38) + 0.22*(\x)}) node[right] {$y = f(x)$};
    
  % Nhãn cận a, b và bề rộng Delta x
  \node[below] at (0.6,0) {$a$};
  \node[below] at (4.1,0) {$b$};
  \draw[<->, red!80!black, thick] (2.0, -0.35) -- (2.7, -0.35) node[midway, below] {$\Delta x$};
\end{tikzpicture}
\caption{Hình 5.1: Xấp xỉ diện tích miền phẳng dưới đường cong y = f(x) bằng Tổng Riemann với các hình chữ nhật.}
\label{fig:tikz-ch5-riemann}
\end{figure}

\section{Ứng dụng Thực tế \& Tối ưu hóa (Kỹ thuật Năng lượng \& Điện tử Công suất)}
\begin{ungdung}{Tích phân Công suất Tức thời thành Điện năng Tiêu thụ (Smart Meter)}
Công tơ điện tử thông minh (Smart Meter) lấy mẫu điện áp u(t) và dòng điện i(t) ở tần số hàng nghìn Hz để tính công suất tức thời p(t) = u(t)$\cdot$ i(t). Tổng điện năng tiêu thụ E (tính bằng kWh) chính là tích phân xác định của p(t) theo thời gian, được vi xử lý tính bằng tổng Riemann hình thang.
\[ E = \int_{t_1}^{t_2} p(t)\,dt = \int_{t_1}^{t_2} u(t)i(t)\,dt \approx \sum_{k=1}^N p(t_k)\Delta t \]
\noindent\textbf{Ví dụ tính toán thực tế:} Nếu công suất phụ tải biến thiên theo ngày là P(t) = 3 + 2sin($\pi$ t/12) (kW) từ t = 0h đến t = 12h, thì điện năng tiêu thụ trong 12 giờ là E = $\int_0^{12} [3 + 2sin(\pi t/12)]$dt = [3t - (24/$\pi$ )cos($\pi$ t/$12)]_0^{12}$ = 36 + 48/$\pi \approx$  51.28 kWh.
\end{ungdung}

\section{Bài tập Trắc nghiệm Củng cố Chương 5}
\begin{enumerate}[label=\textbf{Câu \arabic*:}, leftmargin=2.2em, itemsep=8pt]
  \item Giá trị của tích phân xác định I = $\int_{-2}^{2} \left(x^3 \cos x + \sqrt{4 - x^2}\right)$ dx bằng:
  \begin{itemize}[label={}, leftmargin=1em, itemsep=2pt]
    \item \textbf{A.} 2$\pi$  \quad ($2\pi$)
    \item \textbf{B.} 4$\pi$  \quad ($4\pi$)
    \item \textbf{C.} 0 \quad ($0$)
    \item \textbf{D.} $\pi$  \quad ($\pi$)
  \end{itemize}
  \vspace{2pt}
  \noindent{\color{teal!85!black}\textbf{Đáp án đúng: A.}} \textit{Giải thích:} Tách thành 2 tích phân trên đoạn đối xứng [-2, 2]: Hàm g(x) = $x^3$ cos x là hàm lẻ nên tích phân bằng 0. Hàm h(x) = $\sqrt{4 - x^2}$ biểu diễn nửa trên hình tròn tâm O bán kính R = 2 nên diện tích bằng (1/2)$\pi R^2$ = 2$\pi$ .
  \[ I = \underbrace{\int_{-2}^2 x^3\cos x\,dx}_{=0 \text{ (hàm lẻ)}} + \underbrace{\int_{-2}^2 \sqrt{4-x^2}\,dx}_{=\frac{1}{2}\pi(2^2) = 2\pi} = 2\pi \]
  \item Cho hàm số F(x) = $\int_1^{x^3} \ln$(t)\,dt với x > 0. Giá trị của F'(e) bằng:
  \begin{itemize}[label={}, leftmargin=1em, itemsep=2pt]
    \item \textbf{A.} $9e^2$ \quad ($9e^2$)
    \item \textbf{B.} $3e^2$ \quad ($3e^2$)
    \item \textbf{C.} 3 \quad ($3$)
    \item \textbf{D.} 1/$e^3$ \quad ($e^{-3}$)
  \end{itemize}
  \vspace{2pt}
  \noindent{\color{teal!85!black}\textbf{Đáp án đúng: A.}} \textit{Giải thích:} Theo Định lý Cơ bản của Giải tích Phần 1 kết hợp hàm hợp: F'(x) = $ln(x^3$) $\cdot (x^3$)' = (3 ln x) $\cdot (3x^2$) = $9x^2$ ln x. Thay x = e ta được F'(e) = $9e^2 \cdot$  ln(e) = $9e^2$.
  \[ F'(x) = \ln(x^3)\cdot 3x^2 = 9x^2\ln x \implies F'(e) = 9e^2 \ln e = 9e^2 \]
  \item Tính tích phân I = $\int_0^1$ x $e^{-x^2}$\,dx, ta thu được kết quả:
  \begin{itemize}[label={}, leftmargin=1em, itemsep=2pt]
    \item \textbf{A.} (1 - 1/e) / 2 \quad ($\frac{1 - e^{-1}}{2}$)
    \item \textbf{B.} 1 - 1/e \quad ($1 - e^{-1}$)
    \item \textbf{C.} 1 / (2e) \quad ($\frac{1}{2e}$)
    \item \textbf{D.} (e - 1) / 2 \quad ($\frac{e - 1}{2}$)
  \end{itemize}
  \vspace{2pt}
  \noindent{\color{teal!85!black}\textbf{Đáp án đúng: A.}} \textit{Giải thích:} Đặt u = $-x^2 \Rightarrow$  du = -2x dx $\Rightarrow$  x dx = -du/2. Đổi cận x: 0 $\to$  1 thì u: 0 $\to$  -1. Khi đó I = -(1/2) $\int_0^{-1} e^u$ du = (1/2)(1 - $e^{-1}$).
  \[ I = -\frac{1}{2}\Big[e^{-x^2}\Big]_0^1 = \frac{1 - e^{-1}}{2} \]
\end{enumerate}

% ======================================================================
% CHƯƠNG 6: ỨNG DỤNG CỦA TÍCH PHÂN (TRỌNG TÂM) (Buổi 6)
% ======================================================================
\chapter{Ứng dụng của Tích phân (Trọng tâm)}
\label{chap:chuong-6}

\begin{tcolorbox}[colback=slate!5!white,colframe=slate!75!black,title={\textbf{Mục lục Chương 6 \& Chuẩn đầu ra (Buổi 6)}}]
\textbf{Tổng quan:} Ứng dụng trực tiếp nguyên lý "Chia nhỏ và Cộng dồn" (Slice and Sum) của tích phân xác định để tính diện tích miền phẳng giới hạn bởi các đường cong, thể tích vật thể tròn xoay bằng phương pháp lát cắt (đĩa/vòng đệm) và phương pháp vỏ hình trụ, tính công cơ học của lực biến thiên và giá trị trung bình của hàm số.\\[6pt]
\textbf{Nội dung chính trong chương:}
\begin{itemize}[leftmargin=1.5em, itemsep=2pt]
  \item \textbf{Mục 6.1:} Diện tích miền phẳng giới hạn bởi các đường cong
  \item \textbf{Mục 6.2 - 6.3:} Thể tích vật thể tròn xoay: Phương pháp Đĩa/Vòng đệm \& Vỏ hình trụ
  \item \textbf{Mục 6.4 - 6.5:} Công của lực biến thiên \& Giá trị trung bình của hàm số
\end{itemize}
\textbf{Chuẩn đầu ra (LLO):}
\begin{itemize}[leftmargin=1.5em, itemsep=2pt]
  \item LLO6.1: Áp dụng thành thạo tích phân để tính diện tích miền giữa 2 đường cong, thể tích khối tròn xoay, thể tích vỏ hình trụ, công cơ học và giá trị trung bình của hàm số.
\end{itemize}
\end{tcolorbox}

\section{[Mục 6.1] Diện tích miền phẳng giới hạn bởi các đường cong}
Diện tích miền phẳng giới hạn bởi hai đồ thị y = f(x), y = g(x) và hai đường thẳng x = a, x = b được tính bằng tích phân của trị tuyệt đối hiệu hai hàm. Nếu miền giới hạn thuận lợi hơn khi chiếu lên trục tung Oy (x = u(y) và x = v(y)), ta lấy tích phân theo biến y.

\begin{congthuc}{Các công thức trọng tâm (6.1)}
\noindent\textbf{Diện tích miền phẳng chiếu theo trục Ox hoặc trục Oy:}
\[ S = \int_a^b |f(x) - g(x)|\,dx \quad (\text{dải dọc } dx), \qquad S = \int_c^d |u(y) - v(y)|\,dy \quad (\text{dải ngang } dy) \]
\noindent\textit{Chú giải:} Nếu đề bài chưa cho cận a, b thì cần giải phương trình hoành độ giao điểm f(x) = g(x) để tìm cận.\\[4pt]
\end{congthuc}

\noindent\textbf{Chú giải \& Lưu ý kỹ thuật:}
\begin{itemize}[leftmargin=1.5em, itemsep=2pt]
  \item Kinh nghiệm chọn vi phân dx hay dy: Nếu đường biên trái và biên phải là hàm đơn giản theo y (ví dụ parabol nằm ngang x = $y^2$), việc lấy tích phân theo dy giúp tránh phải chia miền thành nhiều tích phân nhỏ.
\end{itemize}

\begin{vidu}{Tính diện tích miền phẳng giới hạn bởi parabol và đường thẳng}
\textbf{Đề bài:} Tính diện tích miền phẳng D giới hạn bởi parabol y = $x^2$ - 2x và đường thẳng y = x + 4.
\[ y = x^2 - 2x \quad \text{và} \quad y = x + 4 \]
\textbf{Lời giải chi tiết:}
\begin{enumerate}[leftmargin=1.5em, itemsep=2pt]
  \item Bước 1: Giải phương trình hoành độ giao điểm $x^2$ - 2x = x + 4 $\iff x^2$ - 3x - 4 = 0 $\iff$  x = -1 hoặc x = 4.
  \item Bước 2: Trên đoạn [-1, 4], đường thẳng y = x + 4 nằm phía trên parabol y = $x^2$ - 2x.
  \item Bước 3: Thiết lập và tính tích phân xác định.
\end{enumerate}
\[ S = \int_{-1}^4 \left[(x + 4) - (x^2 - 2x)\right] dx = \int_{-1}^4 (-x^2 + 3x + 4)\,dx = \left[-\frac{x^3}{3} + \frac{3x^2}{2} + 4x\right]_{-1}^4 = \frac{125}{6} \]
\noindent\textbf{Kết luận:} Diện tích miền phẳng D bằng 125/6 (đơn vị diện tích).
\end{vidu}

\section{[Mục 6.2 - 6.3] Thể tích vật thể tròn xoay: Phương pháp Đĩa/Vòng đệm \& Vỏ hình trụ}
Có hai phương pháp chuẩn để tính thể tích khối tròn xoay: (1) Phương pháp Đĩa / Vòng đệm (Disk / Washer Method) cắt vật thể bởi các mặt phẳng vuông góc với trục quay; (2) Phương pháp Vỏ hình trụ (Cylindrical Shells) chia vật thể thành các lớp vỏ hình trụ đồng tâm song song với trục quay.

\begin{congthuc}{Các công thức trọng tâm (6.2 - 6.3)}
\noindent\textbf{Phương pháp Vòng đệm (Quay quanh trục Ox):}
\[ V = \pi \int_a^b \left( [R_{\text{ngoài}}(x)]^2 - [r_{\text{trong}}(x)]^2 \right) dx \]
\noindent\textit{Chú giải:} Nếu miền phẳng áp sát trục quay $(r_trong$ = 0), công thức trở thành phương pháp Đĩa: V = $\pi \int [R(x)]^2$ dx.\\[4pt]
\noindent\textbf{Phương pháp Vỏ hình trụ (Quay quanh trục Oy):}
\[ V = 2\pi \int_a^b x \cdot f(x)\,dx \quad \left(\text{Chu vi } 2\pi x \times \text{Chiều cao } f(x) \times \text{Bề dày } dx\right) \]
\noindent\textit{Chú giải:} Nếu quay quanh đường thẳng đứng x = k (với k < a), bán kính vỏ trụ là (x - k).\\[4pt]
\end{congthuc}

\noindent\textbf{Chú giải \& Lưu ý kỹ thuật:}
\begin{itemize}[leftmargin=1.5em, itemsep=2pt]
  \item Khi hàm y = f(x) khó rút ngược x theo y (ví dụ y = $sin(x^2$) hoặc y = $x^3$ + x) mà yêu cầu quay quanh trục tung Oy, phương pháp Vỏ hình trụ là lựa chọn tối ưu nhất.
\end{itemize}

\begin{vidu}{Tính thể tích khối tròn xoay bằng cả hai phương pháp}
\textbf{Đề bài:} Cho miền D giới hạn bởi đường cong y = 2x - $x^2$ và trục hoành y = 0. Tính thể tích khối tròn xoay tạo thành khi quay D quanh trục tung Oy.
\[ D: \; y = 2x - x^2, \; y = 0 \quad (0 \le x \le 2), \text{ quay quanh } Oy \]
\textbf{Lời giải chi tiết:}
\begin{enumerate}[leftmargin=1.5em, itemsep=2pt]
  \item Bước 1: Giao điểm của y = 2x - $x^2$ với y = 0 là x = 0 và x = 2.
  \item Bước 2: Áp dụng phương pháp Vỏ hình trụ quay quanh Oy với bán kính r = x và chiều cao h = 2x - $x^2$.
  \item Bước 3: Khai triển đa thức và tính tích phân từ 0 đến 2.
\end{enumerate}
\[ V_{Oy} = 2\pi \int_0^2 x(2x - x^2)\,dx = 2\pi \int_0^2 (2x^2 - x^3)\,dx = 2\pi \left[\frac{2x^3}{3} - \frac{x^4}{4}\right]_0^2 = 2\pi\left(\frac{16}{3} - 4\right) = \frac{8\pi}{3} \]
\noindent\textbf{Kết luận:} Thể tích khối tròn xoay tạo thành là V = 8$\pi$ /3.
\end{vidu}

\section{[Mục 6.4 - 6.5] Công của lực biến thiên \& Giá trị trung bình của hàm số}
Khi lực tác dụng F(x) biến thiên dọc theo trục chuyển động từ x = a đến x = b (như lực đàn hồi lò xo định luật Hooke F(x) = kx hoặc công bơm chất lỏng ra khỏi bể), công cơ học W được tính bằng tích phân của F(x). Giá trị trung bình $f_ave$ của hàm liên tục trên [a, b] là chiều cao của hình chữ nhật có cùng diện tích với miền dưới đồ thị.

\begin{congthuc}{Các công thức trọng tâm (6.4 - 6.5)}
\noindent\textbf{Công cơ học của lực biến thiên F(x):}
\[ W = \int_a^b F(x)\,dx \quad (\text{Định luật Hooke cho lò xo: } F(x) = kx) \]
\noindent\textit{Chú giải:} Đơn vị chuẩn SI: Lực F tính bằng Newton (N), quãng đường x tính bằng mét (m), công W tính bằng Joule (J = N$\cdot$ m).\\[4pt]
\noindent\textbf{Giá trị trung bình của hàm số \& Định lý Giá trị Trung bình cho Tích phân:}
\[ f_{\text{ave}} = \frac{1}{b - a} \int_a^b f(x)\,dx, \qquad \exists c \in [a, b] : f(c) = f_{\text{ave}} \]
\noindent\textit{Chú giải:} Diện tích miền dưới đường cong đúng bằng diện tích hình chữ nhật: $\int_a^b$ f(x)dx = f(c)(b - a).\\[4pt]
\end{congthuc}

\noindent\textbf{Chú giải \& Lưu ý kỹ thuật:}
\begin{itemize}[leftmargin=1.5em, itemsep=2pt]
  \item Trong bài toán bơm nước ra khỏi bể: Chia khối nước thành các lớp mỏng nằm ngang ở độ sâu y có diện tích mặt cắt A(y), thể tích dV = A(y)dy, trọng lượng dF = $\rho$ g A(y)dy, và quãng đường nâng lớp nước đó lên miệng bể là h(y).
\end{itemize}

\begin{vidu}{Tính công kéo giãn lò xo và giá trị trung bình của hàm số}
\textbf{Đề bài:} Một lò xo có chiều dài tự nhiên 20 cm. Cần một lực 30 N để giữ lò xo ở độ dài 25 cm. Tính công cần thiết để kéo lò xo từ độ dài 25 cm đến độ dài 30 cm.
\[ F(x) = kx, \quad F(0.05) = 30 \text{ N} \implies W = \int_{0.05}^{0.10} kx\,dx \]
\textbf{Lời giải chi tiết:}
\begin{enumerate}[leftmargin=1.5em, itemsep=2pt]
  \item Bước 1: Đổi đơn vị sang mét: Độ giãn khi dài 25 cm là $x_1$ = 0.05 m; khi dài 30 cm là $x_2$ = 0.10 m.
  \item Bước 2: Tìm độ cứng k từ F(0.05) = 30 $\Rightarrow$  k $\cdot$  0.05 = 30 $\Rightarrow$  k = 600 (N/m).
  \item Bước 3: Tính tích phân công W từ $x_1$ = 0.05 m đến $x_2$ = 0.10 m.
\end{enumerate}
\[ W = \int_{0.05}^{0.10} 600x\,dx = \Big[300x^2\Big]_{0.05}^{0.10} = 300(0.01 - 0.0025) = 300(0.0075) = 2.25 \text{ (J)} \]
\noindent\textbf{Kết luận:} Công cơ học cần thực hiện là W = 2.25 Joule.
\end{vidu}

\section{Minh họa Hình học bằng TikZ}
\begin{figure}[htbp]
\centering
\begin{tikzpicture}[scale=1.05, >=stealth]
  % Trục tọa độ
  \draw[->, thick] (-0.5,0) -- (5.2,0) node[right] {$x$};
  \draw[->, thick] (0,-2.2) -- (0,2.4) node[above] {$y$};
  
  % Đường biên trên và dưới của vật thể tròn xoay
  \draw[domain=0.8:4.2, smooth, variable=\x, blue!80!black, very thick] 
    plot ({\x}, {0.65*sqrt(\x) + 0.35}) node[above right] {$y = f(x)$};
  \draw[domain=0.8:4.2, smooth, variable=\x, blue!60!black, thick, dashed] 
    plot ({\x}, {-0.65*sqrt(\x) - 0.35});
    
  % Hai mặt đáy elip ở x = a và x = b
  \draw[blue!80!black, thick] (0.8,0) ellipse (0.18 and 0.93);
  \draw[blue!80!black, thick, fill=blue!8, fill opacity=0.4] (4.2,0) ellipse (0.32 and 1.68);
  
  % Lát cắt vi phân hình đĩa tại x
  \filldraw[fill=red!25, draw=red!80!black, thick, fill opacity=0.7] (2.5,0) ellipse (0.25 and 1.38);
  \draw[->, red!80!black, thick] (2.5,0) -- (2.5,1.38) node[midway, right] {$R(x)$};
  \node[below] at (0.8,-1.0) {$x=a$};
  \node[below] at (4.2,-1.75) {$x=b$};
\end{tikzpicture}
\caption{Hình 6.1: Mô hình Thể tích vật thể tròn xoay quay quanh trục Ox với lát cắt vi phân hình đĩa bán kính R = f(x).}
\label{fig:tikz-ch6-solid}
\end{figure}

\section{Ứng dụng Thực tế \& Tối ưu hóa (Kỹ thuật Điện xoay chiều \& Cơ khí Chế tạo (CAD))}
\begin{ungdung}{Giá trị Hiệu dụng (RMS) của Dòng điện Xoay chiều và Tính Thể tích Chi tiết Tiện CNC}
Giá trị trung bình của dòng điện hình sin i(t) = $I_0$ sin($\omega$ t) trong một chu kỳ bằng 0, do đó trong kỹ thuật điện người ta dùng Giá trị Hiệu dụng (Root Mean Square - RMS), được định nghĩa bằng căn bậc hai của giá trị trung bình của hàm bình phương $i^2$(t) trên một chu kỳ T.
\[ I_{\text{RMS}} = \sqrt{\frac{1}{T} \int_0^T [I_0 \sin(\omega t)]^2\,dt} = \frac{I_0}{\sqrt{2}} \approx 0.707 I_0 \]
\noindent\textbf{Ví dụ tính toán thực tế:} Điện lưới dân dụng tại Việt Nam có điện áp hiệu dụng $U_RMS$ = 220V, nghĩa là biên độ điện áp cực đại thực tế đạt tới $U_0$ = 220$\sqrt{2} \approx$  311.13 V.
\end{ungdung}

\section{Bài tập Trắc nghiệm Củng cố Chương 6}
\begin{enumerate}[label=\textbf{Câu \arabic*:}, leftmargin=2.2em, itemsep=8pt]
  \item Diện tích miền phẳng giới hạn bởi hai đường cong y = $x^2$ và y = $\sqrt{x}$ trên đoạn [0, 1] bằng:
  \begin{itemize}[label={}, leftmargin=1em, itemsep=2pt]
    \item \textbf{A.} 1/3 \quad ($\frac{1}{3}$)
    \item \textbf{B.} 1/2 \quad ($\frac{1}{2}$)
    \item \textbf{C.} 2/3 \quad ($\frac{2}{3}$)
    \item \textbf{D.} 1/6 \quad ($\frac{1}{6}$)
  \end{itemize}
  \vspace{2pt}
  \noindent{\color{teal!85!black}\textbf{Đáp án đúng: A.}} \textit{Giải thích:} Trên đoạn [0, 1], ta có $\sqrt{x} \ge x^2$. Do đó S = $\int_0^1$ ($\sqrt{x} - x^2$)dx = [(2/$3)x^{3/2} - x^3 / 3]_0^1$ = 2/3 - 1/3 = 1/3.
  \[ S = \int_0^1 (\sqrt{x} - x^2)\,dx = \left[\frac{2}{3}x^{3/2} - \frac{x^3}{3}\right]_0^1 = \frac{1}{3} \]
  \item Thể tích khối tròn xoay khi quay miền giới hạn bởi y = $\sqrt{x}$, y = 0, x = 4 quanh trục hoành Ox là:
  \begin{itemize}[label={}, leftmargin=1em, itemsep=2pt]
    \item \textbf{A.} 8$\pi$  \quad ($8\pi$)
    \item \textbf{B.} 4$\pi$  \quad ($4\pi$)
    \item \textbf{C.} 16$\pi$  \quad ($16\pi$)
    \item \textbf{D.} 8 \quad ($8$)
  \end{itemize}
  \vspace{2pt}
  \noindent{\color{teal!85!black}\textbf{Đáp án đúng: A.}} \textit{Giải thích:} Áp dụng công thức phương pháp Đĩa: V = $\pi \int_0^4$ ($\sqrt{x} )^2$ dx = $\pi \int_0^4$ x dx = $\pi  [x^2 / 2]_0^4$ = 8$\pi$ .
  \[ V = \pi \int_0^4 (\sqrt{x})^2\,dx = \pi \left[\frac{x^2}{2}\right]_0^4 = 8\pi \]
  \item Giá trị trung bình của hàm số f(x) = $3x^2$ - 2x trên đoạn [1, 4] bằng:
  \begin{itemize}[label={}, leftmargin=1em, itemsep=2pt]
    \item \textbf{A.} 16 \quad ($16$)
    \item \textbf{B.} 48 \quad ($48$)
    \item \textbf{C.} 12 \quad ($12$)
    \item \textbf{D.} 21 \quad ($21$)
  \end{itemize}
  \vspace{2pt}
  \noindent{\color{teal!85!black}\textbf{Đáp án đúng: A.}} \textit{Giải thích:} Ta có $f_ave$ = [1 / (4 - 1)] $\int_1^4 (3x^2$ - 2x)dx = (1/3) $[x^3 - x^2 ]_1^4$ = (1/3) [(64 - 16) - (1 - 1)] = 48/3 = 16.
  \[ f_{\text{ave}} = \frac{1}{4-1} \Big[x^3 - x^2\Big]_1^4 = \frac{48 - 0}{3} = 16 \]
\end{enumerate}

% ======================================================================
% CHƯƠNG 7: CÁC PHƯƠNG PHÁP TÍNH TÍCH PHÂN & TÍCH PHÂN SUY RỘNG (TRỌNG TÂM) (Buổi 7 - 8)
% ======================================================================
\chapter{Các phương pháp tính Tích phân \& Tích phân suy rộng (Trọng tâm)}
\label{chap:chuong-7}

\begin{tcolorbox}[colback=slate!5!white,colframe=slate!75!black,title={\textbf{Mục lục Chương 7 \& Chuẩn đầu ra (Buổi 7 - 8)}}]
\textbf{Tổng quan:} Chương này hoàn thiện hệ thống kỹ thuật tính tích phân nâng cao: Tích phân từng phần, Tích phân hàm lượng giác, Phép thế lượng giác cho căn thức bậc hai, Phân tích hàm hữu tỉ thành tổng các phân thức đơn giản, và mở rộng khái niệm tích phân sang miền không bị chặn hoặc hàm không bị chặn (Tích phân suy rộng loại 1 và loại 2).\\[6pt]
\textbf{Nội dung chính trong chương:}
\begin{itemize}[leftmargin=1.5em, itemsep=2pt]
  \item \textbf{Mục 7.1:} Phương pháp Tích phân từng phần (Integration by Parts)
  \item \textbf{Mục 7.2 - 7.3:} Tích phân hàm lượng giác \& Phép thay thế lượng giác
  \item \textbf{Mục 7.4 - 7.5:} Tích phân hàm hữu tỉ (Đồng nhất hệ số) \& Chiến lược tính tích phân
  \item \textbf{Mục 7.6:} Tích phân suy rộng (Loại 1: Cận vô cực \& Loại 2: Hàm không bị chặn)
\end{itemize}
\textbf{Chuẩn đầu ra (LLO):}
\begin{itemize}[leftmargin=1.5em, itemsep=2pt]
  \item LLO7.1: Áp dụng thành thạo tích phân từng phần, lượng giác, thế lượng giác và phân tích phân thức hữu tỉ.
  \item LLO8.1: Hiểu định nghĩa tích phân suy rộng loại 1, loại 2 và xét được tính hội tụ / phân kỳ của chúng.
\end{itemize}
\end{tcolorbox}

\section{[Mục 7.1] Phương pháp Tích phân từng phần (Integration by Parts)}
Được suy ra trực tiếp từ quy tắc đạo hàm của tích (uv)' = u'v + uv'. Quy tắc chọn u theo thứ tự ưu tiên LIATE (tiếng Việt: "Nhất Lô - Nhì Đa - Tam Lượng - Tứ Mũ"): Logarithm $\to$  Đa thức $\to$  Lượng giác $\to$  Hàm mũ.

\begin{congthuc}{Các công thức trọng tâm (7.1)}
\noindent\textbf{Công thức Tích phân từng phần:}
\[ \int u\,dv = u v - \int v\,du, \qquad \int_a^b u\,dv = \Big[u v\Big]_a^b - \int_a^b v\,du \]
\noindent\textit{Chú giải:} Với dạng tích phân từng phần lặp I = $\int e^{ax}$ sin(bx)dx hoặc $\int e^{ax}$ cos(bx)dx, từng phần 2 lần sẽ thu lại I ở vế phải.\\[4pt]
\end{congthuc}

\noindent\textbf{Chú giải \& Lưu ý kỹ thuật:}
\begin{itemize}[leftmargin=1.5em, itemsep=2pt]
  \item Khi tính tích phân của một hàm đứng độc lập như $\int$ ln(x)dx, $\int$ arctan(x)dx, hoặc $\int$ arcsin(x)dx, ta luôn đặt u bằng chính hàm đó và dv = dx.
\end{itemize}

\begin{vidu}{Tính tích phân từng phần chứa hàm Logarithm và Đa thức}
\textbf{Đề bài:} Tính tích phân xác định I = $\int_1^e$ x $\ln$ x\,dx.
\[ I = \int_1^e x \ln x\,dx \]
\textbf{Lời giải chi tiết:}
\begin{enumerate}[leftmargin=1.5em, itemsep=2pt]
  \item Bước 1: Theo quy tắc "Nhất Lô - Nhì Đa", đặt u = ln x $\Rightarrow$  du = dx/x và dv = x dx $\Rightarrow$  v = $x^2$ / 2.
  \item Bước 2: Áp dụng công thức tích phân từng phần trên đoạn [1, e].
\end{enumerate}
\[ I = \left[ \frac{x^2}{2} \ln x \right]_1^e - \int_1^e \frac{x^2}{2} \cdot \frac{1}{x}\,dx = \frac{e^2}{2} - \left[\frac{x^2}{4}\right]_1^e = \frac{e^2}{2} - \left(\frac{e^2}{4} - \frac{1}{4}\right) = \frac{e^2 + 1}{4} \]
\noindent\textbf{Kết luận:} Giá trị của tích phân là I = $(e^2$ + 1)/4.
\end{vidu}

\section{[Mục 7.2 - 7.3] Tích phân hàm lượng giác \& Phép thay thế lượng giác}
Khi tích phân chứa các căn thức dạng $\sqrt{a^2 - x^2}$, $\sqrt{a^2 + x^2}$, hoặc $\sqrt{x^2 - a^2}$, ta dùng phép thế lượng giác để khử dấu căn nhờ các hằng đẳng thức Pythagoras: 1 - $sin^2 \theta = cos^2 \theta$  và 1 + $tan^2 \theta = sec^2 \theta$ .

\begin{congthuc}{Các công thức trọng tâm (7.2 - 7.3)}
\noindent\textbf{Bảng 3 phép thế lượng giác chuẩn (với a > 0):}
\[ \begin{aligned} &\sqrt{a^2 - x^2} \implies \text{Đặt } x = a\sin\theta, \; dx = a\cos\theta\,d\theta, \; \theta \in [-\pi/2, \pi/2] \\ &\sqrt{a^2 + x^2} \implies \text{Đặt } x = a\tan\theta, \; dx = a\sec^2\theta\,d\theta, \; \theta \in (-\pi/2, \pi/2) \\ &\sqrt{x^2 - a^2} \implies \text{Đặt } x = a\sec\theta, \; dx = a\sec\theta\tan\theta\,d\theta \end{aligned} \]
\noindent\textit{Chú giải:} Với tích phân $\int sin^m$(x) $cos^n$(x)dx: Nếu m hoặc n lẻ thì tách 1 bậc lẻ đưa vào vi phân; nếu cả m, n đều chẵn thì dùng công thức hạ bậc.\\[4pt]
\end{congthuc}

\noindent\textbf{Chú giải \& Lưu ý kỹ thuật:}
\begin{itemize}[leftmargin=1.5em, itemsep=2pt]
  \item Công thức hạ bậc lượng giác cần nhớ: $sin^2$ x = (1 - cos 2x)/2 và $cos^2$ x = (1 + cos 2x)/2.
\end{itemize}

\begin{vidu}{Tính tích phân chứa căn thức bằng phép thế lượng giác}
\textbf{Đề bài:} Tính tích phân I = $\int_0^{3/2} \frac{dx}{(9 - x^2)^{3/2}}$.
\[ I = \int_0^{3/2} \frac{dx}{(\sqrt{9 - x^2})^3} \]
\textbf{Lời giải chi tiết:}
\begin{enumerate}[leftmargin=1.5em, itemsep=2pt]
  \item Bước 1: Nhận diện dạng $a^2 - x^2$ với a = 3. Đặt x = 3 sin $\theta \Rightarrow$  dx = 3 cos $\theta$  d$\theta$ .
  \item Bước 2: Đổi cận: x = 0 $\Rightarrow \theta$  = 0; x = 3/2 $\Rightarrow$  sin $\theta$  = 1/2 $\Rightarrow \theta = \pi$ /6.
  \item Bước 3: Thay vào tích phân, rút gọn $cos^2 \theta$  ở mẫu thành $sec^2 \theta$  = (tan $\theta$ )'.
\end{enumerate}
\[ I = \int_0^{\pi/6} \frac{3\cos\theta\,d\theta}{(3\cos\theta)^3} = \frac{1}{9} \int_0^{\pi/6} \frac{d\theta}{\cos^2\theta} = \frac{1}{9}\Big[\tan\theta\Big]_0^{\pi/6} = \frac{1}{9} \cdot \frac{1}{\sqrt{3}} = \frac{\sqrt{3}}{27} \]
\noindent\textbf{Kết luận:} Giá trị của tích phân là I = $\sqrt{3}$ / 27.
\end{vidu}

\section{[Mục 7.4 - 7.5] Tích phân hàm hữu tỉ (Đồng nhất hệ số) \& Chiến lược tính tích phân}
Mọi hàm phân thức hữu tỉ thực sự P(x)/Q(x) (với bậc của tử P nhỏ hơn bậc của mẫu Q) đều có thể phân tích thành tổng của các phân thức đơn giản dạng A/(x - $a)^k$ và (Bx + C)/$(x^2$ + px + $q)^m$ (trong đó $\Delta = p^2$ - 4q < 0).

\begin{congthuc}{Các công thức trọng tâm (7.4 - 7.5)}
\noindent\textbf{Phân tích phân thức hữu tỉ (Phương pháp Heaviside che nghiệm):}
\[ \frac{P(x)}{(x - x_1)(x - x_2)} = \frac{A}{x - x_1} + \frac{B}{x - x_2}, \quad A = \frac{P(x_1)}{x_1 - x_2}, \; B = \frac{P(x_2)}{x_2 - x_1} \]
\noindent\textit{Chú giải:} Nếu bậc của P(x) $\ge$  bậc của Q(x), bắt buộc phải thực hiện phép chia đa thức trước khi tách phân thức.\\[4pt]
\end{congthuc}

\noindent\textbf{Chú giải \& Lưu ý kỹ thuật:}
\begin{itemize}[leftmargin=1.5em, itemsep=2pt]
  \item Với phân thức có mẫu bậc 2 vô nghiệm thực: Đưa mẫu về dạng bình phương đủ (x + p/$2)^2 + a^2$ để tách thành ln(mẫu) và arctan.
\end{itemize}

\begin{vidu}{Tích phân hàm phân thức hữu tỉ bằng tách phân thức đơn giản}
\textbf{Đề bài:} Tính tích phân bất định I = $\int \frac{x + 5}{x^2 + x - 2}$\,dx.
\[ I = \int \frac{x + 5}{(x - 1)(x + 2)}\,dx \]
\textbf{Lời giải chi tiết:}
\begin{enumerate}[leftmargin=1.5em, itemsep=2pt]
  \item Bước 1: Phân tích mẫu số $x^2$ + x - 2 = (x - 1)(x + 2). Giả sử (x + 5)/[(x - 1)(x + 2)] = A/(x - 1) + B/(x + 2).
  \item Bước 2: Dùng quy tắc che nghiệm Heaviside: Tại x = 1 $\Rightarrow$  A = (1 + 5)/(1 + 2) = 2; tại x = -2 $\Rightarrow$  B = (-2 + 5)/(-2 - 1) = -1.
  \item Bước 3: Lấy nguyên hàm từng phân thức đơn giản.
\end{enumerate}
\[ I = \int \left( \frac{2}{x - 1} - \frac{1}{x + 2} \right) dx = 2\ln|x - 1| - \ln|x + 2| + C = \ln\left|\frac{(x - 1)^2}{x + 2}\right| + C \]
\noindent\textbf{Kết luận:} Nguyên hàm là I = 2 ln|x - 1| - ln|x + 2| + C.
\end{vidu}

\section{[Mục 7.6] Tích phân suy rộng (Loại 1: Cận vô cực \& Loại 2: Hàm không bị chặn)}
Tích phân suy rộng mở rộng tích phân xác định bằng cách lấy giới hạn khi cận tiến ra vô cực (Loại 1) hoặc khi cận tiến tới điểm kỳ dị làm hàm số tiến ra vô cực (Loại 2). Nếu giới hạn tồn tại hữu hạn thì tích phân HỘI TỤ, ngược lại là PHÂN KỲ.

\begin{congthuc}{Các công thức trọng tâm (7.6)}
\noindent\textbf{Định nghĩa Tích phân suy rộng Loại 1 và Chuẩn p-tích phân tại vô cực:}
\[ \int_a^\infty f(x)\,dx = \lim_{t \to \infty} \int_a^t f(x)\,dx, \qquad \int_1^\infty \frac{dx}{x^p} \;\begin{cases} \text{Hội tụ bằng } \frac{1}{p-1} & \text{khi } p > 1 \\ \text{Phân kỳ} & \text{khi } p \le 1 \end{cases} \]
\noindent\textit{Chú giải:} Với tích phân suy rộng 2 đầu từ -$\infty$ đến +$\infty$, phải tách tại 1 điểm trung gian c và cả 2 nhánh phải độc lập hội tụ.\\[4pt]
\noindent\textbf{Tích phân suy rộng Loại 2 (Điểm kỳ dị tại x = a) \& Chuẩn hội tụ:}
\[ \int_a^b \frac{dx}{(x - a)^q} \;\begin{cases} \text{Hội tụ} & \text{khi } q < 1 \\ \text{Phân kỳ} & \text{khi } q \ge 1 \end{cases} \quad (\text{Ngược chiều với chuẩn } p \text{ ở vô cực!}) \]
\noindent\textit{Chú giải:} Tiêu chuẩn so sánh: Với 0 $\le$  f(x) $\le$  g(x), nếu $\int$ g(x)dx hội tụ thì $\int$ f(x)dx hội tụ.\\[4pt]
\end{congthuc}

\noindent\textbf{Chú giải \& Lưu ý kỹ thuật:}
\begin{itemize}[leftmargin=1.5em, itemsep=2pt]
  \item Cảnh báo bẫy thường gặp: Khi tính $\int_{-1}^1$ (1/$x^2$)dx, không được áp dụng trực tiếp công thức Newton–Leibniz [-1/$x]_{-1}^1$ = -2 vì x = 0 là điểm kỳ dị nằm giữa đoạn [-1, 1] làm tích phân PHÂN KỲ.
\end{itemize}

\begin{vidu}{Xét sự hội tụ và tính giá trị tích phân suy rộng}
\textbf{Đề bài:} Tính các tích phân suy rộng: (a) $I_1 = \int_1^{+\infty}$ x $e^{-x}$\,dx và (b) $I_2 = \int_0^1 \frac{dx}{\sqrt{1 - x}}$.
\[ I_1 = \int_1^{+\infty} x e^{-x}\,dx, \qquad I_2 = \int_0^1 \frac{dx}{\sqrt{1 - x}} \]
\textbf{Lời giải chi tiết:}
\begin{enumerate}[leftmargin=1.5em, itemsep=2pt]
  \item Câu (a): Là tích phân suy rộng loại 1. Tính nguyên hàm từng phần $\int$ x $e^{-x}$ dx = -(x + $1)e^{-x}$, sau đó lấy giới hạn khi t $\to + \infty$.
  \item Câu (b): Là tích phân suy rộng loại 2 có điểm kỳ dị tại x = 1 (với q = 1/2 < 1 nên hội tụ).
\end{enumerate}
\[ I_1 = \lim_{t \to +\infty} \Big[ -(x + 1)e^{-x} \Big]_1^t = \lim_{t \to +\infty} \left( -\frac{t + 1}{e^t} + \frac{2}{e} \right) = 0 + \frac{2}{e} = \frac{2}{e} \;(\text{Hội tụ}) \]
\[ I_2 = \lim_{t \to 1^-} \Big[ -2\sqrt{1 - x} \Big]_0^t = \lim_{t \to 1^-} \left( -2\sqrt{1 - t} + 2 \right) = 2 \;(\text{Hội tụ}) \]
\noindent\textbf{Kết luận:} Cả hai tích phân suy rộng đều hội tụ với $I_1$ = 2/e và $I_2$ = 2.
\end{vidu}

\section{Minh họa Hình học bằng TikZ}
\begin{figure}[htbp]
\centering
\begin{tikzpicture}[scale=1.1, >=stealth]
  % Miền diện tích hội tụ dưới đường y = 1/x^2
  \fill[teal!20] (1,0) -- plot[domain=1:4.6, smooth, variable=\x] ({\x}, {1.6/(\x*\x)}) -- (4.6,0) -- cycle;
  
  % Trục tọa độ
  \draw[->, thick] (-0.3,0) -- (5.0,0) node[right] {$x$};
  \draw[->, thick] (0,-0.3) -- (0,3.2) node[above] {$y$};
  
  % Đồ thị y = 1/x (Phân kỳ) và y = 1/x^2 (Hội tụ)
  \draw[domain=0.55:4.6, smooth, variable=\x, red!80!black, very thick, dashed] 
    plot ({\x}, {1.6/(\x)}) node[above right] {$y = \frac{1}{x}$ (Phân kỳ)};
  \draw[domain=0.75:4.6, smooth, variable=\x, teal!85!black, very thick] 
    plot ({\x}, {1.6/(\x*\x)}) node[right] {$y = \frac{1}{x^2}$ (Hội tụ)};
    
  \draw[thick] (1,0) node[below] {$1$} -- (1,1.6);
  \node[teal!85!black] at (2.1, 0.45) {$\int_1^\infty \frac{dx}{x^2} = 1$};
\end{tikzpicture}
\caption{Hình 7.1: So sánh tốc độ suy giảm của 1/x (Tích phân suy rộng phân kỳ) và 1/$x^2$ (Hội tụ về diện tích hữu hạn bằng 1) trên [1, +$\infty$).}
\label{fig:tikz-ch7-parts}
\end{figure}

\section{Ứng dụng Thực tế \& Tối ưu hóa (Xử lý Tín hiệu Số (DSP) \& Lý thuyết Điều khiển)}
\begin{ungdung}{Biến đổi Laplace và Biến đổi Fourier qua Tích phân Suy rộng}
Trong kỹ thuật điều khiển và viễn thông, phép Biến đổi Laplace L{f(t)} chuyển phương trình vi phân theo thời gian t sang phương trình đại số trong miền tần số phức s thông qua một tích phân suy rộng loại 1 kết hợp kỹ thuật tích phân từng phần.
\[ \mathcal{L}\{f(t)\}(s) = \int_0^{+\infty} f(t) e^{-st}\,dt, \qquad \mathcal{L}\{f'(t)\} = s\mathcal{L}\{f(t)\} - f(0) \]
\noindent\textbf{Ví dụ tính toán thực tế:} Biến đổi Laplace của tín hiệu bậc thang đơn vị f(t) = 1 (với s > 0) là $\int_0^\infty e^{-st}$dt = $\lim_{T\to \infty} [-e^{-st}/s]_0^T$ = 1/s.
\end{ungdung}

\section{Bài tập Trắc nghiệm Củng cố Chương 7}
\begin{enumerate}[label=\textbf{Câu \arabic*:}, leftmargin=2.2em, itemsep=8pt]
  \item Nguyên hàm I = $\int \arctan$ x\,dx bằng:
  \begin{itemize}[label={}, leftmargin=1em, itemsep=2pt]
    \item \textbf{A.} x arctan(x) - (1/2) ln(1 + $x^2$) + C \quad ($x\arctan x - \frac{1}{2}\ln(1+x^2) + C$)
    \item \textbf{B.} x arctan(x) + (1/2) ln(1 + $x^2$) + C \quad ($x\arctan x + \frac{1}{2}\ln(1+x^2) + C$)
    \item \textbf{C.} 1 / (1 + $x^2$) + C \quad ($\frac{1}{1+x^2} + C$)
    \item \textbf{D.} x / (1 + $x^2$) - arctan(x) + C \quad ($\frac{x}{1+x^2} - \arctan x + C$)
  \end{itemize}
  \vspace{2pt}
  \noindent{\color{teal!85!black}\textbf{Đáp án đúng: A.}} \textit{Giải thích:} Đặt u = arctan x $\Rightarrow$  du = dx/(1 + $x^2$) và dv = dx $\Rightarrow$  v = x. Khi đó I = x arctan x - $\int [x / (1 + x^2)]$dx = x arctan x - (1/2)ln(1 + $x^2$) + C.
  \[ I = x\arctan x - \int \frac{x}{1+x^2}\,dx = x\arctan x - \frac{1}{2}\ln(1+x^2) + C \]
  \item Trong các tích phân suy rộng sau đây, tích phân nào HỘI TỤ?
  \begin{itemize}[label={}, leftmargin=1em, itemsep=2pt]
    \item \textbf{A.} $\int_1^{+\infty}$ dx / $x^{3/2}$ \quad ($\int_1^{+\infty} \frac{dx}{x^{3/2}}$)
    \item \textbf{B.} $\int_1^{+\infty}$ dx / $\sqrt{x}$ \quad ($\int_1^{+\infty} \frac{dx}{\sqrt{x}}$)
    \item \textbf{C.} $\int_0^1$ dx / $x^2$ \quad ($\int_0^1 \frac{dx}{x^2}$)
    \item \textbf{D.} $\int_0^{+\infty}$ cos(x) dx \quad ($\int_0^{+\infty} \cos x\,dx$)
  \end{itemize}
  \vspace{2pt}
  \noindent{\color{teal!85!black}\textbf{Đáp án đúng: A.}} \textit{Giải thích:} Tích phân loại 1 dạng $\int_1^\infty$ dx/$x^p$ hội tụ khi và chỉ khi p > 1. Ở đáp án A, p = 3/2 = 1.5 > 1 nên tích phân hội tụ (và bằng 2). Đáp án B có p = 1/2 < 1 nên phân kỳ; đáp án C là loại 2 tại x = 0 với q = 2 $\ge$  1 nên phân kỳ.
  \[ p = \frac{3}{2} > 1 \implies \int_1^{+\infty} x^{-3/2}\,dx = \Big[-2x^{-1/2}\Big]_1^{+\infty} = 2 \]
  \item Để tính tích phân I = $\int \sqrt{4 - x^2}$\,dx, phép đổi biến lượng giác chuẩn là:
  \begin{itemize}[label={}, leftmargin=1em, itemsep=2pt]
    \item \textbf{A.} Đặt x = 2 sin $\theta$  với $\theta \in  [-\pi /2, \pi /2]$ \quad ($x = 2\sin\theta$)
    \item \textbf{B.} Đặt x = 2 tan $\theta$  với $\theta \in$  (-$\pi$ /2, $\pi$ /2) \quad ($x = 2\tan\theta$)
    \item \textbf{C.} Đặt x = 4 sin $\theta$  \quad ($x = 4\sin\theta$)
    \item \textbf{D.} Đặt x = 2 / cos $\theta$  \quad ($x = 2\sec\theta$)
  \end{itemize}
  \vspace{2pt}
  \noindent{\color{teal!85!black}\textbf{Đáp án đúng: A.}} \textit{Giải thích:} Với căn thức $\sqrt{a^2 - x^2}$ trong đó $a^2$ = 4 $\Rightarrow$  a = 2, ta đặt x = 2 sin $\theta$  để $\sqrt{4 - 4sin^2 \theta }$ = 2 cos $\theta$ .
  \[ \sqrt{4 - x^2} = \sqrt{2^2 - x^2} \implies x = 2\sin\theta \]
\end{enumerate}

% ======================================================================
% CHƯƠNG 8: CÁC ỨNG DỤNG KHÁC CỦA TÍCH PHÂN (Buổi 8 - 9)
% ======================================================================
\chapter{Các ứng dụng khác của Tích phân}
\label{chap:chuong-8}

\begin{tcolorbox}[colback=slate!5!white,colframe=slate!75!black,title={\textbf{Mục lục Chương 8 \& Chuẩn đầu ra (Buổi 8 - 9)}}]
\textbf{Tổng quan:} Mở rộng ứng dụng của tích phân sang hình học vi phân (độ dài đường cong, diện tích mặt tròn xoay) và các bài toán liên ngành: Vật lý \& Kỹ thuật (áp lực thủy tĩnh lên đập nước, mô-men và trọng tâm bản phẳng), Kinh tế (thặng dư tiêu dùng), Sinh học (lưu lượng máu Poiseuille \& cung lượng tim) và Lý thuyết Xác suất (hàm mật độ xác suất, kỳ vọng và phân phối chuẩn Gauss).\\[6pt]
\textbf{Nội dung chính trong chương:}
\begin{itemize}[leftmargin=1.5em, itemsep=2pt]
  \item \textbf{Mục 8.1 - 8.2:} Độ dài cung đường cong \& Diện tích mặt tròn xoay
  \item \textbf{Mục 8.3:} Ứng dụng trong Vật lý \& Kỹ thuật: Áp lực thủy tĩnh \& Trọng tâm (Centroid)
  \item \textbf{Mục 8.4 - 8.5:} Ứng dụng trong Kinh tế, Sinh học \& Xác suất (Phân phối chuẩn)
\end{itemize}
\textbf{Chuẩn đầu ra (LLO):}
\begin{itemize}[leftmargin=1.5em, itemsep=2pt]
  \item LLO8.2: Áp dụng tích phân tính độ dài cung và diện tích bề mặt khối tròn xoay.
  \item LLO9.1: Hiểu và ứng dụng tích phân giải các bài toán thực tế trong vật lý (áp suất thủy tĩnh, trọng tâm), kinh tế (thặng dư tiêu dùng), sinh học và xác suất (giá trị trung bình, phân phối chuẩn).
\end{itemize}
\end{tcolorbox}

\section{[Mục 8.1 - 8.2] Độ dài cung đường cong \& Diện tích mặt tròn xoay}
Vi phân độ dài cung ds được xấp xỉ bởi cạnh huyền của tam giác vuông vi phân có hai cạnh góc vuông dx và dy: ds = $\sqrt{dx^2 + dy^2}$. Khi quay cung ds quanh trục tọa độ, nó quét nên một dải hình nón cụt có diện tích vi phân dS = 2$\pi$  r ds.

\begin{congthuc}{Các công thức trọng tâm (8.1 - 8.2)}
\noindent\textbf{Vi phân cung ds và Độ dài cung đường cong y = f(x) trên [a, b]:}
\[ ds = \sqrt{1 + [f'(x)]^2}\,dx \implies L = \int_a^b \sqrt{1 + [f'(x)]^2}\,dx \]
\noindent\textit{Chú giải:} Nếu đường cong cho dưới dạng x = g(y) với c $\le$  y $\le$  d thì ds = $\sqrt{1 + [g'(y} ]^2$) dy.\\[4pt]
\noindent\textbf{Diện tích mặt tròn xoay (Surface Area of Revolution):}
\[ S_{Ox} = 2\pi \int_a^b y\,ds = 2\pi \int_a^b f(x)\sqrt{1 + [f'(x)]^2}\,dx, \qquad S_{Oy} = 2\pi \int_a^b x\,ds \]
\noindent\textit{Chú giải:} Luôn ghi nhớ công thức tổng quát S = $\int$ 2$\pi$ R ds, trong đó R là khoảng cách từ điểm trên cung tới trục quay.\\[4pt]
\end{congthuc}

\noindent\textbf{Chú giải \& Lưu ý kỹ thuật:}
\begin{itemize}[leftmargin=1.5em, itemsep=2pt]
  \item Trong các bài toán tính độ dài cung có kết quả giải tích đẹp, biểu thức 1 + $[f'(x)]^2$ thường có dạng bình phương hoàn hảo (như (a + $bx)^2$ hoặc $cosh^2$ x).
\end{itemize}

\begin{vidu}{Tính độ dài cung của đường cong bán lập phương}
\textbf{Đề bài:} Tính độ dài cung của đường cong y = $\frac{2}{3}$(x - $1)^{3/2}$ từ điểm có hoành độ x = 1 đến x = 4.
\[ y = \frac{2}{3}(x - 1)^{3/2}, \quad 1 \le x \le 4 \]
\textbf{Lời giải chi tiết:}
\begin{enumerate}[leftmargin=1.5em, itemsep=2pt]
  \item Bước 1: Tính đạo hàm y' = (x - $1)^{1/2} = \sqrt{x - 1}$, suy ra 1 + $(y')^2$ = 1 + (x - 1) = x.
  \item Bước 2: Thay vào công thức độ dài cung L = $\int_1^4 \sqrt{x}$ dx.
\end{enumerate}
\[ L = \int_1^4 \sqrt{1 + (x - 1)}\,dx = \int_1^4 x^{1/2}\,dx = \left[\frac{2}{3}x^{3/2}\right]_1^4 = \frac{2}{3}(8 - 1) = \frac{14}{3} \]
\noindent\textbf{Kết luận:} Độ dài cung của đường cong trên đoạn [1, 4] bằng 14/3.
\end{vidu}

\section{[Mục 8.3] Ứng dụng trong Vật lý \& Kỹ thuật: Áp lực thủy tĩnh \& Trọng tâm (Centroid)}
Áp suất thủy tĩnh tại độ sâu d dưới mặt chất lỏng là P = $\rho$ gd. Vì áp suất thay đổi theo độ sâu, tổng lực thủy tĩnh tác dụng lên thành đập thẳng đứng phải được tính bằng tích phân các dải nằm ngang. Trọng tâm ($\bar{x}$, $\bar{y}$) của bản phẳng đồng chất là điểm cân bằng mô-men tĩnh đối với hai trục tọa độ.

\begin{congthuc}{Các công thức trọng tâm (8.3)}
\noindent\textbf{Tổng lực áp suất thủy tĩnh lên bản phẳng thẳng đứng:}
\[ F = \int_a^b \rho g \cdot h(y) \cdot w(y)\,dy \quad (\rho g \text{ của nước } \approx 9800 \text{ N/m}^3) \]
\noindent\textit{Chú giải:} Trong đó h(y) là độ sâu của dải ngang tính từ mặt thoáng và w(y) là chiều rộng của bản tại vị trí y.\\[4pt]
\noindent\textbf{Tọa độ Trọng tâm (Centroid) của miền phẳng D có diện tích A:}
\[ \bar{x} = \frac{1}{A} \int_a^b x[f(x) - g(x)]\,dx, \qquad \bar{y} = \frac{1}{A} \int_a^b \frac{1}{2}\left([f(x)]^2 - [g(x)]^2\right) dx \]
\noindent\textit{Chú giải:} Định lý Pappus: Thể tích khối tròn xoay bằng diện tích miền phẳng nhân với chu vi đường tròn mà trọng tâm vẽ nên (V = A $\cdot$  2$\pi$ d).\\[4pt]
\end{congthuc}

\noindent\textbf{Chú giải \& Lưu ý kỹ thuật:}
\begin{itemize}[leftmargin=1.5em, itemsep=2pt]
  \item Nếu miền phẳng D có trục đối xứng thì trọng tâm ($\bar{x}$, $\bar{y}$) bắt buộc phải nằm trên trục đối xứng đó.
\end{itemize}

\begin{vidu}{Tìm tọa độ trọng tâm của nửa hình tròn}
\textbf{Đề bài:} Tìm tọa độ trọng tâm ($\bar{x}$, $\bar{y}$) của miền phẳng nửa hình tròn trên: y = $\sqrt{R^2 - x^2}$, y $\ge$ 0.
\[ D: \; y = \sqrt{R^2 - x^2} \quad (-R \le x \le R), \quad A = \frac{1}{2}\pi R^2 \]
\textbf{Lời giải chi tiết:}
\begin{enumerate}[leftmargin=1.5em, itemsep=2pt]
  \item Bước 1: Do miền D đối xứng qua trục Oy nên hoành độ trọng tâm $\bar{x}$ = 0.
  \item Bước 2: Áp dụng công thức tung độ trọng tâm $\bar{y}$ = (1/A) $\int_{-R}^R$ (1/$2)[f(x)]^2$ dx.
\end{enumerate}
\[ \bar{y} = \frac{1}{\frac{1}{2}\pi R^2} \int_{-R}^R \frac{1}{2}(R^2 - x^2)\,dx = \frac{2}{\pi R^2} \int_0^R (R^2 - x^2)\,dx = \frac{2}{\pi R^2}\left(\frac{2R^3}{3}\right) = \frac{4R}{3\pi} \]
\noindent\textbf{Kết luận:} Trọng tâm của nửa hình tròn nằm tại điểm (0, 4R / (3$\pi$ )) $\approx$  (0, 0.4244 R).
\end{vidu}

\section{[Mục 8.4 - 8.5] Ứng dụng trong Kinh tế, Sinh học \& Xác suất (Phân phối chuẩn)}
Trong kinh tế học, Thặng dư tiêu dùng (Consumer Surplus) đo tổng lợi ích tiết kiệm được của người mua khi giá thị trường P thấp hơn giá họ sẵn sàng trả p(x). Trong lý thuyết xác suất, biến ngẫu nhiên liên tục X được đặc trưng bởi hàm mật độ xác suất f(x) $\ge$  0 thỏa mãn tích phân trên toàn trục số bằng 1.

\begin{congthuc}{Các công thức trọng tâm (8.4 - 8.5)}
\noindent\textbf{Thặng dư Tiêu dùng (Kinh tế) \& Lưu lượng máu (Sinh học):}
\[ \text{CS} = \int_0^X [p(x) - P]\,dx, \qquad F_{\text{flux}} = \int_0^R 2\pi r \cdot v(r)\,dr = \frac{\pi P R^4}{8\eta l} \;(\text{Định luật Poiseuille}) \]
\noindent\textit{Chú giải:} Lưu lượng máu qua động mạch tỉ lệ thuận với lũy thừa bậc 4 của bán kính mạch máu $R^4$.\\[4pt]
\noindent\textbf{Hàm mật độ xác suất (PDF), Giá trị trung bình (Kỳ vọng $\mu$ ) \& Phân phối chuẩn:}
\[ \int_{-\infty}^{+\infty} f(x)\,dx = 1, \quad P(a \le X \le b) = \int_a^b f(x)\,dx, \quad \mu = \int_{-\infty}^{+\infty} x f(x)\,dx, \quad f_{\mathcal{N}}(x) = \frac{1}{\sigma\sqrt{2\pi}}e^{-\frac{(x-\mu)^2}{2\sigma^2}} \]
\noindent\textit{Chú giải:} Hàm mật độ mũ (Exponential distribution) mô hình hóa thời gian chờ: f(t) = $\lambda e^{-\lambda t}$ (với t $\ge$  0) có kỳ vọng $\mu$  = 1/$\lambda$ .\\[4pt]
\end{congthuc}

\noindent\textbf{Chú giải \& Lưu ý kỹ thuật:}
\begin{itemize}[leftmargin=1.5em, itemsep=2pt]
  \item Với phân phối chuẩn N($\mu$ , $\sigma ^2$), quy tắc thực nghiệm 68–95–99.7 cho biết khoảng 68.27\% dữ liệu nằm trong [$\mu - \sigma$ , $\mu + \sigma$ ] và 95.45\% nằm trong [$\mu$  - 2$\sigma$ , $\mu$  + 2$\sigma$ ].
\end{itemize}

\begin{vidu}{Kiểm tra hàm mật độ xác suất và tính kỳ vọng tuổi thọ linh kiện}
\textbf{Đề bài:} Tuổi thọ X (tính bằng nghìn giờ) của một vi mạch điện tử có hàm mật độ xác suất f(x) = k $e^{-x/5}$ khi x $\ge$ 0 (và f(x) = 0 khi x < 0). Tìm hằng số k, xác suất mạch hoạt động trên 5 nghìn giờ và tuổi thọ trung bình $\mu$.
\[ f(x) = \begin{cases} k e^{-x/5} & \text{khi } x \ge 0 \\ 0 & \text{khi } x < 0 \end{cases} \]
\textbf{Lời giải chi tiết:}
\begin{enumerate}[leftmargin=1.5em, itemsep=2pt]
  \item Bước 1: Giải điều kiện chuẩn hóa $\int_0^{+\infty}$ k $e^{-x/5}$ dx = 1 $\Rightarrow$  5k = 1 $\Rightarrow$  k = 1/5 = 0.2.
  \item Bước 2: Tính xác suất P(X $\ge$  5) = $\int_5^{+\infty}$ (1/$5)e^{-x/5}$ dx.
  \item Bước 3: Tính kỳ vọng $\mu = \int_0^{+\infty}$ x $\cdot$  (1/$5)e^{-x/5}$ dx bằng tích phân từng phần.
\end{enumerate}
\[ \int_0^{+\infty} k e^{-x/5}\,dx = 5k = 1 \implies k = \frac{1}{5}, \qquad P(X \ge 5) = \int_5^{+\infty} \frac{1}{5}e^{-x/5}\,dx = e^{-1} \approx 36.79\% \]
\[ \mu = \int_0^{+\infty} \frac{x}{5}e^{-x/5}\,dx = 5 \text{ (nghìn giờ)} \]
\noindent\textbf{Kết luận:} Hằng số k = 1/5, xác suất hoạt động trên 5000 giờ là 1/e $\approx$  36.79\% và tuổi thọ trung bình là 5000 giờ.
\end{vidu}

\section{Minh họa Hình học bằng TikZ}
\begin{figure}[htbp]
\centering
\begin{tikzpicture}[scale=1.15, >=stealth]
  % Trục tọa độ
  \draw[->, thick] (-0.3,0) -- (4.8,0) node[right] {$x$};
  \draw[->, thick] (0,-0.3) -- (0,3.5) node[above] {$y$};
  
  % Đường cong y = f(x)
  \draw[domain=0.5:4.2, smooth, variable=\x, blue!80!black, very thick] 
    plot ({\x}, {0.22*(\x)^2 + 0.4}) node[above left] {$y = f(x)$};
    
  % Tam giác vi phân ds, dx, dy
  \draw[dashed, thick, teal!80!black] (1.8, 1.11) -- (3.3, 1.11) node[midway, below] {$dx$};
  \draw[dashed, thick, teal!80!black] (3.3, 1.11) -- (3.3, 2.80) node[midway, right] {$dy = f'(x)dx$};
  \draw[red!80!black, very thick] (1.8, 1.11) -- (3.3, 2.80) node[midway, above left] {$ds = \sqrt{dx^2+dy^2}$};
  
  \filldraw[black] (1.8, 1.11) circle (2pt) node[left] {$P_i$};
  \filldraw[black] (3.3, 2.80) circle (2pt) node[above] {$P_{i+1}$};
\end{tikzpicture}
\caption{Hình 8.1: Vi phân độ dài cung ds = $\sqrt{dx^2 + dy^2}$ và đường cong phân phối chuẩn Gauss.}
\label{fig:tikz-ch8-arclength}
\end{figure}

\section{Ứng dụng Thực tế \& Tối ưu hóa (Kỹ thuật Xây dựng Thủy lợi \& Khoa học Dữ liệu)}
\begin{ungdung}{Thiết kế Đập Thủy điện và Độ tin cậy Hệ thống Server theo Phân phối Xác suất}
Khi thiết kế đập thủy điện, kỹ sư kết cấu phải tính tích phân áp lực thủy tĩnh để xác định bề dày chân đập và vị trí điểm đặt hợp lực (tâm áp lực) nhằm chống lật đập. Trong Khoa học Dữ liệu và AI, tích phân hàm mật độ chuẩn Gauss (Error Function - erf) là nền tảng của kiểm định giả thuyết thống kê và hàm kích hoạt GELU trong mô hình ngôn ngữ lớn.
\[ \text{GELU}(x) = x \cdot \Phi(x) = x \int_{-\infty}^x \frac{1}{\sqrt{2\pi}} e^{-t^2/2}\,dt \]
\noindent\textbf{Ví dụ tính toán thực tế:} Một cửa van đập nước hình chữ nhật rộng w = 4 m, cao 3 m đặt ngập sát mặt nước chịu tổng lực đẩy thủy tĩnh F = $\int_0^3$ 9800 $\cdot$  y $\cdot$  4 dy = 19600 $\times  [y^2]_0^3$ = 176,400 N = 176.4 kN.
\end{ungdung}

\section{Bài tập Trắc nghiệm Củng cố Chương 8}
\begin{enumerate}[label=\textbf{Câu \arabic*:}, leftmargin=2.2em, itemsep=8pt]
  \item Công thức tính diện tích mặt tròn xoay tạo thành khi quay đường cong y = $e^x$ (0 $\le$  x $\le$  1) quanh trục hoành Ox là:
  \begin{itemize}[label={}, leftmargin=1em, itemsep=2pt]
    \item \textbf{A.} S = 2$\pi \int_0^1 e^x \sqrt{1 + e^{2x}}$ dx \quad ($S = 2\pi \int_0^1 e^x \sqrt{1 + e^{2x}}\,dx$)
    \item \textbf{B.} S = $\pi \int_0^1 e^{2x}$ dx \quad ($S = \pi \int_0^1 e^{2x}\,dx$)
    \item \textbf{C.} S = 2$\pi \int_0^1$ x $\sqrt{1 + e^{2x}}$ dx \quad ($S = 2\pi \int_0^1 x \sqrt{1 + e^{2x}}\,dx$)
    \item \textbf{D.} S = $\int_0^1 \sqrt{1 + e^{2x}}$ dx \quad ($S = \int_0^1 \sqrt{1 + e^{2x}}\,dx$)
  \end{itemize}
  \vspace{2pt}
  \noindent{\color{teal!85!black}\textbf{Đáp án đúng: A.}} \textit{Giải thích:} Với y = f(x) = $e^x \Rightarrow$  f'(x) = $e^x$. Diện tích mặt tròn xoay quanh Ox là S = 2$\pi \int_0^1$ f(x)$\sqrt{1 + [f'(x} ]^2$)dx = 2$\pi \int_0^1 e^x \sqrt{1 + e^{2x}}$dx.
  \[ S_{Ox} = 2\pi \int_a^b y\sqrt{1+(y')^2}\,dx = 2\pi \int_0^1 e^x\sqrt{1+e^{2x}}\,dx \]
  \item Hàm số f(x) = c x(2 - x) trên đoạn [0, 2] (và bằng 0 ngoài đoạn đó) là một hàm mật độ xác suất khi hằng số c bằng:
  \begin{itemize}[label={}, leftmargin=1em, itemsep=2pt]
    \item \textbf{A.} c = 3/4 \quad ($c = \frac{3}{4}$)
    \item \textbf{B.} c = 4/3 \quad ($c = \frac{4}{3}$)
    \item \textbf{C.} c = 1/2 \quad ($c = \frac{1}{2}$)
    \item \textbf{D.} c = 3/2 \quad ($c = \frac{3}{2}$)
  \end{itemize}
  \vspace{2pt}
  \noindent{\color{teal!85!black}\textbf{Đáp án đúng: A.}} \textit{Giải thích:} Điều kiện chuẩn hóa của hàm mật độ xác suất: $\int_0^2$ c(2x - $x^2$)dx = 1 $\iff$  c $[x^2 - x^3 / 3]_0^2$ = 1 $\iff$  c(4 - 8/3) = 1 $\iff$  (4/3)c = 1 $\iff$  c = 3/4.
  \[ c \int_0^2 (2x - x^2)\,dx = c\left(4 - \frac{8}{3}\right) = \frac{4}{3}c = 1 \implies c = \frac{3}{4} \]
  \item Độ dài cung của đường dây xích (Catenary) y = $\cos$h x từ x = 0 đến x = $\ln$ 2 bằng:
  \begin{itemize}[label={}, leftmargin=1em, itemsep=2pt]
    \item \textbf{A.} 3/4 \quad ($\frac{3}{4}$)
    \item \textbf{B.} 5/4 \quad ($\frac{5}{4}$)
    \item \textbf{C.} ln 2 \quad ($\ln 2$)
    \item \textbf{D.} 1/2 \quad ($\frac{1}{2}$)
  \end{itemize}
  \vspace{2pt}
  \noindent{\color{teal!85!black}\textbf{Đáp án đúng: A.}} \textit{Giải thích:} Ta có y' = sinh x, suy ra 1 + $(y')^2$ = 1 + $sinh^2$ x = $cosh^2$ x. Do đó L = $\int_0^{ln 2}$ cosh x dx = [sinh $x]_0^{ln 2} = (e^{ln 2} - e^{-ln 2}$)/2 = (2 - 1/2)/2 = 3/4.
  \[ L = \int_0^{\ln 2} \sqrt{1 + \sinh^2 x}\,dx = \Big[\sinh x\Big]_0^{\ln 2} = \frac{2 - 1/2}{2} = \frac{3}{4} \]
\end{enumerate}

% ======================================================================
% CHƯƠNG 9: PHƯƠNG TRÌNH VI PHÂN CẤP 1 (Buổi 10 - 11)
% ======================================================================
\chapter{Phương trình Vi phân Cấp 1}
\label{chap:chuong-9}

\begin{tcolorbox}[colback=slate!5!white,colframe=slate!75!black,title={\textbf{Mục lục Chương 9 \& Chuẩn đầu ra (Buổi 10 - 11)}}]
\textbf{Tổng quan:} Phương trình vi phân (PTVP) là ngôn ngữ toán học mô tả các hệ thống động lực học biến thiên theo thời gian. Chương này trình bày cách mô hình hóa bài toán thực tế bằng PTVP cấp 1, biểu diễn hình học qua Trường hướng (Direction Fields), xấp xỉ nghiệm số bằng Phương pháp Euler, và giải giải tích chính xác cho PTVP tách biến, mô hình tăng trưởng mũ / Logistic cùng PTVP tuyến tính cấp 1.\\[6pt]
\textbf{Nội dung chính trong chương:}
\begin{itemize}[leftmargin=1.5em, itemsep=2pt]
  \item \textbf{Mục 9.1 - 9.2:} Mô hình hóa bằng PTVP, Trường hướng \& Phương pháp số Euler
  \item \textbf{Mục 9.3:} Phương trình vi phân tách biến (Separable Equations) \& Bài toán pha trộn
  \item \textbf{Mục 9.4 - 9.5:} Mô hình tăng trưởng/suy giảm hàm mũ \& Phương trình tăng trưởng Logistic
  \item \textbf{Mục 9.6:} Phương trình vi phân tuyến tính cấp 1 (Phương pháp Thừa số tích phân)
\end{itemize}
\textbf{Chuẩn đầu ra (LLO):}
\begin{itemize}[leftmargin=1.5em, itemsep=2pt]
  \item LLO10.1: Hiểu và vận dụng được các mô hình toán học bằng phương trình vi phân (tăng trưởng, pha trộn, mạch điện).
  \item LLO10.2: Khảo sát dáng điệu của nghiệm qua trường hướng và tính nghiệm xấp xỉ bằng phương pháp Euler.
  \item LLO10.3 \& LLO11.1: Giải thành thạo phương trình vi phân tách biến, phương trình Logistic và phương trình vi phân tuyến tính cấp 1.
\end{itemize}
\end{tcolorbox}

\section{[Mục 9.1 - 9.2] Mô hình hóa bằng PTVP, Trường hướng \& Phương pháp số Euler}
Với bài toán giá trị ban đầu (Initial Value Problem - IVP) y' = F(x, y), $y(x_0$) = $y_0$: tại mỗi điểm (x, y) trên mặt phẳng, ta vẽ một đoạn thẳng ngắn có hệ số góc F(x, y) để tạo thành Trường hướng. Khi không giải được công thức nghiệm giải tích, phương pháp Euler xấp xỉ đường cong nghiệm bằng chuỗi các đoạn tiếp tuyến liên tiếp với bước nhảy h.

\begin{congthuc}{Các công thức trọng tâm (9.1 - 9.2)}
\noindent\textbf{Thuật toán lặp Euler giải gần đúng bài toán Cauchy y' = F(x, y), $y(x_0$) = $y_0$:}
\[ x_{n+1} = x_n + h, \qquad y_{n+1} = y_n + h \cdot F(x_n, y_n), \quad n = 0, 1, 2, \dots \]
\noindent\textit{Chú giải:} Bước nhảy h càng nhỏ thì sai số tích lũy càng giảm nhưng khối lượng tính toán tăng tỉ lệ nghịch với h.\\[4pt]
\end{congthuc}

\noindent\textbf{Chú giải \& Lưu ý kỹ thuật:}
\begin{itemize}[leftmargin=1.5em, itemsep=2pt]
  \item Các nghiệm cân bằng (Equilibrium solutions) của phương trình tự trị y' = g(y) là các hàm hằng y = c thỏa mãn g(c) = 0.
\end{itemize}

\begin{vidu}{Tính nghiệm xấp xỉ bằng Phương pháp Euler}
\textbf{Đề bài:} Cho bài toán giá trị ban đầu y' = x + y, y(0) = 1. Sử dụng phương pháp Euler với bước nhảy h = 0.1 để tính xấp xỉ giá trị y(0.2).
\[ y' = x + y, \quad y(0) = 1, \quad h = 0.1 \implies y(0.2) \approx y_2 \]
\textbf{Lời giải chi tiết:}
\begin{enumerate}[leftmargin=1.5em, itemsep=2pt]
  \item Bước 1: Khởi tạo $x_0$ = 0, $y_0$ = 1 và hàm F(x, y) = x + y.
  \item Bước 2: Tính bước 1 tại $x_1$ = 0.1: $y_1 = y_0$ + h $\cdot F(x_0$, $y_0$).
  \item Bước 3: Tính bước 2 tại $x_2$ = 0.2: $y_2 = y_1$ + h $\cdot F(x_1$, $y_1$).
\end{enumerate}
\[ y_1 = 1 + 0.1 \times (0 + 1) = 1.10 \quad (\text{tại } x_1 = 0.1) \]
\[ y_2 = 1.10 + 0.1 \times (0.1 + 1.10) = 1.10 + 0.12 = 1.22 \quad (\text{tại } x_2 = 0.2) \]
\noindent\textbf{Kết luận:} Giá trị xấp xỉ Euler là y(0.2) $\approx$  1.22 (nghiệm đúng là y(x) = $2e^x$ - x - 1 cho y(0.2) $\approx$  1.2428).
\end{vidu}

\section{[Mục 9.3] Phương trình vi phân tách biến (Separable Equations) \& Bài toán pha trộn}
Phương trình tách biến có dạng dy/dx = g(x)$\cdot$ f(y). Khi f(y) $\neq$  0, ta chuyển tất cả các biểu thức chứa y về vế trái đi kèm dy, tất cả biểu thức chứa x về vế phải đi kèm dx, rồi lấy tích phân hai vế.

\begin{congthuc}{Các công thức trọng tâm (9.3)}
\noindent\textbf{Phương pháp giải PTVP Tách biến:}
\[ \frac{dy}{dx} = g(x)f(y) \implies \int \frac{dy}{f(y)} = \int g(x)\,dx + C \]
\noindent\textit{Chú giải:} Đừng quên kiểm tra riêng các nghiệm hằng y = $y_0$ làm cho $f(y_0$) = 0 trước khi chia hai vế cho f(y).\\[4pt]
\noindent\textbf{Mô hình Bài toán Pha trộn dung dịch (Mixing Problems):}
\[ \frac{dA}{dt} = R_{\text{vào}} - R_{\text{ra}} = (C_{\text{vào}} \cdot v_{\text{vào}}) - \left(\frac{A(t)}{V(t)} \cdot v_{\text{ra}}\right) \]
\noindent\textit{Chú giải:} Trong đó A(t) là lượng chất tan (kg hoặc gam) trong bồn tại thời điểm t, V(t) là thể tích dung dịch.\\[4pt]
\end{congthuc}

\noindent\textbf{Chú giải \& Lưu ý kỹ thuật:}
\begin{itemize}[leftmargin=1.5em, itemsep=2pt]
  \item Định luật làm nguội Newton (Newton's Law of Cooling) dT/dt = k(T - $T_m$) cũng là một phương trình tách biến điển hình, cho nghiệm T(t) = $T_m + (T_0 - T_m)e^{kt}$.
\end{itemize}

\begin{vidu}{Giải bài toán giá trị ban đầu của PTVP tách biến}
\textbf{Đề bài:} Giải bài toán giá trị ban đầu: $\frac{dy}{dx} = \frac{x}{y(1 + x^2)}$ với điều kiện y(0) = -2.
\[ \frac{dy}{dx} = \frac{x}{y(1 + x^2)}, \quad y(0) = -2 \]
\textbf{Lời giải chi tiết:}
\begin{enumerate}[leftmargin=1.5em, itemsep=2pt]
  \item Bước 1: Tách biến về hai vế: y dy = [x / (1 + $x^2$)] dx.
  \item Bước 2: Lấy nguyên hàm hai vế để tìm nghiệm tổng quát dạng ẩn.
  \item Bước 3: Thay x = 0, y = -2 để tìm hằng số C và chọn nhánh dấu âm phù hợp với y(0) = -2 < 0.
\end{enumerate}
\[ \int y\,dy = \int \frac{x}{1 + x^2}\,dx \iff \frac{y^2}{2} = \frac{1}{2}\ln(1 + x^2) + C_1 \iff y^2 = \ln(1 + x^2) + C \]
\[ \text{Thay } x = 0, y = -2 \implies (-2)^2 = \ln 1 + C \implies C = 4 \implies y(x) = -\sqrt{\ln(1 + x^2) + 4} \]
\noindent\textbf{Kết luận:} Nghiệm duy nhất của bài toán là y(x) = -$\sqrt{ln(1 + x^2}$ + 4).
\end{vidu}

\section{[Mục 9.4 - 9.5] Mô hình tăng trưởng/suy giảm hàm mũ \& Phương trình tăng trưởng Logistic}
Khi tài nguyên dồi dào, quần thể tăng trưởng theo hàm mũ dP/dt = kP. Trong thực tế, môi trường có sức chứa tối đa M (Carrying capacity), nên tốc độ tăng trưởng giảm dần khi P tiến gần tới M, được mô tả chính xác bởi phương trình vi phân Logistic của Verhulst.

\begin{congthuc}{Các công thức trọng tâm (9.4 - 9.5)}
\noindent\textbf{Mô hình Tăng trưởng / Phân rã phóng xạ hàm mũ:}
\[ \frac{dP}{dt} = kP \implies P(t) = P_0 e^{kt}, \qquad T_{1/2} = \frac{\ln 2}{|k|} \;(\text{Chu kỳ bán rã khi } k < 0) \]
\noindent\textit{Chú giải:} Tốc độ biến thiên tức thời tỉ lệ thuận với quy mô hiện tại của đại lượng.\\[4pt]
\noindent\textbf{Phương trình vi phân Logistic \& Nghiệm giải tích:}
\[ \frac{dP}{dt} = kP\left(1 - \frac{P}{M}\right) \implies P(t) = \frac{M}{1 + A e^{-kt}}, \quad \text{với } A = \frac{M - P_0}{P_0} \]
\noindent\textit{Chú giải:} Quần thể tăng nhanh nhất (điểm uốn của đường cong chữ S) khi đạt đúng một nửa sức chứa: P = M / 2.\\[4pt]
\end{congthuc}

\noindent\textbf{Chú giải \& Lưu ý kỹ thuật:}
\begin{itemize}[leftmargin=1.5em, itemsep=2pt]
  \item Khi t $\to + \infty$ (với k > 0), số hạng $e^{-kt} \to$  0 nên nghiệm Logistic luôn hội tụ về sức chứa cân bằng: $\lim_{t\to \infty}$ P(t) = M bất kể giá trị ban đầu $P_0$ > 0.
\end{itemize}

\begin{vidu}{Phân tích mô hình tăng trưởng Logistic}
\textbf{Đề bài:} Số lượng cá thể P(t) trong một khu bảo tồn tuân theo phương trình vi phân $\frac{dP}{dt}$ = 0.04 P $\left$(1 - $\frac{P}{5000} \right$) với P(0) = 1000. Hãy tìm công thức P(t), giới hạn của P(t) khi t $\to \infty$ và quy mô quần thể khi tốc độ tăng trưởng đạt cực đại.
\[ \frac{dP}{dt} = 0.04 P \left(1 - \frac{P}{5000}\right), \quad P(0) = 1000 \]
\textbf{Lời giải chi tiết:}
\begin{enumerate}[leftmargin=1.5em, itemsep=2pt]
  \item Bước 1: Xác định các tham số chuẩn: k = 0.04, sức chứa tối đa M = 5000, và $P_0$ = 1000.
  \item Bước 2: Tính hằng số A = (M - $P_0$)/$P_0$ = (5000 - 1000)/1000 = 4.
  \item Bước 3: Viết nghiệm P(t) và xác định điểm uốn P = M/2.
\end{enumerate}
\[ P(t) = \frac{5000}{1 + 4e^{-0.04t}}, \qquad \lim_{t \to +\infty} P(t) = 5000 \]
\[ \text{Tốc độ tăng } \frac{dP}{dt} \text{ đạt cực đại khi } P = \frac{M}{2} = \frac{5000}{2} = 2500 \text{ (cá thể)} \]
\noindent\textbf{Kết luận:} Công thức quần thể là P(t) = 5000 / (1 + $4e^{-0.04t}$), tăng nhanh nhất khi đạt 2500 cá thể và tiến tới ổn định ở 5000 cá thể.
\end{vidu}

\section{[Mục 9.6] Phương trình vi phân tuyến tính cấp 1 (Phương pháp Thừa số tích phân)}
Phương trình vi phân tuyến tính cấp 1 có dạng chuẩn y' + P(x)y = Q(x). Để giải, ta nhân cả hai vế với Thừa số tích phân (Integrating Factor) I(x) = $e^{\int P(x)dx}$ nhằm biến vế trái thành đạo hàm đúng của tích [I(x) $\cdot$  y]'.

\begin{congthuc}{Các công thức trọng tâm (9.6)}
\noindent\textbf{Thừa số tích phân I(x) và Công thức nghiệm tổng quát:}
\[ y' + P(x)y = Q(x) \implies I(x) = e^{\int P(x)\,dx}, \qquad y(x) = \frac{1}{I(x)} \left( \int I(x)Q(x)\,dx + C \right) \]
\noindent\textit{Chú giải:} Lưu ý: Nếu phương trình có dạng a(x)y' + b(x)y = c(x), bắt buộc phải chia hai vế cho a(x) để đưa hệ số của y' về 1 trước khi tìm P(x).\\[4pt]
\end{congthuc}

\noindent\textbf{Chú giải \& Lưu ý kỹ thuật:}
\begin{itemize}[leftmargin=1.5em, itemsep=2pt]
  \item Ứng dụng kinh điển trong kỹ thuật điện: Mạch RL nối tiếp tuân theo định luật Kirchhoff L(dI/dt) + RI = E(t), là một PTVP tuyến tính cấp 1 đối với dòng điện I(t).
\end{itemize}

\begin{vidu}{Giải phương trình vi phân tuyến tính cấp 1 bằng thừa số tích phân}
\textbf{Đề bài:} Tìm nghiệm tổng quát của phương trình vi phân x y' + 2y = $4x^2$ với x > 0.
\[ x y' + 2y = 4x^2 \quad (x > 0) \]
\textbf{Lời giải chi tiết:}
\begin{enumerate}[leftmargin=1.5em, itemsep=2pt]
  \item Bước 1: Chia hai vế cho x > 0 để đưa về dạng chuẩn: y' + (2/x)y = 4x. Xác định P(x) = 2/x, Q(x) = 4x.
  \item Bước 2: Tính thừa số tích phân I(x) = $e^{\int (2/x)dx} = e^{2 ln x} = x^2$.
  \item Bước 3: Nhân hai vế dạng chuẩn với $x^2$, đưa vế trái về $(x^2$ y)' và lấy nguyên hàm.
\end{enumerate}
\[ y' + \frac{2}{x}y = 4x \xrightarrow{\times I(x) = x^2} x^2 y' + 2xy = 4x^3 \iff (x^2 y)' = 4x^3 \]
\[ x^2 y = \int 4x^3\,dx = x^4 + C \implies y(x) = x^2 + \frac{C}{x^2} \]
\noindent\textbf{Kết luận:} Nghiệm tổng quát của phương trình trên khoảng (0, +$\infty$) là y(x) = $x^2$ + C / $x^2$.
\end{vidu}

\section{Minh họa Hình học bằng TikZ}
\begin{figure}[htbp]
\centering
\begin{tikzpicture}[scale=1.05, >=stealth]
  % Trục tọa độ
  \draw[->, thick] (-0.3,0) -- (5.2,0) node[right] {$t$};
  \draw[->, thick] (0,-0.3) -- (0,3.6) node[above] {$P(t)$};
  
  % Đường tiệm cận ngang sức chứa tối đa M = 2.8
  \draw[dashed, red!80!black, thick] (0, 2.8) node[left] {$M$} -- (5.0, 2.8);
  \draw[dotted, teal!80!black] (0, 1.4) node[left] {$M/2$} -- (5.0, 1.4);
  
  % Trường hướng rút gọn (các vector chỉ phương)
  \foreach \x in {0.5, 1.3, 2.1, 2.9, 3.7, 4.5} {
    \foreach \y in {0.4, 1.0, 1.8, 2.4, 3.2} {
      \pgfmathsetmacro{\slope}{0.65*\y*(1 - \y/2.8)}
      \draw[->, gray!60, thin] (\x-0.12, \y-0.12*\slope) -- (\x+0.12, \y+0.12*\slope);
    }
  }
  
  % Đường cong nghiệm Logistic chữ S
  \draw[domain=0:4.9, smooth, variable=\t, blue!85!black, very thick] 
    plot ({\t}, {2.8 / (1 + 6*exp(-1.1*\t))}) node[below right] {$P(t) = \frac{M}{1+Ae^{-kt}}$};
    
  \filldraw[teal!80!black] (1.63, 1.4) circle (2.2pt) node[below right] {Điểm uốn};
\end{tikzpicture}
\caption{Hình 9.1: Trường hướng (Direction Field) và đường cong nghiệm chữ S của Phương trình vi phân Logistic tiến về sức chứa M.}
\label{fig:tikz-ch9-logistic}
\end{figure}

\section{Ứng dụng Thực tế \& Tối ưu hóa (Kỹ thuật Mạch điện (Định luật Kirchhoff) \& Dịch tễ học)}
\begin{ungdung}{Quá trình Quá độ trong Mạch điện RL / RC và Mô hình Lan truyền Thông tin}
Khi đóng khóa K trong mạch điện gồm cuộn cảm L và điện trở R mắc vào nguồn một chiều E, dòng điện i(t) không tăng tức thời mà tuân theo PTVP tuyến tính L(di/dt) + Ri = E. Nghiệm giải tích cho thấy dòng điện tiến dần theo hàm mũ về giá trị xác lập E/R với hằng số thời gian $\tau$  = L/R.
\[ L\frac{di}{dt} + R i = E, \; i(0) = 0 \implies i(t) = \frac{E}{R}\left(1 - e^{-\frac{R}{L}t}\right) \]
\noindent\textbf{Ví dụ tính toán thực tế:} Với cuộn dây có L = 2 H, điện trở R = 10 $\Omega$  và nguồn E = 20 V: Dòng điện xác lập là $I_max$ = 20/10 = 2 A, và sau thời gian t = $\tau$  = L/R = 0.2 giây, dòng điện đạt 2(1 - $e^{-1}$) $\approx$  1.264 A (63.2\% giá trị cực đại).
\end{ungdung}

\section{Bài tập Trắc nghiệm Củng cố Chương 9}
\begin{enumerate}[label=\textbf{Câu \arabic*:}, leftmargin=2.2em, itemsep=8pt]
  \item Thừa số tích phân I(x) của phương trình vi phân tuyến tính y' - $3x^2$ y = $e^{x^3}$ là:
  \begin{itemize}[label={}, leftmargin=1em, itemsep=2pt]
    \item \textbf{A.} I(x) = $e^{-x^3}$ \quad ($I(x) = e^{-x^3}$)
    \item \textbf{B.} I(x) = $e^{x^3}$ \quad ($I(x) = e^{x^3}$)
    \item \textbf{C.} I(x) = $-x^3$ \quad ($I(x) = -x^3$)
    \item \textbf{D.} I(x) = $e^{-3x^2}$ \quad ($I(x) = e^{-3x^2}$)
  \end{itemize}
  \vspace{2pt}
  \noindent{\color{teal!85!black}\textbf{Đáp án đúng: A.}} \textit{Giải thích:} Phương trình có dạng chuẩn y' + P(x)y = Q(x) với P(x) = $-3x^2$. Vậy thừa số tích phân là I(x) = $e^{\int (-3x^2)dx} = e^{-x^3}$.
  \[ I(x) = \exp\left(\int -3x^2\,dx\right) = e^{-x^3} \]
  \item Nghiệm của bài toán giá trị ban đầu y' = 2xy, y(0) = 3 là:
  \begin{itemize}[label={}, leftmargin=1em, itemsep=2pt]
    \item \textbf{A.} y(x) = 3 $e^{x^2}$ \quad ($y(x) = 3e^{x^2}$)
    \item \textbf{B.} y(x) = $e^{x^2}$ + 2 \quad ($y(x) = e^{x^2} + 2$)
    \item \textbf{C.} y(x) = 3 $e^{2x}$ \quad ($y(x) = 3e^{2x}$)
    \item \textbf{D.} y(x) = 3 $x^2$ \quad ($y(x) = 3x^2$)
  \end{itemize}
  \vspace{2pt}
  \noindent{\color{teal!85!black}\textbf{Đáp án đúng: A.}} \textit{Giải thích:} Tách biến: dy/y = 2x dx $\Rightarrow$  ln|y| = $x^2 + C_1 \Rightarrow$  y = C $e^{x^2}$. Thay y(0) = 3 $\Rightarrow$  C = 3. Vậy y(x) = $3e^{x^2}$.
  \[ \frac{dy}{y} = 2x\,dx \implies y = Ce^{x^2} \xrightarrow{y(0)=3} y = 3e^{x^2} \]
  \item Cho phương trình Logistic dP/dt = 0.2 P (1 - P/1000). Các nghiệm cân bằng (hàm hằng) của phương trình là:
  \begin{itemize}[label={}, leftmargin=1em, itemsep=2pt]
    \item \textbf{A.} P = 0 và P = 1000 \quad ($P = 0 \text{ và } P = 1000$)
    \item \textbf{B.} P = 500 và P = 1000 \quad ($P = 500 \text{ và } P = 1000$)
    \item \textbf{C.} Chỉ có P = 1000 \quad ($P = 1000$)
    \item \textbf{D.} P = 0 và P = 200 \quad ($P = 0 \text{ và } P = 200$)
  \end{itemize}
  \vspace{2pt}
  \noindent{\color{teal!85!black}\textbf{Đáp án đúng: A.}} \textit{Giải thích:} Nghiệm cân bằng là các giá trị P làm cho dP/dt = 0 $\iff$  0.2 P (1 - P/1000) = 0 $\iff$  P = 0 hoặc P = 1000.
  \[ \frac{dP}{dt} = 0 \iff 0.2P\left(1 - \frac{P}{1000}\right) = 0 \iff P = 0 \text{ hoặc } P = 1000 \]
\end{enumerate}

% ======================================================================
% CHƯƠNG 10: PHƯƠNG TRÌNH THAM SỐ VÀ TỌA ĐỘ CỰC (Buổi 12 - 13)
% ======================================================================
\chapter{Phương trình Tham số và Tọa độ Cực}
\label{chap:chuong-10}

\begin{tcolorbox}[colback=slate!5!white,colframe=slate!75!black,title={\textbf{Mục lục Chương 10 \& Chuẩn đầu ra (Buổi 12 - 13)}}]
\textbf{Tổng quan:} Nhiều quỹ đạo chuyển động và đường cong hình học phức tạp (như đường Cycloid, đường hình tim Cardioid, đường hoa hồng) không thể viết dưới dạng hàm đơn trị y = f(x). Chương này trang bị phương pháp biểu diễn đường cong bằng Phương trình tham số (x(t), y(t)) và Hệ tọa độ cực (r, $\theta$ ), cùng các công thức giải tích tính tiếp tuyến, diện tích, độ dài cung và phương trình các đường Conic.\\[6pt]
\textbf{Nội dung chính trong chương:}
\begin{itemize}[leftmargin=1.5em, itemsep=2pt]
  \item \textbf{Mục 10.1 - 10.2:} Đường cong tham số \& Giải tích trên đường cong tham số (Tiếp tuyến, Độ dài cung, Diện tích)
  \item \textbf{Mục 10.3 - 10.4:} Hệ tọa độ cực (r, $\theta$ ) \& Tính diện tích, độ dài cung trong tọa độ cực
  \item \textbf{Mục 10.5 - 10.6:} Các đường Conic (Parabol, Elip, Hyperbol) \& Phương trình Conic trong tọa độ cực
\end{itemize}
\textbf{Chuẩn đầu ra (LLO):}
\begin{itemize}[leftmargin=1.5em, itemsep=2pt]
  \item LLO12.1: Vận dụng được phương trình tham số của đường cong để giải các bài toán tiếp tuyến, độ dài cung và diện tích.
  \item LLO13.1: Vận dụng thành thạo hệ tọa độ cực để tính diện tích, độ dài cung và thiết lập phương trình các đường conic (Elip, Parabol, Hyperbol).
\end{itemize}
\end{tcolorbox}

\section{[Mục 10.1 - 10.2] Đường cong tham số \& Giải tích trên đường cong tham số (Tiếp tuyến, Độ dài cung, Diện tích)}
Khi tọa độ điểm M(x, y) đều là hàm của biến thứ ba t (gọi là tham số, thường là thời gian): x = f(t), y = g(t), ta có thể tính đạo hàm dy/dx, đạo hàm cấp hai $d^2y / dx^2$, diện tích và độ dài cung trực tiếp theo tham số t mà không cần khử t.

\begin{congthuc}{Các công thức trọng tâm (10.1 - 10.2)}
\noindent\textbf{Đạo hàm cấp 1 và cấp 2 của hàm cho bởi phương trình tham số:}
\[ \frac{dy}{dx} = \frac{dy/dt}{dx/dt} = \frac{y'(t)}{x'(t)}, \qquad \frac{d^2y}{dx^2} = \frac{\frac{d}{dt}\left(\frac{dy}{dx}\right)}{dx/dt} = \frac{x'(t)y''(t) - y'(t)x''(t)}{[x'(t)]^3} \]
\noindent\textit{Chú giải:} Tiếp tuyến nằm ngang khi y'(t) = 0 và x'(t) $\neq$  0; tiếp tuyến thẳng đứng khi x'(t) = 0 và y'(t) $\neq$  0.\\[4pt]
\noindent\textbf{Diện tích và Độ dài cung của đường cong tham số ($\alpha \le$  t $\le \beta$ ):}
\[ S = \int_\alpha^\beta |y(t) \cdot x'(t)|\,dt, \qquad L = \int_\alpha^\beta \sqrt{[x'(t)]^2 + [y'(t)]^2}\,dt \]
\noindent\textit{Chú giải:} Vi phân cung tổng quát: ds = $\sqrt{dx^2 + dy^2} = \sqrt{[x'(t} ]^2 + [y'(t)]^2$) dt.\\[4pt]
\end{congthuc}

\noindent\textbf{Chú giải \& Lưu ý kỹ thuật:}
\begin{itemize}[leftmargin=1.5em, itemsep=2pt]
  \item Đường Cycloid (quỹ đạo 1 điểm trên vành bánh xe bán kính R lăn không trượt trên đường thẳng) có phương trình tham số: x = R(t - sin t), y = R(1 - cos t). Độ dài 1 nhịp (0 $\le$  t $\le$  2$\pi$ ) bằng đúng 8R.
\end{itemize}

\begin{vidu}{Tính tiếp tuyến và độ dài cung của đường Cycloid / Đường tham số}
\textbf{Đề bài:} Cho đường cong tham số x(t) = $t^2$, y(t) = $t^3$ - 3t. (a) Tìm các điểm trên đường cong mà tiếp tuyến nằm ngang. (b) Tính đạo hàm cấp hai $d^2y / dx^2$ tại t = 2.
\[ x(t) = t^2, \quad y(t) = t^3 - 3t \]
\textbf{Lời giải chi tiết:}
\begin{enumerate}[leftmargin=1.5em, itemsep=2pt]
  \item Bước 1: Tính x'(t) = 2t và y'(t) = $3t^2$ - 3. Suy ra dy/dx = $(3t^2$ - 3)/(2t).
  \item Bước 2: Tiếp tuyến nằm ngang khi y'(t) = 0 (và x'(t) $\neq$  0) $\iff 3t^2$ - 3 = 0 $\iff$  t = $\pm$ 1.
  \item Bước 3: Đạo hàm biểu thức dy/dx = (3/2)t - (3/$2)t^{-1}$ theo t rồi chia cho x'(t) = 2t.
\end{enumerate}
\[ y'(t) = 0 \iff t = \pm 1 \implies \text{Hai điểm tiếp tuyến ngang là } (1, -2) \text{ và } (1, 2) \]
\[ \frac{d^2y}{dx^2} = \frac{\frac{3}{2} + \frac{3}{2t^2}}{2t} = \frac{3(t^2 + 1)}{4t^3} \implies \left.\frac{d^2y}{dx^2}\right|_{t=2} = \frac{3(4 + 1)}{4(8)} = \frac{15}{32} \]
\noindent\textbf{Kết luận:} Hai điểm có tiếp tuyến ngang là (1, -2), (1, 2) và đạo hàm cấp hai tại t = 2 bằng 15/32.
\end{vidu}

\section{[Mục 10.3 - 10.4] Hệ tọa độ cực (r, $\theta$ ) \& Tính diện tích, độ dài cung trong tọa độ cực}
Trong hệ tọa độ cực, mỗi điểm P được xác định bởi khoảng cách có hướng r từ gốc cực O và góc định hướng $\theta$  hợp với trục cực Ox. Để tính diện tích miền phẳng giới hạn bởi đường cong cực r = r($\theta$ ), ta chia miền thành các hình quạt tròn vi phân có diện tích dA = (1/2) $r^2$ d$\theta$ .

\begin{congthuc}{Các công thức trọng tâm (10.3 - 10.4)}
\noindent\textbf{Công thức chuyển đổi giữa Tọa độ Descartes (x, y) và Tọa độ cực (r, $\theta$ ):}
\[ x = r\cos\theta, \quad y = r\sin\theta, \qquad r^2 = x^2 + y^2, \quad \tan\theta = \frac{y}{x} \]
\noindent\textit{Chú giải:} Đường tròn đi qua gốc cực: r = 2a cos $\theta$  (tâm (a, 0)) và r = 2a sin $\theta$  (tâm (0, a)).\\[4pt]
\noindent\textbf{Diện tích hình quạt cực và Độ dài cung của đường cong r = r($\theta$ ):}
\[ S = \frac{1}{2} \int_\alpha^\beta [r(\theta)]^2\,d\theta, \qquad L = \int_\alpha^\beta \sqrt{[r(\theta)]^2 + [r'(\theta)]^2}\,d\theta \]
\noindent\textit{Chú giải:} Diện tích miền nằm giữa hai đường cong cực $r_ngo$ài($\theta$ ) và $r_trong$($\theta$ ) là (1/2) $\int ([r_ngo$à$i]^2 - [r_trong]^2$) d$\theta$ .\\[4pt]
\end{congthuc}

\noindent\textbf{Chú giải \& Lưu ý kỹ thuật:}
\begin{itemize}[leftmargin=1.5em, itemsep=2pt]
  \item Đường hoa hồng r = a cos(n$\theta$ ) hoặc r = a sin(n$\theta$ ): Nếu n lẻ thì đường cong có đúng n cánh hoa; nếu n chẵn thì đường cong có 2n cánh hoa.
\end{itemize}

\begin{vidu}{Tính diện tích miền giới hạn bởi đường hình tim (Cardioid)}
\textbf{Đề bài:} Tính diện tích miền phẳng nằm bên trong đường hình tim (Cardioid) r = 1 + $\cos \theta$ (với 0 $\le \theta \le$ 2$\pi$).
\[ r = 1 + \cos\theta \implies S = \frac{1}{2}\int_0^{2\pi} (1 + \cos\theta)^2\,d\theta \]
\textbf{Lời giải chi tiết:}
\begin{enumerate}[leftmargin=1.5em, itemsep=2pt]
  \item Bước 1: Áp dụng công thức diện tích trong tọa độ cực S = (1/2) $\int_0^{2\pi }$ (1 + cos $\theta )^2$ d$\theta$ .
  \item Bước 2: Khai triển hằng đẳng thức và hạ bậc $cos^2 \theta$  = (1 + cos 2$\theta$ )/2.
  \item Bước 3: Lấy tích phân trên chu kỳ [0, 2$\pi$ ] (các tích phân của cos $\theta$  và cos 2$\theta$  trên cả chu kỳ đều bằng 0).
\end{enumerate}
\[ S = \frac{1}{2} \int_0^{2\pi} \left(1 + 2\cos\theta + \frac{1 + \cos 2\theta}{2}\right) d\theta = \frac{1}{2} \int_0^{2\pi} \left(\frac{3}{2} + 2\cos\theta + \frac{1}{2}\cos 2\theta\right) d\theta = \frac{1}{2}\left(\frac{3}{2}\cdot 2\pi\right) = \frac{3\pi}{2} \]
\noindent\textbf{Kết luận:} Diện tích toàn bộ đường hình tim Cardioid là S = 3$\pi$ /2.
\end{vidu}

\section{[Mục 10.5 - 10.6] Các đường Conic (Parabol, Elip, Hyperbol) \& Phương trình Conic trong tọa độ cực}
Ba đường Conic (Parabol, Elip, Hyperbol) có thể được định nghĩa thống nhất thông qua tâm sai e (Eccentricity): Tập hợp các điểm M sao cho tỉ số khoảng cách từ M tới tiêu điểm F và khoảng cách từ M tới đường chuẩn l là hằng số e = |MF| / d(M, l).

\begin{congthuc}{Các công thức trọng tâm (10.5 - 10.6)}
\noindent\textbf{Phương trình chính tắc của Elip và Hyperbol trong tọa độ Descartes:}
\[ \frac{x^2}{a^2} + \frac{y^2}{b^2} = 1 \;\left(c^2 = a^2 - b^2, \; e = \frac{c}{a} < 1\right), \qquad \frac{x^2}{a^2} - \frac{y^2}{b^2} = 1 \;\left(c^2 = a^2 + b^2, \; e = \frac{c}{a} > 1\right) \]
\noindent\textit{Chú giải:} Với Parabol $y^2$ = 4px (tiêu điểm F(p, 0), đường chuẩn x = -p), tâm sai luôn bằng e = 1.\\[4pt]
\noindent\textbf{Phương trình thống nhất của đường Conic trong Tọa độ cực (Tiêu điểm tại gốc cực O):}
\[ r = \frac{ed}{1 \pm e\cos\theta} \quad \text{hoặc} \quad r = \frac{ed}{1 \pm e\sin\theta} \]
\noindent\textit{Chú giải:} Phân loại nhanh theo tâm sai e: Nếu e < 1 là Elip; nếu e = 1 là Parabol; nếu e > 1 là Hyperbol.\\[4pt]
\end{congthuc}

\noindent\textbf{Chú giải \& Lưu ý kỹ thuật:}
\begin{itemize}[leftmargin=1.5em, itemsep=2pt]
  \item Với quỹ đạo hành tinh hình Elip nhận Mặt Trời làm tiêu điểm tại gốc cực O: Khoảng cách cận điểm (Perihelion - gần nhất) là $r_min$ = a(1 - e) và viễn điểm (Aphelion - xa nhất) là $r_max$ = a(1 + e).
\end{itemize}

\begin{vidu}{Nhận diện đường Conic và tìm các yếu tố hình học trong tọa độ cực}
\textbf{Đề bài:} Cho đường conic có phương trình tọa độ cực r = $\frac{10}{3 - 2\cos\theta}$. Hãy xác định tâm sai e, tên gọi đường conic và độ dài trục lớn 2a.
\[ r = \frac{10}{3 - 2\cos\theta} \]
\textbf{Lời giải chi tiết:}
\begin{enumerate}[leftmargin=1.5em, itemsep=2pt]
  \item Bước 1: Chia cả tử và mẫu cho 3 để đưa hằng số đầu tiên ở mẫu về 1 chuẩn tắc.
  \item Bước 2: Đọc hệ số của cos $\theta$  để suy ra tâm sai e = 2/3 < 1 (Elip).
  \item Bước 3: Tính $r_max$ (khi $\theta$  = 0) và $r_min$ (khi $\theta = \pi$ ), từ đó suy ra trục lớn 2a = $r_max + r_min$.
\end{enumerate}
\[ r = \frac{10/3}{1 - \frac{2}{3}\cos\theta} \implies e = \frac{2}{3} < 1 \implies \text{Đường cong là một Elip} \]
\[ r_{\max} = r(0) = \frac{10}{3 - 2} = 10, \quad r_{\min} = r(\pi) = \frac{10}{3 + 2} = 2 \implies 2a = r_{\max} + r_{\min} = 12 \implies a = 6 \]
\noindent\textbf{Kết luận:} Đường đã cho là Elip có tâm sai e = 2/3 và bán trục lớn a = 6.
\end{vidu}

\section{Minh họa Hình học bằng TikZ}
\begin{figure}[htbp]
\centering
\begin{tikzpicture}[scale=1.15, >=stealth]
  % Lưới vòng tròn đồng tâm tọa độ cực
  \foreach \r in {0.8, 1.6, 2.4} {
    \draw[very thin, gray!30] (0,0) circle (\r);
  }
  \draw[->, thick] (-2.6,0) -- (2.8,0) node[right] {Trục cực};
  \draw[->, thick] (0,-2.4) -- (0,2.4) node[above] {$\theta = \pi/2$};
  
  % Đường hình tim Cardioid r = 1.1(1 + cos theta)
  \draw[domain=0:360, samples=120, smooth, variable=\t, blue!85!black, very thick] 
    plot ({1.1*(1 + cos(\t))*cos(\t)}, {1.1*(1 + cos(\t))*sin(\t)});
    
  % Đường hoa hồng r = 1.9 cos(2 theta)
  \draw[domain=0:360, samples=160, smooth, variable=\t, red!80!black, thick, dashed] 
    plot ({1.9*cos(2*\t)*cos(\t)}, {1.9*cos(2*\t)*sin(\t)});
    
  \node[blue!85!black, above right] at (1.3, 1.35) {$r = 1 + \cos\theta$};
  \node[red!80!black, below right] at (1.2, -1.45) {$r = a\cos(2\theta)$};
\end{tikzpicture}
\caption{Hình 10.1: Đường hình tim Cardioid r = 1 + cos $\theta$  và đường hoa hồng 4 cánh r = 2 cos(2$\theta$ ) trong hệ tọa độ cực.}
\label{fig:tikz-ch10-polar}
\end{figure}

\section{Ứng dụng Thực tế \& Tối ưu hóa (Cơ học Vũ trụ (Quỹ đạo Vệ tinh) \& Kỹ thuật Âm thanh (Microphone))}
\begin{ungdung}{Định luật Kepler về Quỹ đạo Vệ tinh và Giản đồ Thu âm Cardioid}
Theo Định luật 1 Kepler, mọi vệ tinh nhân tạo và hành tinh đều chuyển động theo quỹ đạo đường Conic trong tọa độ cực r($\theta$ ) = ed / (1 + e cos $\theta$ ) với tâm Trái Đất tại tiêu điểm gốc O. Trong kỹ thuật âm thanh studio, microphone hướng thu Cardioid có độ nhạy theo góc đúng bằng r($\theta$ ) = a(1 + cos $\theta$ ), giúp thu rõ giọng hát phía trước ($\theta$  = 0) và triệt tiêu hoàn toàn tạp âm từ phía sau ($\theta = \pi$ ).
\[ r(\theta) = \frac{a(1 - e^2)}{1 + e\cos\theta}, \qquad G_{\text{mic}}(\theta) = G_0 (1 + \cos\theta) \]
\noindent\textbf{Ví dụ tính toán thực tế:} Vệ tinh có bán trục lớn quỹ đạo a = 10,000 km và tâm sai e = 0.2 sẽ có khoảng cách gần tâm Trái Đất nhất là $r_min$ = 10,000(1 - 0.2) = 8,000 km và xa nhất là $r_max$ = 12,000 km.
\end{ungdung}

\section{Bài tập Trắc nghiệm Củng cố Chương 10}
\begin{enumerate}[label=\textbf{Câu \arabic*:}, leftmargin=2.2em, itemsep=8pt]
  \item Độ dài cung của đường cong tham số x(t) = 3$\cos$ t, y(t) = 3$\sin$ t với 0 $\le$ t $\le \pi$ bằng:
  \begin{itemize}[label={}, leftmargin=1em, itemsep=2pt]
    \item \textbf{A.} 3$\pi$  \quad ($3\pi$)
    \item \textbf{B.} 6$\pi$  \quad ($6\pi$)
    \item \textbf{C.} 9$\pi$  \quad ($9\pi$)
    \item \textbf{D.} 3 \quad ($3$)
  \end{itemize}
  \vspace{2pt}
  \noindent{\color{teal!85!black}\textbf{Đáp án đúng: A.}} \textit{Giải thích:} Ta có $[x'(t)]^2 + [y'(t)]^2$ = (-3sin $t)^2$ + (3cos $t)^2 = 9(sin^2$ t + $cos^2$ t) = 9. Do đó L = $\int_0^\pi \sqrt{9}$ dt = 3$\pi$  (nửa chu vi đường tròn bán kính R = 3).
  \[ L = \int_0^\pi \sqrt{9\sin^2 t + 9\cos^2 t}\,dt = \int_0^\pi 3\,dt = 3\pi \]
  \item Trong hệ tọa độ cực, phương trình r = 4$\sin \theta$ biểu diễn đường cong nào trong hệ tọa độ Descartes?
  \begin{itemize}[label={}, leftmargin=1em, itemsep=2pt]
    \item \textbf{A.} Đường tròn $x^2$ + (y - $2)^2$ = 4 \quad ($x^2 + (y - 2)^2 = 4$)
    \item \textbf{B.} Đường tròn (x - $2)^2 + y^2$ = 4 \quad ($(x - 2)^2 + y^2 = 4$)
    \item \textbf{C.} Đường tròn $x^2 + y^2$ = 16 \quad ($x^2 + y^2 = 16$)
    \item \textbf{D.} Đường thẳng y = 4 \quad ($y = 4$)
  \end{itemize}
  \vspace{2pt}
  \noindent{\color{teal!85!black}\textbf{Đáp án đúng: A.}} \textit{Giải thích:} Nhân cả hai vế với r: $r^2$ = 4 r sin $\theta \iff x^2 + y^2$ = 4y $\iff x^2$ + (y - $2)^2 = 2^2$ (đường tròn tâm (0, 2) bán kính R = 2).
  \[ r^2 = 4r\sin\theta \iff x^2 + y^2 = 4y \iff x^2 + (y - 2)^2 = 4 \]
  \item Đường conic có phương trình tọa độ cực r = $\frac{6}{2 + 3\sin\theta}$ là đường gì?
  \begin{itemize}[label={}, leftmargin=1em, itemsep=2pt]
    \item \textbf{A.} Hyperbol (vì tâm sai e = 3/2 > 1) \quad ($\text{Hyperbol } (e = 1.5 > 1)$)
    \item \textbf{B.} Elip (vì tâm sai e = 2/3 < 1) \quad ($\text{Elip } (e = 2/3 < 1)$)
    \item \textbf{C.} Parabol (vì tâm sai e = 1) \quad ($\text{Parabol } (e = 1)$)
    \item \textbf{D.} Đường tròn \quad ($\text{Đường tròn}$)
  \end{itemize}
  \vspace{2pt}
  \noindent{\color{teal!85!black}\textbf{Đáp án đúng: A.}} \textit{Giải thích:} Chia cả tử và mẫu cho 2: r = 3 / [1 + (3/2)sin $\theta$ ]. Suy ra tâm sai e = 3/2 = 1.5 > 1, vậy đây là một nhánh của Hyperbol.
  \[ r = \frac{3}{1 + \frac{3}{2}\sin\theta} \implies e = \frac{3}{2} > 1 \implies \text{Hyperbol} \]
\end{enumerate}

% ======================================================================
% CHƯƠNG 11: DÃY SỐ VÀ CHUỖI SỐ (CHUỖI LŨY THỪA & TAYLOR) (Buổi 14 - 15)
% ======================================================================
\chapter{Dãy số và Chuỗi số (Chuỗi Lũy thừa \& Taylor)}
\label{chap:chuong-11}

\begin{tcolorbox}[colback=slate!5!white,colframe=slate!75!black,title={\textbf{Mục lục Chương 11 \& Chuẩn đầu ra (Buổi 14 - 15)}}]
\textbf{Tổng quan:} Chương cuối cùng của Giải tích 1 nghiên cứu sự hội tụ của dãy số vô hạn, chuỗi số vô hạn cùng bộ tiêu chuẩn kiểm tra sự hội tụ kinh điển (Tích phân, So sánh, Leibniz cho chuỗi đan dấu, D'Alembert và Cauchy). Đỉnh cao của chương là biểu diễn các hàm giải tích bằng Chuỗi lũy thừa, Chuỗi Taylor và Chuỗi Maclaurin để tính toán xấp xỉ độ chính xác tùy ý.\\[6pt]
\textbf{Nội dung chính trong chương:}
\begin{itemize}[leftmargin=1.5em, itemsep=2pt]
  \item \textbf{Mục 11.1 - 11.2:} Dãy số, Chuỗi số, Chuỗi cấp số nhân \& Điều kiện cần của sự hội tụ
  \item \textbf{Mục 11.3 - 11.4:} Chuỗi số dương: Tiêu chuẩn Tích phân \& Tiêu chuẩn So sánh
  \item \textbf{Mục 11.5 - 11.7:} Chuỗi đan dấu (Leibniz), Hội tụ tuyệt đối, Tiêu chuẩn D'Alembert \& Cauchy
  \item \textbf{Mục 11.8 - 11.11:} Chuỗi lũy thừa, Bán kính hội tụ \& Khai triển chuỗi Taylor - Maclaurin
\end{itemize}
\textbf{Chuẩn đầu ra (LLO):}
\begin{itemize}[leftmargin=1.5em, itemsep=2pt]
  \item LLO14.1: Giải được các bài tập tìm giới hạn của dãy số bằng định lý đơn điệu bị chặn và định lý kẹp.
  \item LLO14.2: Vận dụng thành thạo các tiêu chuẩn tích phân, so sánh, Leibniz, D'Alembert và Cauchy để xét sự hội tụ/phân kỳ và hội tụ tuyệt đối của chuỗi số.
  \item LLO15.1: Tìm được miền hội tụ của chuỗi lũy thừa và khai triển thành thạo chuỗi Taylor, Maclaurin của hàm số.
\end{itemize}
\end{tcolorbox}

\section{[Mục 11.1 - 11.2] Dãy số, Chuỗi số, Chuỗi cấp số nhân \& Điều kiện cần của sự hội tụ}
Dãy số {$a_n$} hội tụ nếu tồn tại giới hạn hữu hạn $\lim_{n\to \infty} a_n$ = L. Mọi dãy đơn điệu và bị chặn đều hội tụ (Định lý Weierstrass). Với chuỗi số $\sum_{n=1}^\infty a_n$, nếu dãy tổng riêng $S_n = a_1 + a_2$ + ... + $a_n$ có giới hạn hữu hạn S khi n $\to \infty$ thì chuỗi hội tụ và có tổng bằng S.

\begin{congthuc}{Các công thức trọng tâm (11.1 - 11.2)}
\noindent\textbf{Chuỗi Cấp số nhân (Geometric Series) \& Điều kiện cần để chuỗi hội tụ:}
\[ \sum_{n=0}^\infty a r^n = \begin{cases} \frac{a}{1 - r} & \text{khi } |r| < 1 \text{ (Hội tụ)} \\ \text{Phân kỳ} & \text{khi } |r| \ge 1 \end{cases}, \qquad \sum_{n=1}^\infty a_n \text{ hội tụ} \implies \lim_{n \to \infty} a_n = 0 \]
\noindent\textit{Chú giải:} Tiêu chuẩn phân kỳ (Test for Divergence): Nếu $\lim_{n\to \infty} a_n \neq$  0 (hoặc không tồn tại) thì chuỗi $\sum a_n$ chắc chắn PHÂN KỲ.\\[4pt]
\end{congthuc}

\noindent\textbf{Chú giải \& Lưu ý kỹ thuật:}
\begin{itemize}[leftmargin=1.5em, itemsep=2pt]
  \item Cảnh báo quan trọng: Điều kiện lim $a_n$ = 0 chỉ là điều kiện CẦN, không phải điều kiện đủ! Chuỗi điều hòa (Harmonic series) $\sum_{n=1}^\infty$ (1/n) có lim (1/n) = 0 nhưng PHÂN KỲ.
\end{itemize}

\begin{vidu}{Tính tổng chuỗi số bằng phân tích vi phân (Chuỗi rút gọn - Telescoping Series)}
\textbf{Đề bài:} Chứng minh chuỗi số $\sum_{n=1}^\infty \frac{1}{n(n + 1)}$ hội tụ và tính tổng chính xác của nó.
\[ S = \sum_{n=1}^\infty \frac{1}{n(n + 1)} \]
\textbf{Lời giải chi tiết:}
\begin{enumerate}[leftmargin=1.5em, itemsep=2pt]
  \item Bước 1: Phân tích số hạng tổng quát $a_n$ = 1/[n(n + 1)] = 1/n - 1/(n + 1).
  \item Bước 2: Lập tổng riêng thứ N và khử liên tiếp các số hạng đối nhau.
  \item Bước 3: Lấy giới hạn $\lim_{N\to \infty} S_N$.
\end{enumerate}
\[ S_N = \sum_{n=1}^N \left(\frac{1}{n} - \frac{1}{n + 1}\right) = \left(1 - \frac{1}{2}\right) + \left(\frac{1}{2} - \frac{1}{3}\right) + \cdots + \left(\frac{1}{N} - \frac{1}{N + 1}\right) = 1 - \frac{1}{N + 1} \]
\[ S = \lim_{N \to \infty} S_N = \lim_{N \to \infty} \left(1 - \frac{1}{N + 1}\right) = 1 \]
\noindent\textbf{Kết luận:} Chuỗi hội tụ và có tổng bằng S = 1.
\end{vidu}

\section{[Mục 11.3 - 11.4] Chuỗi số dương: Tiêu chuẩn Tích phân \& Tiêu chuẩn So sánh}
Với chuỗi số dương $(a_n$ > 0), dãy tổng riêng $S_n$ luôn tăng nên chuỗi hội tụ khi và chỉ khi $S_n$ bị chặn trên. Ta kiểm tra bằng cách so sánh với tích phân suy rộng (Tiêu chuẩn Tích phân Cauchy–Maclaurin) hoặc so sánh với chuỗi chuẩn p-series $\sum$ 1/$n^p$.

\begin{congthuc}{Các công thức trọng tâm (11.3 - 11.4)}
\noindent\textbf{Tiêu chuẩn Tích phân \& Chuẩn hội tụ của p-Chuỗi (Riemann Series):}
\[ \sum_{n=1}^\infty a_n \text{ cùng tính chất hội tụ với } \int_1^\infty f(x)\,dx, \qquad \sum_{n=1}^\infty \frac{1}{n^p} \;\begin{cases} \text{Hội tụ} & \text{khi } p > 1 \\ \text{Phân kỳ} & \text{khi } p \le 1 \end{cases} \]
\noindent\textit{Chú giải:} Trong tiêu chuẩn tích phân, hàm f(x) (với f(n) = $a_n$) phải liên tục, dương và nghịch biến trên [1, +$\infty$).\\[4pt]
\noindent\textbf{Tiêu chuẩn So sánh Giới hạn (Limit Comparison Test):}
\[ a_n, b_n > 0, \quad \lim_{n \to \infty} \frac{a_n}{b_n} = c \in (0, +\infty) \implies \sum a_n \text{ và } \sum b_n \text{ cùng hội tụ hoặc cùng phân kỳ} \]
\noindent\textit{Chú giải:} Rất hiệu quả cho các chuỗi phân thức đại số hoặc lượng giác khi n $\to \infty$ (ví dụ sin(1/$n^2$) ~ 1/$n^2$).\\[4pt]
\end{congthuc}

\noindent\textbf{Chú giải \& Lưu ý kỹ thuật:}
\begin{itemize}[leftmargin=1.5em, itemsep=2pt]
  \item Chuỗi Bertrand dạng $\sum_{n=2}^\infty$ 1 / [n (ln $n)^p$] (xét bằng tiêu chuẩn tích phân với phép đổi biến u = ln x) hội tụ khi và chỉ khi p > 1.
\end{itemize}

\begin{vidu}{Xét sự hội tụ của chuỗi số bằng Tiêu chuẩn So sánh Giới hạn}
\textbf{Đề bài:} Xét sự hội tụ của chuỗi số $\sum_{n=1}^\infty \frac{2n + 1}{\sqrt{n^5 + 3n^2 + 2}}$.
\[ \sum_{n=1}^\infty a_n = \sum_{n=1}^\infty \frac{2n + 1}{\sqrt{n^5 + 3n^2 + 2}} \]
\textbf{Lời giải chi tiết:}
\begin{enumerate}[leftmargin=1.5em, itemsep=2pt]
  \item Bước 1: Khi n $\to \infty$, tử thức tương đương 2n (bậc 1), mẫu thức tương đương $\sqrt{n^5} = n^{5/2}$.
  \item Bước 2: Chọn chuỗi chuẩn so sánh $b_n$ = n / $n^{5/2}$ = 1 / $n^{3/2}$.
  \item Bước 3: Tính giới hạn lim $(a_n / b_n$) và kết luận theo p-chuỗi.
\end{enumerate}
\[ \lim_{n \to \infty} \frac{a_n}{b_n} = \lim_{n \to \infty} \frac{(2n + 1)n^{3/2}}{\sqrt{n^5 + 3n^2 + 2}} = \lim_{n \to \infty} \frac{2 + 1/n}{\sqrt{1 + 3/n^3 + 2/n^5}} = 2 \in (0, +\infty) \]
\noindent\textbf{Kết luận:} Vì chuỗi $\sum b_n = \sum$ (1/$n^{3/2}$) là p-chuỗi với p = 3/2 > 1 (hội tụ), nên chuỗi đã cho HỘI TỤ.
\end{vidu}

\section{[Mục 11.5 - 11.7] Chuỗi đan dấu (Leibniz), Hội tụ tuyệt đối, Tiêu chuẩn D'Alembert \& Cauchy}
Chuỗi đan dấu $\sum (-1)^{n-1} b_n (b_n$ > 0) được xét bằng định lý Leibniz. Với chuỗi có dấu bất kỳ, nếu chuỗi trị tuyệt đối $\sum |a_n$| hội tụ thì ta nói chuỗi HỘI TỤ TUYỆT ĐỐI (và suy ra chuỗi ban đầu cũng hội tụ). Hai công cụ mạnh nhất để xét hội tụ tuyệt đối là Tiêu chuẩn tỉ số D'Alembert và Tiêu chuẩn căn Cauchy.

\begin{congthuc}{Các công thức trọng tâm (11.5 - 11.7)}
\noindent\textbf{Tiêu chuẩn Leibniz cho Chuỗi đan dấu $\sum (-1)^{n-1} b_n$ (với $b_n$ > 0):}
\[ \begin{cases} b_{n+1} \le b_n \;(\text{dãy giảm}) \\ \lim_{n \to \infty} b_n = 0 \end{cases} \implies \sum_{n=1}^\infty (-1)^{n-1}b_n \text{ hội tụ, và sai số } |R_n| = |S - S_n| \le b_{n+1} \]
\noindent\textit{Chú giải:} Chuỗi đan dấu điều hòa $\sum (-1)^{n-1}$/n hội tụ (bằng ln 2) nhưng không hội tụ tuyệt đối (gọi là Bán hội tụ / Hội tụ có điều kiện).\\[4pt]
\noindent\textbf{Tiêu chuẩn Tỉ số D'Alembert (Ratio Test) \& Tiêu chuẩn Căn Cauchy (Root Test):}
\[ L = \lim_{n \to \infty} \left|\frac{a_{n+1}}{a_n}\right| \quad \text{hoặc} \quad L = \lim_{n \to \infty} \sqrt[n]{|a_n|} \implies \begin{cases} L < 1 : \text{Hội tụ tuyệt đối} \\ L > 1 : \text{Phân kỳ} \\ L = 1 : \text{Chưa kết luận được} \end{cases} \]
\noindent\textit{Chú giải:} Kinh nghiệm: Dùng D'Alembert khi $a_n$ chứa giai thừa (n!) hoặc lũy thừa; dùng Cauchy khi $a_n$ có dạng lũy thừa bậc n tất cả $[f(n)]^n$.\\[4pt]
\end{congthuc}

\noindent\textbf{Chú giải \& Lưu ý kỹ thuật:}
\begin{itemize}[leftmargin=1.5em, itemsep=2pt]
  \item Khi gặp bài toán chuỗi số bất kỳ, chiến lược 4 bước là: (1) Kiểm tra nhanh lim $a_n \neq$  0; (2) Nếu có n! hoặc $a^n$ dùng D'Alembert; (3) Nếu là phân thức của n dùng So sánh giới hạn với 1/$n^p$; (4) Nếu đan dấu và p $\le$  1 dùng Leibniz.
\end{itemize}

\begin{vidu}{Xét tính hội tụ tuyệt đối bằng Tiêu chuẩn D'Alembert và Cauchy}
\textbf{Đề bài:} Xét sự hội tụ của hai chuỗi số: (a) $\sum_{n=1}^\infty \frac{(-3)^n n!}{(2n)!}$ và (b) $\sum_{n=1}^\infty \left(\frac{2n + 3}{3n + 1}\right)^n$.
\[ (a)\; \sum_{n=1}^\infty \frac{(-3)^n n!}{(2n)!}, \qquad (b)\; \sum_{n=1}^\infty \left(\frac{2n + 3}{3n + 1}\right)^n \]
\textbf{Lời giải chi tiết:}
\begin{enumerate}[leftmargin=1.5em, itemsep=2pt]
  \item Câu (a): Có chứa giai thừa n! và (2n)!, áp dụng tiêu chuẩn tỉ số D'Alembert.
  \item Câu (b): Có dạng lũy thừa bậc n, áp dụng tiêu chuẩn căn Cauchy.
\end{enumerate}
\[ L_a = \lim_{n \to \infty} \left|\frac{a_{n+1}}{a_n}\right| = \lim_{n \to \infty} \frac{3(n+1)}{(2n+2)(2n+1)} = 0 < 1 \implies \text{Chuỗi (a) hội tụ tuyệt đối} \]
\[ L_b = \lim_{n \to \infty} \sqrt[n]{|a_n|} = \lim_{n \to \infty} \frac{2n + 3}{3n + 1} = \frac{2}{3} < 1 \implies \text{Chuỗi (b) hội tụ tuyệt đối} \]
\noindent\textbf{Kết luận:} Cả hai chuỗi đều hội tụ tuyệt đối vì giới hạn tiêu chuẩn L < 1.
\end{vidu}

\section{[Mục 11.8 - 11.11] Chuỗi lũy thừa, Bán kính hội tụ \& Khai triển chuỗi Taylor - Maclaurin}
Chuỗi lũy thừa có dạng $\sum_{n=0}^\infty c_n$ (x - $a)^n$. Luôn tồn tại Bán kính hội tụ R $\ge$  0 sao cho chuỗi hội tụ tuyệt đối khi |x - a| < R và phân kỳ khi |x - a| > R. Nếu hàm f(x) khả vi vô hạn lần tại a, chuỗi Taylor của nó biểu diễn chính xác hàm số thông qua các đạo hàm cấp cao tại tâm a.

\begin{congthuc}{Các công thức trọng tâm (11.8 - 11.11)}
\noindent\textbf{Công thức tìm Bán kính hội tụ R của chuỗi lũy thừa $\sum c_n$ (x - $a)^n$:}
\[ R = \lim_{n \to \infty} \left|\frac{c_n}{c_{n+1}}\right| \quad \text{hoặc} \quad R = \frac{1}{\lim_{n \to \infty} \sqrt[n]{|c_n|}} \]
\noindent\textit{Chú giải:} Để tìm Tập (Miền) hội tụ hoàn chỉnh, sau khi tìm khoảng (a - R, a + R), bắt buộc phải thay trực tiếp 2 đầu mút x = a - R và x = a + R vào chuỗi để xét riêng!\\[4pt]
\noindent\textbf{Công thức Chuỗi Taylor tại x = a và Chuỗi Maclaurin (tại a = 0):}
\[ f(x) = \sum_{n=0}^\infty \frac{f^{(n)}(a)}{n!}(x - a)^n = f(a) + f'(a)(x - a) + \frac{f''(a)}{2!}(x - a)^2 + \cdots \]
\noindent\textit{Chú giải:} Chuỗi lũy thừa có thể lấy đạo hàm và lấy tích phân từng số hạng bên trong khoảng hội tụ (-R, R).\\[4pt]
\end{congthuc}

\noindent\textbf{Chú giải \& Lưu ý kỹ thuật:}
\begin{itemize}[leftmargin=1.5em, itemsep=2pt]
  \item Để tìm khai triển Maclaurin của các hàm phức tạp, KHÔNG nên tính đạo hàm cấp cao trực tiếp mà hãy biến đổi từ 5 khai triển Maclaurin chuẩn $(e^x$, sin x, cos x, 1/(1-x), ln(1+x)) kết hợp phép thế biến, đạo hàm hoặc tích phân từng số hạng.
\end{itemize}

\begin{vidu}{Tìm miền hội tụ của chuỗi lũy thừa và khai triển Maclaurin}
\textbf{Đề bài:} Tìm bán kính hội tụ và miền hội tụ của chuỗi lũy thừa $\sum_{n=1}^\infty \frac{(x - 2)^n}{n \cdot 3^n}$.
\[ \sum_{n=1}^\infty \frac{(x - 2)^n}{n \cdot 3^n} \]
\textbf{Lời giải chi tiết:}
\begin{enumerate}[leftmargin=1.5em, itemsep=2pt]
  \item Bước 1: Xác định tâm chuỗi a = 2 và hệ số $c_n$ = 1 / (n $\cdot 3^n$). Tính bán kính hội tụ R = lim $|c_n / c_{n+1}$|.
  \item Bước 2: Khoảng hội tụ là (2 - 3, 2 + 3) = (-1, 5).
  \item Bước 3: Xét riêng tại x = -1 (được chuỗi đan dấu điều hòa $\to$  hội tụ) và tại x = 5 (được chuỗi điều hòa $\to$  phân kỳ).
\end{enumerate}
\[ R = \lim_{n \to \infty} \frac{(n + 1)3^{n+1}}{n \cdot 3^n} = 3 \implies \text{Khoảng hội tụ: } |x - 2| < 3 \iff -1 < x < 5 \]
\[ \text{Tại } x = -1: \sum_{n=1}^\infty \frac{(-1)^n}{n} \;(\text{Hội tụ theo Leibniz}); \quad \text{Tại } x = 5: \sum_{n=1}^\infty \frac{1}{n} \;(\text{Phân kỳ}) \]
\noindent\textbf{Kết luận:} Bán kính hội tụ R = 3 và miền hội tụ của chuỗi là nửa khoảng [-1, 5).
\end{vidu}

\section{Minh họa Hình học bằng TikZ}
\begin{figure}[htbp]
\centering
\begin{tikzpicture}[scale=0.95, >=stealth]
  % Trục tọa độ
  \draw[->, thick] (-3.8,0) -- (3.8,0) node[right] {$x$};
  \draw[->, thick] (0,-2.2) -- (0,2.2) node[above] {$y$};
  
  % Đồ thị chuẩn y = sin(x)
  \draw[domain=-3.5:3.5, samples=100, smooth, variable=\x, blue!85!black, very thick] 
    plot ({\x}, {sin(\x r)}) node[right] {$\sin x$};
    
  % Đa thức Taylor bậc 1: T_1(x) = x
  \draw[domain=-1.9:1.9, smooth, variable=\x, gray!80, thick, dotted] 
    plot ({\x}, {\x}) node[above right] {$T_1(x) = x$};
    
  % Đa thức Taylor bậc 3: T_3(x) = x - x^3/6
  \draw[domain=-2.7:2.7, smooth, variable=\x, orange!90!black, thick, dashed] 
    plot ({\x}, {\x - (\x)^3/6}) node[below right] {$T_3(x)$};
    
  % Đa thức Taylor bậc 5: T_5(x) = x - x^3/6 + x^5/120
  \draw[domain=-3.3:3.3, smooth, variable=\x, red!80!black, thick] 
    plot ({\x}, {\x - (\x)^3/6 + (\x)^5/120}) node[above left] {$T_5(x)$};
\end{tikzpicture}
\caption{Hình 11.1: Sự hội tụ của các đa thức Taylor $T_1$(x), $T_3$(x), $T_5$(x) ôm sát đồ thị hàm f(x) = sin(x) khi tăng bậc khai triển.}
\label{fig:tikz-ch11-taylor}
\end{figure}

\section{Ứng dụng Thực tế \& Tối ưu hóa (Vi xử lý Máy tính (FPU) \& Giải tích Số)}
\begin{ungdung}{Cách CPU/Máy tính Cầm tay Tính hàm Siêu việt \& Tích phân Không có Nguyên hàm Sơ cấp}
Bộ xử lý toán học trong máy tính chỉ thực hiện trực tiếp các phép cộng, trừ, nhân, chia nên không thể tính trực tiếp sin(x), $e^x$ hay ln(x). Chúng sử dụng đa thức xấp xỉ Taylor/Maclaurin (hoặc đa thức Chebyshev rút gọn từ chuỗi Taylor). Ngoài ra, chuỗi Maclaurin giúp tính dễ dàng các tích phân không có nguyên hàm sơ cấp như tích phân xác suất $\int e^{-x^2}$ dx hay tích phân quang học Fresnel $\int sin(x^2$) dx.
\[ \int_0^x e^{-t^2}\,dt = \int_0^x \sum_{n=0}^\infty (-1)^n \frac{t^{2n}}{n!}\,dt = \sum_{n=0}^\infty (-1)^n \frac{x^{2n+1}}{n!(2n+1)} = x - \frac{x^3}{3} + \frac{x^5}{10} - \frac{x^7}{42} + \cdots \]
\noindent\textbf{Ví dụ tính toán thực tế:} Để tính $\int_0^{0.5} e^{-t^2}$ dt: Chỉ cần lấy 3 số hạng đầu của chuỗi Maclaurin ta có 0.5 - (0.$5)^3$/3 + (0.$5)^5$/10 = 0.5 - 0.04167 + 0.003125 = 0.46146 (chính xác tới 5 chữ số thập phân!).
\end{ungdung}

\section{Bài tập Trắc nghiệm Củng cố Chương 11}
\begin{enumerate}[label=\textbf{Câu \arabic*:}, leftmargin=2.2em, itemsep=8pt]
  \item Tổng của chuỗi cấp số nhân hội tụ $\sum_{n=1}^\infty \frac{3}{4^n}$ bằng:
  \begin{itemize}[label={}, leftmargin=1em, itemsep=2pt]
    \item \textbf{A.} 1 \quad ($1$)
    \item \textbf{B.} 4 \quad ($4$)
    \item \textbf{C.} 3/4 \quad ($\frac{3}{4}$)
    \item \textbf{D.} 4/3 \quad ($\frac{4}{3}$)
  \end{itemize}
  \vspace{2pt}
  \noindent{\color{teal!85!black}\textbf{Đáp án đúng: A.}} \textit{Giải thích:} Đây là chuỗi cấp số nhân bắt đầu từ n = 1 có số hạng đầu a = 3/4 và công bội r = 1/4 (< 1). Tổng S = a / (1 - r) = (3/4) / (1 - 1/4) = 1.
  \[ S = \frac{a_1}{1 - r} = \frac{3/4}{1 - 1/4} = 1 \]
  \item Khai triển Maclaurin của hàm số f(x) = $e^{-x^2}$ là:
  \begin{itemize}[label={}, leftmargin=1em, itemsep=2pt]
    \item \textbf{A.} $\sum_{n=0}^\infty (-1)^n x^{2n}$ / n! \quad ($\sum_{n=0}^\infty (-1)^n \frac{x^{2n}}{n!}$)
    \item \textbf{B.} $\sum_{n=0}^\infty (-1)^n x^{2n}$ / (2n)! \quad ($\sum_{n=0}^\infty (-1)^n \frac{x^{2n}}{(2n)!}$)
    \item \textbf{C.} $\sum_{n=0}^\infty x^{2n}$ / n! \quad ($\sum_{n=0}^\infty \frac{x^{2n}}{n!}$)
    \item \textbf{D.} $\sum_{n=0}^\infty (-1)^n x^n$ / n! \quad ($\sum_{n=0}^\infty (-1)^n \frac{x^n}{n!}$)
  \end{itemize}
  \vspace{2pt}
  \noindent{\color{teal!85!black}\textbf{Đáp án đúng: A.}} \textit{Giải thích:} Từ khai triển chuẩn $e^u = \sum_{n=0}^\infty u^n$ / n!, thay u = $-x^2$ vào ta thu được $e^{-x^2} = \sum_{n=0}^\infty (-x^2 )^n$ / n! = $\sum_{n=0}^\infty (-1)^n x^{2n}$ / n!.
  \[ e^u = \sum_{n=0}^\infty \frac{u^n}{n!} \xrightarrow{u = -x^2} e^{-x^2} = \sum_{n=0}^\infty (-1)^n \frac{x^{2n}}{n!} \]
  \item Sử dụng khai triển Maclaurin, giới hạn L = \$\lim_{x \to 0} \frac{\sin x - x + x^3/6}{x^5}$ có giá trị bằng:
  \begin{itemize}[label={}, leftmargin=1em, itemsep=2pt]
    \item \textbf{A.} 1/120 \quad ($\frac{1}{120}$)
    \item \textbf{B.} -1/120 \quad ($-\frac{1}{120}$)
    \item \textbf{C.} 1/24 \quad ($\frac{1}{24}$)
    \item \textbf{D.} 0 \quad ($0$)
  \end{itemize}
  \vspace{2pt}
  \noindent{\color{teal!85!black}\textbf{Đáp án đúng: A.}} \textit{Giải thích:} Khai triển Maclaurin của sin x đến bậc 5 là: sin x = x - $x^3$/3! + $x^5$/5! - $O(x^7$) = x - $x^3$/6 + $x^5$/120 - ... Thay vào tử thức rồi chia cho $x^5$ và cho x $\to$  0 ta được ngay 1/120 (nhanh hơn rất nhiều so với 5 lần L'Hospital).
  \[ \sin x = x - \frac{x^3}{6} + \frac{x^5}{120} - \cdots \implies \lim_{x \to 0} \frac{\frac{x^5}{120} + o(x^5)}{x^5} = \frac{1}{120} \]
\end{enumerate}

\end{document}
